% concatenated from main.tex
\documentclass[10pt,reqno]{amsart}
\usepackage{setspace}
\onehalfspacing
\usepackage[utf8]{inputenc}
\usepackage{graphicx}
\graphicspath{ {./images/} }
\setcounter{tocdepth}{2}
\usepackage[pagewise]{lineno}
%\linenumbers
\setcounter{tocdepth}{3}
\usepackage{amssymb}
\usepackage[foot]{amsaddr}
\usepackage[all,arc]{xy}
\usepackage{enumerate}
\usepackage{mathrsfs}
%\usepackage{pst-node}
\usepackage[margin=1in]{geometry} 
\usepackage{amsmath,amsthm}
\usepackage{mathtools}
\usepackage[T1]{fontenc}
\usepackage{fourier}
%\usepackage{baskervald}
%\usepackage{showkeys}
\usepackage{stmaryrd}
\usepackage{amsfonts}
%\usepackage[english,russian]{babel}
\usepackage{cite}


\usepackage[backref=page]{hyperref}
\usepackage[capitalise,nameinlink]{cleveref}
\newcommand*\bb{\mathbb}
\newcommand*\wh{\widehat}
\newcommand *\w{^\wedge}
\newcommand*\g{\mathfrak{g}}
\newcommand*\diff{\delta_{x,y}}
\DeclareMathOperator*{\esup}{ess\,sup}
\DeclareMathOperator*{\einf}{ess\,inf}
\DeclareMathOperator*{\eosc}{ess\,osc}
\newcommand{\cprime}{\('\)}

\newcommand*\Om{\Omega}
\newcommand*\de{\partial}
\newcommand*\test{$C_c^{\infty}$}

\newcommand*\tl{\mathrm{Tail}}

\newtheorem{thm}{Theorem}[section]
\theoremstyle{definition}
\newtheorem{defn}[thm]{Definition}
\newtheorem{example}[thm]{Example}
\newtheorem{cor}[thm]{Corollary}
\newtheorem{prop}[thm]{Proposition}
\newtheorem{lem}[thm]{Lemma}
\newtheorem{rem}[thm]{Remark}

\newtheorem{theorem}{Theorem}[section]
\newtheorem{lemma}[theorem]{Lemma}
\newtheorem{proposition}[theorem]{Proposition}
\newtheorem{definition}[theorem]{Definition}
\newtheorem{remark}[theorem]{Remark}

%%%%%%%%%%%%%%%%%%%%%%%%%%%%%%%CREF%%%%%%%%%%%%%%%%%%%%%
\usepackage{xcolor}
\hypersetup{
	colorlinks,
	linkcolor={blue!100!black},
	citecolor={red!100!black},
	urlcolor={green!100!black}
}


%%%%%%%%%%%%%%%%%%%%%%%%%%%%%%%%%%%%%%%%%%%%%%
%%%%%    PRASASD 19 AUG 2024 % %%%% %%%%%%% %
%%%%%%%%%%%%%%%%%%%%%%%%%%%%%%%%%%%%%%%%%%%%%%


%Inner Product
\DeclarePairedDelimiterX{\inp}[2]{\langle}{\rangle}{#1, #2}
\newcommand{\norm}[1]{\left\lVert#1\right\rVert}


%\makeatletter
%\let\c@equation\c@thm
%\makeatother

\makeatletter
\def\@tocline#1#2#3#4#5#6#7{\relax
	\ifnum #1>\c@tocdepth % then omit
	\else
	\par \addpenalty\@secpenalty\addvspace{#2}%
	\begingroup \hyphenpenalty\@M
	\@ifempty{#4}{%
		\@tempdima\csname r@tocindent\number#1\endcsname\relax
	}{%
		\@tempdima#4\relax
	}%
	\parindent\z@ \leftskip#3\relax \advance\leftskip\@tempdima\relax
	\rightskip\@pnumwidth plus4em \parfillskip-\@pnumwidth
	#5\leavevmode\hskip-\@tempdima
	\ifcase #1
	\or\or \hskip 1em \or \hskip 2em \else \hskip 3em \fi%
	#6\nobreak\relax
	\dotfill\hbox to\@pnumwidth{\@tocpagenum{#7}}\par
	\nobreak
	\endgroup
	\fi}
\makeatother

\usepackage{hyperref}

\newcommand{\overbar}[1]{\mkern 1.5mu\overline{\mkern-1.5mu#1\mkern-1.5mu}\mkern 1.5mu}

\numberwithin{equation}{section}


\usepackage{epigraph}

\title{A Class of De Giorgi Type and H\"older Continuity for Some Problems in Musielak-Orlicz-Sobolev Spaces}

\author{Hlel Missaoui$^{1}$}
\author{Anouar Bahrouni$^{2}$}
\author{Hichem Ounaies$^{3}$}
\address{$^{1}$ Mathematics Department,  Higher Institute of Computer Science of Mahdia, University of Monastir, 5111 Mahdia, Tunisia.}
\address{$^{2}$ $^{3}$ Mathematics Department, Faculty of Sciences, University of Monastir, 5019 Monastir, Tunisia.}
\email{$^{1}$ \tt Hlel.Missaoui@fsm.rnu.tn}
\email{$^{2}$ \tt bahrounianouar@yahoo.fr}
\email{$^{3}$ \tt Hichem.Ounaies@fsm.rnu.tn}
%\thanks{$^{\ast}$ Support from the National Research Foundation of Korea (NRF) through IRTG 2235/NRF-2016K2A9A2A13003815 at Seoul National University.}
%\thanks{$^{\dagger}$ Support from funding of the Deutsche Forschungsgemeinschaft through GRK 2235/2 2021 - 282638148.}
%\thanks{$^{\S}$ Support from Deutsche Forschungsgemeinschaft (DFG, German Research Foundation) – Project-ID 317210226 – SFB 1283 is gratefully acknowledged}
\date{}
\begin{document}
\begin{abstract}
In this paper, we introduce a new class of De Giorgi type functions, denoted by \(\mathcal{B}_{G(x,t)}\), and establish the H\"older continuity of its elements under suitable additional assumptions on the generalized \textnormal{N}-function \(G(x,t)\). As an application, we prove the H\"older continuity of solutions to quasilinear equations whose principal part is in divergence form with \(G(x,t)\)-growth conditions, including both critical and standard growth cases. The novelty of our work lies in the generalization of the H\"older continuity results previously known for variable exponent \cite[X, Fan and D. Zhao]{Fan1999} and Orlicz \cite[G. M. Lieberman]{Li1991} problems. Moreover, our results encompass a wide variety of quasilinear equations.

\end{abstract}
\maketitle
\noindent \textbf{Keywor{\rm d}s.} Class of De Giorgi Type; H\"older Continuity; Musielak-Orlicz-Sobolev Spaces; Weak solutions.

\noindent\textbf{\textup{2020} Mathematics Subject Classification.} {Primary: 
35R09, %Integro-partial differential equations
35B65; %Smoothness and regularity of solutions to PDEs
Secondary: 
47G20, %Integro-differential operators
35K55%Nonlinear parabolic equations
}

\smallskip
	
\tableofcontents


\section{Introduction}
\label{sec:intro}
\subsection{Overview}
The qualitative theory of partial differential equations studies the behavior of weak solutions
$u \in W^{1,1}_{\mathrm{loc}}(\Omega)$ of problems of the form
\begin{equation}\label{main_eq}
-\operatorname{div} A(x,u,D u) = B(x,u,D u) \quad \text{in } \Omega.
\end{equation}
Existence results, based on the calculus of variations and topological methods, usually give solutions only in Sobolev spaces. 
To use these solutions in applications, it is important to know if they are more regular. 
In particular, one wants Hölder continuity $u \in C^{0,\alpha}_{\mathrm{loc}}(\Omega)$. 
This property lies between boundedness ($L^\infty$) and differentiability ($C^{1,\alpha}$), and is crucial for uniqueness, stability, and free boundary problems.\\

The study of regularity began with the famous works of De Giorgi \cite{deGiorgi1957} and Nash \cite{Nash1958}, later completed by Moser \cite{Moser1960,Moser1961}. 
They proved that weak solutions of linear elliptic equations in divergence form,
\[
-\operatorname{div}(A(x)D u)=0,
\]
with bounded measurable coefficients $A(x)$ satisfying uniform ellipticity
$\lambda |\xi|^2 \le \langle A(x)\xi,\xi\rangle \le \Lambda |\xi|^2$,
are H\"older continuous. 
This result was very surprising, because it did not require smoothness of $A(x)$. 
It introduced two powerful techniques, the Moser iteration and the De Giorgi measure argument.
Later, attention moved to quasilinear equations, such as the $p$-Laplacian
\[
\Delta_p u := \operatorname{div}(|D u|^{p-2}D u).
\]
Uraltseva \cite{Uraltseva1968}, Uhlenbeck \cite{Uhlenbeck1977}, and Lieberman \cite{Li1988} proved that solutions of $-\Delta_p u=0$ are $C^{1,\alpha}$. 
But this regularity is delicate for $-\Delta_p u=f$, if $f \in L^q$ with $q<n/p$, the solutions may not even be continuous \cite{MG2,marc1,Mar1}. 
This shows that the regularity program must proceed in two steps: first, prove $u\in L^\infty$ (by De Giorgi’s lemma or Moser iteration), then prove $u\in C^{0,\alpha}$.\\

These ideas and techniques were systematically extended to tackle the more general non-polynomial growth of operators in the Orlicz setting. The foundational work by Donaldson and Trudinger \cite{DonaldsonTrudinger1971} first established the boundedness ($L^\infty$) of solutions.  Lieberman \cite{Li1991}, who proved that weak solutions to quasilinear equations of the form $-\operatorname{div}(a(|D u|)D u) = 0$ are locally H\"older continuous ($C^{0,\alpha}$). This required the associated \textnormal{N}-function $G(t)$ to satisfy the condition $\Delta_2$, which is essential for the Sobolev-Poincar\'e inequalities and the scaling arguments of the De Giorgi-Nash-Moser theory. These contributions laid the foundation of the regularity theory in the Orlicz framework.\\

Although Orlicz spaces encompass many classical examples, there remain important cases that they do not cover. For this reason, considerable research has focused on the variable exponent framework, where the growth is governed by a function $p(x)$. The analysis of problems involving terms of the type $t^{p(x)}$ was initiated by Zhikov; see \cite{Zhikov1987,Zhikov1995}. Subsequently, Fan and Zhao \cite{Fan1999} uncovered a new regularity phenomenon by showing that weak solutions of equations of the form
\[
-\operatorname{div}\big(|Du|^{p(x)-2}Du\big)=0
\]
are locally H\"older continuous, provided that the exponent function $p(x)$ is log-H\"older continuous. This condition, namely
\[
|p(x)-p(y)| \le \frac{C}{-\ln |x-y|} \quad \text{for } |x-y|<1,
\]
was shown to be not only sufficient but essentially necessary for regularity, a fact later clarified and reinforced by the work of Acerbi and Mingione \cite{Acerbi2007}.\\

A further major advancement in the theory of non-standard growth problems came with the analysis of double-phase functionals, which introduce a new type of heterogeneity: the growth exponent remains fixed, but the structure of the operator changes depending on the location in the domain. For functionals of the form
\[
\operatorname{div}\!\left(|D u|^{p-2}D u + a(x)|D u|^{q-2}D u\right),
\]
with $1 < p \leq q$, the regularity theory reveals a novel feature: the regularity of solutions depends directly on the regularity of the coefficient $a(x)$. Colombo and Mingione~\cite{Colombo2015, Colombo2015b} proved that in the natural growth case, boundedness and H\"older continuity of minimizers and solutions hold provided $a(x)$ is H\"older continuous, i.e., $a(x) \in C^{0,\alpha}$. This condition is sharp: if $a(x)$ is merely continuous, minimizers may fail to be locally bounded~\cite{Colombo2015}. This phenomenon highlights a key principle: in problems with non-uniformly elliptic structures, the regularity of the coefficients controlling the growth is often necessary for the regularity of solutions. The boundedness result in this setting was also independently obtained by Baroni, Colombo, and Mingione~\cite{Baroni2016}.\\

Due to the vast scope and rapid development of regularity theory, we cannot cite all contributions individually. We acknowledge with gratitude the many valuable works not mentioned here and instead refer the reader to comprehensive surveys and bibliographies, such as~\cite{mr3, Mingione2006}. For detailed discussions of specific regularity results, we direct the reader to the following foundational works and the extensive literature contained therein: for the classical Laplacian and linear theory, see~\cite{deGiorgi1957, Nash1958, Moser1960, Moser1961}; for the $p$-Laplacian and $(p,q)$-growth, see~\cite{ GmeinederKristensen2024, Uraltseva1968,  KristensenTaheri2003, Li1988, Uhlenbeck1977,  KristensenMingione2008, DeFilippisKochKristensen2024,  KristensenMingione2010, Manfredi1986, Manfredi1988, marc1, marc2, Ac1, Ac2, mg3, MG3, NF, LE, MH, Pucci}; for Orlicz and other non-standard growth conditions, see~\cite{Li1991, Simon1,HsK2,HsK1,HHL21, Harjulehto2017, Cupini2017, Eleuteri2016, Chlebicka2021, Chlebicka2022, mgg}; for double-phase problems, see~\cite{Colombo2015, Colombo2015b, Baroni2016,defi,Simon1}; and for variable exponent $p(x)$-Laplacian problems, see~\cite{Fan1999, FanZhao1999, Acerbi2007, Coscia1999, KKPZ}.\\

 The Musielak--Orlicz growth framework represents the most general setting, as it
encompasses all previously studied models, including the standard $p$-growth,
Orlicz growth, variable exponent growth, and double-phase (or $(p,q)$) growth. A major advance in the regularity theory for generalized Orlicz growth was achieved
by Harjulehto, Hästö, and Lee \cite{HHL21}, who established Hölder continuity of
local minimizers of integral functionals with Musielak--Orlicz growth under very
general structural assumptions, namely (A0), (A1), and (A2) (see
Subsection~\ref{sec2}). In particular, their results do not require any logarithmic
continuity condition on the growth function, such as \eqref{GG3} (see
Subsection~\ref{sec1.2}), and they cover important models including the double-phase
case. This highlights the strength and generality of the variational approach
developed in \cite{HHL21} for the regularity of minimizers, which relies on
comparison estimates and harmonic approximation techniques.

The present paper is concerned with a different class of problems. We study weak
solutions of quasilinear elliptic equations and systems in divergence form \eqref{main_eq},
possibly including lower-order terms, for which a variational structure is not
available in general. Consequently, the methods developed in \cite{HHL21} cannot
be applied directly in our setting. Instead, we adopt a De Giorgi--type iteration
adapted to Musielak--Orlicz growth conditions. Within this non-variational
framework, an additional control on the spatial oscillation of the effective
growth, formulated in assumption \eqref{GG3}, is required in order to derive the
key structural estimates and to obtain Hölder regularity of solutions. We also
emphasize that our analysis covers Neumann boundary conditions, which are not
addressed in the variational framework of \cite{HHL21}.

To the best of our knowledge, the literature contains only one other work
specifically devoted to the Hölder continuity of weak solutions to divergence-form
elliptic problems under Musielak--Orlicz growth conditions, namely the paper by
Wang, Liu, and Zhao \cite{WaLiuZhao}. In their analysis, the authors impose a
restrictive structural assumption on the growth function $G(x,t)$. More
precisely, they assume the existence of a strictly increasing differentiable
function $\xi : [0,+\infty) \to [0,+\infty)$ satisfying
\[
n\,\xi(s) > s\,\xi'(s) > \xi(s) \quad \text{for all } s \ge 0,
\]
such that
\begin{itemize}
    \item[(A11)] $G(x, s t) \ge \xi(s)\, G(x,t)$ for all $s \ge 0$, $t \in \mathbb{R}$,
    and $x \in \Omega$;
    \item[(A12)] $\xi, \xi^{-1}, \xi_{\ast}, \widehat{\widetilde{\xi}} \in
    \Delta_{\mathbb{R}^{+}}$.
\end{itemize}
A significant limitation of this framework is that assumption (A11) excludes
several fundamental models in the theory of non-standard growth. In particular, it
is not satisfied by variable exponent growth, $(p,q)$-growth equations, or
double-phase functionals, which are among the main motivating examples of the
Musielak--Orlicz setting. As a consequence, while the results of \cite{WaLiuZhao}
are obtained under a specific and interesting set of hypotheses, their approach
does not provide a unified treatment of these central growth models.
\\

Motivated by the preceding discussion, in this paper we establish the local H\"older continuity of weak solutions to~\eqref{main_eq} under general Musielak--Orlicz growth conditions on the operators $A$ and $B$.
The principal novelty of our work lies in providing a unified framework that simultaneously covers the classical cases of the $p$-Laplacian, variable exponent, and Orlicz-type equations, while also treating both subcritical and natural growth (critical) cases with either Dirichlet or Neumann boundary conditions. Note that equation \eqref{main_eq} involves a critical term, which makes the study of regularity more delicate and necessitates a careful analysis of the associated Sobolev conjugate for generalized Young functions.
 The study of such conjugates was initiated by Fan~\cite{Fan2012a}, who introduced the concept under the assumption that $A$ satisfies a Lipschitz condition in the spatial variable. This theory was significantly advanced by Cianchi and Diening~\cite{Cianchi2024}, who established a sharper formulation under weaker regularity assumptions. However, despite these advances, a complete theory of Sobolev conjugates for generalized Young functions remains incomplete, lacking crucial structural properties such as the $\Delta_2$-condition, which is notoriously difficult to verify.\\
 
 As explained above, our work has two main contributions. First, we establish the local Hölder continuity of weak solutions to \eqref{main_eq} under general Musielak–Orlicz growth conditions on the operators $A$, which in particular generalizes the variable-exponent type operators treated in \cite{Fan1999}. Moreover, our work closes the gap left in \cite{Fan1999}, where the critical cases were excluded. At the same time, it extends another class of equations in the framework of $(p,g)$-growth with critical term; see Subsection~1.3.\\
 %This cann`t works in unbounded domain due the behaviour the Sobolev conjugate in neighboord of $0$. To pass this problem we suppose an additional  assumption concerning the Sobolev  conjugate $\Phi_d$, see \eqref{Hast}.\\

Four main difficulties arise from the Musielak--Orlicz structure when establishing regularity.  
First, the energy is nonhomogeneous and exhibits explicit $x$-dependence alongside the natural $t$-dependence, which prevents the direct application of classical scaling and iteration techniques such as the De Giorgi method.  
Second, the regularity proofs are highly sensitive to the precise parameters controlling the growth and structure of the function $G(x,t)$, requiring careful balancing of assumptions to ensure the validity of key inequalities.  
Third, even when adapted versions of De Giorgi-type methods are applicable, they may not guarantee regularity, as demonstrated by the counterexamples of Giaquinta~\cite{MG2} and Marcellini~\cite{marc2, marc1}, which show that weak solutions may fail to be regular even under standard structural assumptions.  
Finally, the critical case introduces additional challenges compared to the subcritical case, due to the lack of uniform integrability and the need for more refined estimates.  
These aspects collectively indicate that establishing regularity under non-standard growth conditions requires new ideas and a delicate analysis beyond the classical framework.\\

Our approach is motivated by the works of Lieberman~\cite{Li1991} and Fan--Zhao~\cite{FanZhao1999}, and in particular by the classical techniques of De~Giorgi~\cite{deGiorgi1957} and Ladyzhenskaya--Ural'tseva~\cite{OL}. 
These foundational contributions introduced powerful tools for proving Hölder continuity of solutions. 
In particular, they defined the class $\mathcal{B}(\Omega, M, \gamma, \gamma_1, \delta, 1/q)$ and showed that all functions in this class are Hölder-continuous. 
This framework is especially effective under standard $m$-growth conditions (see, e.g.~\cite{MG1, MG3, OL, CM}). 
For the treatment of critical cases, we further rely on the method developed by Ho, Kim, Winkert, and Zhang~\cite{KKPZ}, and we also make use of some fundamental properties of the Sobolev conjugate established in \cite{Ala3}.
%----------------
\subsection{Assumptions and Main Results}\label{sec1.2}
For the sake of clarity, in this subsection, we present only the main assumptions, key results, and core definitions. Additional definitions and technical details will be introduced in the subsequent sections.

Let $\Omega$ be a bounded domain in $\mathbb{R}^n$ with \( n \geq 2 \), and let 
\[G(x,t) = \int_0^{\vert t\vert} g(x,s)\, \mathrm{d}s,\]
be a generalized \textnormal{N}-function (see Section~\ref{sec2} for the precise definitions) satisfying the following structural conditions:\\
There exist constants $1 < g^- \leq g^+ < n$ such that for all $x \in \mathbb{R}^n$, $t > 0$:
\begin{equation}\label{D22}
1 < g^- \leq \frac{g(x,t)t}{G(x,t)} \leq g^+ .
\tag{G0}
\end{equation}
There exists $F \geq 1$ such that for almost every $x \in \mathbb{R}^n$:
\begin{equation}\label{GG1}
F^{-1} \leq G(x,1) \leq F.
\tag{G1}
\end{equation}
There exists $\mu \in (0,1]$ such that for every ball $B \subset \mathbb{R}^n$ with $|B| \leq 1$, for every $t \in [1, 1/|B|]$, and for almost all $x,y \in B$:
\begin{equation}\label{GG2}
\mu G^{-1}(x,t) \leq G^{-1}(y,t).
\tag{G2}
\end{equation}
There exists $L_0 > 1$ such that, for every $x_0 \in \Omega$, for all sufficiently small $R > 0$ and for all $t>1$, we heve
\begin{equation}\label{GG3} 
\eta(R):= \sup_{x,y \in B_R(x_0)} \left| r\left(x,t\right) - r\left(y,t\right) \right| 
\leq \frac{L_0}{|\ln (2R)|},
\tag{G3}
\end{equation}
where $r(x,t) := \dfrac{t g(x,t)}{G(x,t)}$. For the definition of $G^{-1}$ see Definition \ref{definv}.\\

    
We now introduce a De Giorgi-type class of functions tailored to our framework:
\begin{definition}\label{Definition 2.1}
Let \( M, \gamma, \gamma_1 \), and \( \delta \) be positive constants with \( \delta \leq 2 \). A function \( u \) in the Musielak-Orlicz-Sobolev space \(W^{1,G(x,t)}(\Omega) \) (The definition of Musielak–Sobolev spaces is provided in Subsection 2.2), satisfying \( \max_{\Omega} |u(x)| \leq M \), is said to belong to the class \( \mathcal{B}_{G(x,t)}(\Omega, M, \gamma, \gamma_1, \delta) \) if the functions \( w(x) := \pm u(x) \) satisfy the inequality
\begin{equation}\label{2.4}
\int_{A_{k,\rho - \sigma \rho}} G(x,|Dw|) \, \mathrm{d}x \leq \gamma \int_{A_{k,\rho}} G\left(x,\left| \frac{w(x) - k}{\sigma \rho} \right|\right) \, \mathrm{d}x + \gamma_1 |A_{k,\rho}|,
\end{equation}
for every ball \( B_\rho \subset \Omega \), every \( \sigma \in (0,1) \), and every threshold \( k \) satisfying
\begin{equation}\label{2.5}
k \geq \max_{B_\rho} w(x) - \delta M,
\end{equation}
where \( A_{k,\rho} := \{x \in B_\rho : w(x) > k\} \).
\end{definition}
One of the central results of this work is the following theorem:
\begin{theorem}\label{Theorem 2.1}
Suppose that the function \( G(x,t) \) satisfies conditions \eqref{D22}--\eqref{GG1} and \eqref{GG3}. Then the De Giorgi-type class belongs to a Hölder space:
\[
\mathcal{B}_{G(x,t)}(\Omega, M, \gamma, \gamma_1, \delta) \subset C^{0,\alpha}(\Omega),
\]
where the Hölder exponent \( \alpha \in (0,1] \) depends only on the parameters \( n, g^-, g^+, F, L:=\exp{L_0}, \gamma \), and \( \delta \), and is independent of $M$ and \( \gamma_1 \).
\end{theorem}
\begin{definition}[see Sect. 1, Chapter 1 \cite{OL}]\label{Definition 2.3}
We will say that \(u \in \mathcal{B}_{G(x,t)}(\overline{\Omega}, M, \gamma, \gamma_1, \delta)\) if \(u \in \mathcal{B}_{G(x,t)}(\Omega, M, \gamma, \gamma_1, \delta)\) and the inequalities \eqref{2.4} hold also for arbitrary ball \(B_\rho\) with \(B_\rho \cap \partial \Omega \neq \phi\), \(\sigma \in (0, 1)\) and
\[
k \geq \max \left\{ \max_{\Omega_\rho} w(x) - \delta M, \max_{\partial \Omega_\rho} w(x) \right\},
\]
where \(\Omega_\rho = B_\rho \cap \Omega\), \(\partial \Omega_\rho = \partial \Omega \cap B_\rho\), \(A_{k, \rho} = \{x \in \Omega_\rho: w(x) > k\}\).
\end{definition}
\begin{definition}[see Sect. 1, of Chap. I of \cite{OL}]\label{Definition 2.4}
We will say that the boundary \(\partial \Omega\) of \(\Omega\) satisfies the condition \eqref{AO}, if there exist two positive constants \(m_0\) and \(\theta_0\) such that for any ball \(B_\rho\) with center on \(\partial \Omega\) and radius \(\rho \leq m_0\) and for any connected branch \(\Omega'_\rho\) of \(B_\rho \cap \Omega\), the following inequality hol{\rm d}s:
\begin{equation}\label{AO}
|\Omega'_\rho| \leq (1 - \theta_0)|B_\rho|.\tag{A}
\end{equation}
\end{definition}
Our second main result is the following theorem.
\begin{theorem}\label{Theorem 2.2}
Let \(u \in \mathcal{B}_{G(x,t)}(\overline{\Omega}, \gamma, \gamma_1, \delta)\). If \(G(x,t)\) satisfies conditions \eqref{D22}--\eqref{GG1} and \eqref{GG3}, \(\partial \Omega\) satisfies the condition \eqref{AO} and \(u|_{\partial \Omega} \in C^{0,\alpha_1}(\partial \Omega)\), then \(u \in C^{0,\alpha}(\overline{\Omega})\) where \(\alpha = \alpha(n, g^-, g^+, F, L, \delta, \alpha_1)\).
\end{theorem}

As an application of Theorems~\ref{Theorem 2.1} and \ref{Theorem 2.2}, we investigate the Hölder continuity of weak solutions to the elliptic equation in divergence form:
\begin{equation}\label{P}
\operatorname{div} A(x, u, Du) + B(x, u, Du) = 0, \quad x \in \Omega, \tag{$\mathcal{P}$}
\end{equation}
with Dirichlet boundary condition:
\begin{equation}\label{PD}
\left\lbrace\begin{array}{rcl}
\operatorname{div} A(x, u, Du) + B(x, u, Du) &=& 0, \quad x \in \Omega,\\
u &=& 0, \quad \text{on } \partial\Omega,
\end{array}\right. \tag{$\mathcal{PD}$}
\end{equation}
or Neumann boundary condition:
\begin{equation}\label{PN}
\left\lbrace\begin{array}{rcl}
\operatorname{div} A(x, u, Du) + B(x, u, Du) &=& 0, \quad x \in \Omega,\\
A(x, u, Du)\cdot\nu &=& C(x,u), \quad \text{on } \partial\Omega,
\end{array}\right. \tag{$\mathcal{PN}$}
\end{equation}
where \( \nu(x) \) denotes the outward unit normal vector to \( \partial\Omega \) at \( x \in \partial\Omega \). The functions
\[
A : \Omega \times \mathbb{R} \times \mathbb{R}^n \rightarrow \mathbb{R}^n,\quad
B : \Omega \times \mathbb{R} \times \mathbb{R}^n \rightarrow \mathbb{R},\quad
C : \partial\Omega \times \mathbb{R} \rightarrow \mathbb{R}
\]
are Carathéodory functions satisfying the following \( G(x,t) \)-structure conditions:
\begin{align}
A(x, u, \eta) \cdot \eta &\geq a_4\, G(x,|\eta|) - a_5\, H_1(x,|u|) - a_6, \tag{$\mathcal{A}_1$}\label{6.31}\\
|A(x, u, \eta)| &\leq a_1\, g(x,|\eta|) + a_2\, P(x,|u|) + a_3,  \tag{$\mathcal{A}_2$}\label{6.41}\\
|B(x, u, \eta)| &\leq b_1\, M(x,|\eta|) + b_2\, h_1(x,|u|) + b_3, \tag{$\mathcal{B}$}\label{6.51}\\
|C(x, u)| &\leq c_1\, h_2(x,|u|) + c_2. \tag{$\mathcal{C}$}\label{6.61}
\end{align}
The auxiliary functions \( M(x,t) \) and \( P(x,t) \) are defined as:
\begin{align}
M(x,t) &= \widetilde{H_1}^{-1}(x,G(x,t)), \label{6.71}\\
P(x,t) &= \widetilde{G}^{-1}(x,H_1(x,t)). \label{6.81}
\end{align}
For the definition of $\widetilde{H_1}^{-1}$ see Definitions \ref{definv} and \ref{CF}. The constants \( a_i \), \( b_j \), and \( c_k \) are all positive, and we fix parameters \( F_1 > 1 \) and \( F_2 > 1 \). We also introduce the generalized \textnormal{N}-functions:
\[
H_1(x,t) = \int_0^t h_1(x,s)\, \mathrm{d}s,\qquad H_2(x,t) = \int_0^t h_2(x,s)\, \mathrm{d}s,
\]
which satisfy the growth conditions:
\begin{align}
F_i^{-1} \leq H_i(x,1) \leq F_i, &\qquad \text{for all } x \in \Omega,\quad i=1,2, \tag{$\mathcal{H}_1$} \label{H111}\\
g^+ \leq h_i^- \leq \frac{h_i(x,t)\, t}{H_i(x,t)} \leq h_i^+\leq g_*^-:=\frac{ng^-}{n-g^-}, &\qquad \text{for all } x \in \Omega,\ t > 0,\quad i=1,2, \tag{$\mathcal{H}_2$} \label{H211}
\end{align}
\begin{equation}
\int_{\Omega} H_i(x,t) \, \mathrm{d}x < \infty,\ \text{for } t>0,\quad i=1,2, \tag{$\mathcal{H}_3$}\label{YoungTriple4}
\end{equation}
and
\begin{equation}
G(x,t) \prec H_i(x,t) \prec\prec G^*(x,t),\quad i = 1,2, \tag{$\mathcal{H}_4$}\label{YoungTriple1}
\end{equation}
where the symbol "$\prec\prec$" will be introduced in Definition \ref{prec}, and
\[
G^*(x,t) = \int_0^t g^*(x,s)\, \mathrm{d}s
\]
is the Sobolev conjugate of \( G(x,t) \) (see Subsection \ref{sec2}).

\begin{remark}
The generalized \textnormal{N}-functions \( H_1(x,t) \) and \( H_2(x,t) \) may exhibit critical growth corresponding to the conjugate \( G^*(x,t) \). In such a case, we have \( h_i^+ = g_*^+ := \frac{n g^+}{n - g^+} \) and \( h_i^- = g_*^- := \frac{n g^-}{n - g^-} \) for \( i = 1,2 \).
\end{remark}

\begin{definition}\label{Definition 4.1}
Under assumptions \eqref{D22}--\eqref{GG2}, \eqref{6.31}--\eqref{6.41}, \eqref{6.51} and \eqref{H111}--\eqref{YoungTriple1}, a function \( u \in W^{1,G(x,t)}(\Omega) \) is a \emph{weak solution} of problem~\eqref{PD} if
\begin{equation}\label{4.9}
\int_{\Omega} A(x, u, Du) \cdot Dv \, \mathrm{d}x - \int_{\Omega} B(x, u, Du) v \, \mathrm{d}x = 0, \tag{$\mathcal{FDV}$}
\end{equation}
for all test functions \( v \in W_0^{1,G(x,t)}(\Omega) \).
\end{definition}

\begin{definition}\label{Definition 4.1w}
Under the hypotheses of Definition~\ref{Definition 4.1} and condition~\eqref{6.61}, a function \( u \in W^{1,G(x,t)}(\Omega) \) is said to be a \emph{weak solution} of problem~\eqref{PN} if
\begin{equation}\label{FNV}
\int_{\Omega} A(x, u, Du) \cdot Dv \, \mathrm{d}x - \int_{\partial\Omega} C(x, u)v \, \mathrm{d}\sigma - \int_{\Omega} B(x, u, Du) v \, \mathrm{d}x = 0, \tag{$\mathcal{FNV}$}
\end{equation}
for all \( v \in W^{1,G(x,t)}(\Omega) \), where \( \mathrm{d}\sigma \) denotes the surface measure on \( \partial\Omega \).
\end{definition}

Since the boundedness of weak solutions is essential for establishing their regularity, in particular, Hölder continuity, we first prove the boundedness of weak solutions to problem~\eqref{P} under both Dirichlet and Neumann conditions.

\begin{theorem}\label{Theorem 4.1w}
Assume that conditions \eqref{D22}--\eqref{GG2}, \eqref{6.31}--\eqref{6.41}, \eqref{6.51}, and \eqref{H111}--\eqref{YoungTriple1} hold. If \( u \in W^{1,G(x,t)}(\Omega) \) is a weak solution to problem~\eqref{PD}, then \( u \in L^{\infty}(\Omega) \).
\end{theorem}
\begin{theorem}\label{Theorem 6.1w}
Under the assumptions of Theorem~\ref{Theorem 4.1w} and condition~\eqref{6.61}, any weak solution \( u \in W^{1,G(x,t)}(\Omega) \) of problem~\eqref{PN} satisfies \( u \in L^{\infty}(\Omega) \cap L^{\infty}(\partial\Omega) \).
\end{theorem}
Combining the previous theorems, we establish the Hölder continuity of weak solutions to problem~\eqref{P}, subject to either Dirichlet boundary conditions~\eqref{PD} or Neumann boundary conditions~\eqref{PN}.
\begin{theorem}\label{Theorem 4.2}
Assume that conditions \eqref{D22}--\eqref{GG3}, \eqref{6.31}--\eqref{6.41}, \eqref{6.51} and \eqref{H111}--\eqref{YoungTriple1}, hold. Let \( u \in W^{1,G(x,t)}(\Omega) \) be a weak solution to the problem~\eqref{PD}, and suppose that
\begin{equation}\label{M}
\mathop{\mathrm{ess\,sup}}_{x \in \Omega} |u(x)| \leq M. \tag{$\mathcal{M}$}
\end{equation}
Then \( u \in \mathcal{B}_{G(x,t)}(\overline{\Omega}, M, \gamma, \gamma_1, \delta) \), where
\[
\gamma = \gamma\big(a_2, a_3, a_4, g^-, g^+\big), \quad
\gamma_1 = \gamma_1\big(a_1,a_2, a_3, a_5,a_6, b_1,b_2,b_3, g^+,h^+_1,F_1,M\big), \quad
\delta = \min\left\{ \frac{a_4}{4 b_3M}, 2 \right\}.
\]
\end{theorem}

\begin{theorem}\label{Theorem 4.21}
Assume that conditions of Theorem~\ref{Theorem 4.2} and condition~\eqref{6.61} hold. Let \( u \in W^{1,G(x,t)}(\Omega) \) be a weak solution to problem~\eqref{PN}, and suppose that \eqref{M} holds. Then \( u \in \mathcal{B}_{G(x,t)}(\overline{\Omega}, M, \gamma, \gamma_1, \delta) \), where
\[
\gamma = \gamma\big(a_2, a_3, a_4,g^-, g^+\big), \quad
\gamma_1 = \gamma_1\big(a_1,a_2, a_3, a_5,a_6, b_1,b_2,b_3,c_1,c_2, g^+, h^+_1,F_1,M\big), \quad
\delta = \min\left\{ \frac{a_4}{4 b_3M}, 2 \right\}.
\]
\end{theorem}

\begin{theorem}\label{Theorem 4.3}
Suppose that the assumptions of Theorem~\ref{Theorem 4.2} or Theorem~\ref{Theorem 4.21} are satisfied. Then \( u \in C^{0,\alpha}(\Omega) \), where
\[
\alpha = \alpha\big(a_2, a_3, a_4,M, n, g^-, g^+, F, L\big) \in (0,1].
\]
\end{theorem}

\begin{theorem}\label{Theorem 4.4}
Let \( u \in W^{1,G(x,t)}(\Omega) \) be a weak solution to problem~\eqref{PD} (respectively,~\eqref{PN}), under the assumptions of Theorem~\ref{Theorem 4.2} (respectively, Theorem~\ref{Theorem 4.21}). Then \( u \) is locally Hölder continuous in \( \Omega \). Moreover, if the boundary \( \partial \Omega \) satisfies condition~(A) and the trace \( u|_{\partial \Omega} \) is Hölder continuous, then \( u \) is Hölder continuous up to the boundary, i.e., \( u \in C^{0,\alpha}(\overline{\Omega}) \) for some \( \alpha \in (0,1] \).
\end{theorem}
%----------------
\subsection{Examples}
The following discussion of customary generalized \textnormal{N}-functions satisfying \eqref{D22}--\eqref{GG3} serves to illustrate the range of problems covered by our theorems.


\begin{example}{\rm\bf{[\(p\)-growth conditions]}} \label{ex1}
Let \(p \in (1, n)\). Consider the case where \(G(x,t) = |t|^p\). In this setting, \(G(x,t)\) is a generalized \textnormal{N}-function that satisfies the conditions \eqref{D22}--\eqref{GG3}. Consequently, the main results established in this work apply to the classical \(p\)-growth case, which has already been extensively investigated by Lady{\v{z}}enskaja-Solonnikov-Ural{\cprime}ceva \cite{OL} and G. Lieberman~\cite{Li1988}.
\end{example}

\begin{example}{\rm\bf{[\(G(t)\)-growth conditions]}} \label{ex2}
A more general framework than the classical \(p\)-growth is obtained by considering a generalized \textnormal{N}-function \(G(x,t)\) that does not depend on the spatial variable \(x\). In this case, we write \(G(t)\), and the problem reduces to the so-called Orlicz case. If 
\[ 
1 < g^- \leq \frac{g(t)t}{G(t)} \leq g^+ <n,\ \text{for all}\ t>0, 
\]
then, the function \(G(t)\) still satisfies the assumptions \eqref{D22}--\eqref{GG3}, thereby falling within the scope of our results. Various aspects of this case have been previously studied in depth, notably by G. Lieberman~\cite{Li1991}.
\end{example}

\begin{example}{\rm\bf{[\(p(x)\)-growth conditions]}} \label{ex3}
Let \(p:\Omega \rightarrow \mathbb{R}\) be a measurable function such that
\begin{equation}\label{var111}
1 < p^- := \inf_{x \in \overline{\Omega}} p(x) \leq p^+ := \sup_{x \in \overline{\Omega}} p(x) < \infty.
\end{equation}
Assume that \(p\) satisfies the log-Hölder continuity condition:
\begin{equation}\label{var1}
    \left|\frac{1}{p(x)} - \frac{1}{p(y)}\right| \leq \frac{c}{\ln\left(e + \frac{1}{|x - y|}\right)} \quad \text{for all } x, y \in \overline{\Omega},
\end{equation}
for some constant \(c > 0\). Then the function
\[
G(x,t) = |t|^{p(x)}
\]
defines a generalized \textnormal{N}-function that satisfies the structural assumptions \eqref{D22}--\eqref{GG3}; see \cite[Section 7.1]{Harjulehto2019}. Therefore, our results also apply to the variable exponent setting, which has been extensively studied by X. Fan and D. Zhao~\cite{Fan1999}.
\end{example}
\begin{example}{\rm\bf{[Logarithmic growth conditions]}} \label{ex4}
Let \(p: \Omega \rightarrow \mathbb{R}\) be a measurable function satisfying conditions \eqref{var111}--\eqref{var1}. Then the function
\[
G(x,t) = |t|^{p(x)} \ln(e + |t|)
\]
defines a generalized \textnormal{N}-function that satisfies the structural conditions \eqref{D22}--\eqref{GG3}. This example extends the case of the classical variable exponent by incorporating a logarithmic perturbation that reflects more refined material responses in heterogeneous media. This type of growth has recently attracted attention in the context of non-standard regularity theory and is encompassed by our main results.
\end{example}
\begin{example}{\rm\bf{[\((p(x),q(x))\)-growth conditions]}} \label{ex6}
Let \(p, q: \Omega \rightarrow \mathbb{R}\) be measurable functions such that
\begin{equation}\label{var20}
p,q\in C^{0,1}(\overline{\Omega}),
\end{equation}
\begin{equation}\label{var21}
1<p(x)<n\quad \text{and}\quad p(x)<q(x)\quad \text{for all}\ x \in \overline{\Omega},
\end{equation}
\begin{equation}\label{var22}
\frac{q^+}{p^-}<1+\frac{1}{n}, \quad \text{where}\quad p^- := \inf_{x \in \overline{\Omega}} p(x)\ \text{and}\ q^+ := \sup_{x \in \overline{\Omega}} q(x),
\end{equation}
and
\begin{equation}\label{double22}
   q(x)-p(x)=c>0,\ \text{for all}\ x\in \overline{\Omega}.
\end{equation}
Under these conditions, the function
\[
G(x,t) = |t|^{p(x)} +|t|^{q(x)}
\]
defines a generalized \textnormal{N}-function satisfying assumptions \eqref{D22}--\eqref{GG3}. Consequently, our results also apply to the variable exponent double phase growth setting, which has recently been the focus of active research.
\end{example}
\begin{example}{\rm\bf{[\((p,q)\)-type growth conditions]}} \label{ex5}
Let \(1 < p < q < \infty\) and let \(a \in L^\infty(\Omega)\) be a positive function satisfying
\begin{equation}\label{double2}
    q \leq \frac{(n+1)p}{n},
\end{equation}
as well as the Hölder continuity condition
\begin{equation}\label{double3}
    \frac{1}{a} \in C^{0, \frac{n}{p}(q - p)}(\Omega).
\end{equation}
In particular, the following estimate holds:
\begin{equation}\label{double5}
    \frac{1}{a(x)} \leq  \frac{1}{a(y)} + |x - y|^{\frac{n}{p}(q - p)}  \quad \text{for all } x, y \in \Omega,
\end{equation}
for some constant \(c > 0\). Under these conditions, the function
\[
G(x,t) = |t|^p + a(x)|t|^q
\]
defines a generalized function \textnormal{N} satisfying the assumptions \eqref{D22}--\eqref{GG3}. Hence, our results also cover the semi-double phase growth case. Unfortunately, our results do not cover the classical double phase growth case, which has attracted significant attention in recent years (see ~\cite{Colombo2015}).
\end{example}
\begin{example}{\rm\bf{[Logarithmic \((p,q)\)-type growth conditions]}} \label{ex8}
Let \(1 < p < q < \infty\) and let \(a \in L^\infty(\Omega)\) be a positive function satisfying the conditions \eqref{double2}--\eqref{double5}. Then the function
\[
G(x,t) = |t|^{p} + a(x)|t|^{q} \ln(e + |t|)
\]
defines a generalized \textnormal{N}-function that satisfies the structural conditions \eqref{D22}--\eqref{GG3}. This example extends the classical double-phase case by incorporating a logarithmic perturbation in the higher-order term, which reflects more refined material responses in heterogeneous media. This type of growth has recently attracted attention in the context of non-standard regularity theory and is encompassed by our main results.
\end{example}
Following the same structure, many additional examples of growth conditions can be generated that fall within the scope of our results.\\

The following remark clarifies the logical relationship between assumptions
\eqref{GG2} and \eqref{GG3}. Although both conditions impose continuity-type
requirements on the generalized $\mathrm{N}$-function $G(x,t)$ with respect to the
spatial variable, they control fundamentally different aspects of the growth.
Condition \eqref{GG2} is a local comparability assumption, which is mainly used to
derive functional-analytic tools such as Poincaré and Sobolev-type inequalities in
Musielak--Orlicz spaces. In contrast, condition \eqref{GG3} provides a quantitative
control on the spatial oscillation of the effective growth exponent
$\frac{t g(x,t)}{G(x,t)}$, which plays a crucial role in De Giorgi--type iteration
arguments. The remark below shows, by means of explicit examples, that these two
assumptions are independent in general.

\begin{remark}\label{RemH}
Conditions \eqref{GG2} and \eqref{GG3} are independent in general.

\smallskip
\noindent
\textbf{(1) \eqref{GG2} does not imply \eqref{GG3}.}
Consider the variable exponent case
\[
G(x,t)=t^{p(x)},
\]
where $p:\Omega\to\mathbb{R}$ is bounded and measurable but not log-Hölder continuous.
Then condition \eqref{GG2} holds (see, e.g., \cite{Harjulehto2019}), since
$G^{-1}(x,t)=t^{1/p(x)}$ is locally comparable on balls. However,
\[
\frac{t g(x,t)}{G(x,t)}=p(x),
\]
and thus condition \eqref{GG3} fails whenever $p(\cdot)$ lacks log-Hölder
continuity. Hence \eqref{GG2} does not imply \eqref{GG3}.

\smallskip
\noindent
\textbf{(2) \eqref{GG3} does not imply \eqref{GG2}.}
Let
\[
G(x,t)=t^p(1+a(x)),
\quad p>1,
\]
where $a:\Omega\to[0,\infty)$ is unbounded or highly oscillatory.
Then
\[
\frac{t g(x,t)}{G(x,t)}= p,
\]
so condition \eqref{GG3} holds trivially. However, the lack of local comparability
of $G(x,t)$ prevents \eqref{GG2} from being satisfied in general.

\smallskip
Therefore, \eqref{GG2} and \eqref{GG3} are logically independent assumptions.
\end{remark}


%----------------
\subsection{Structure of the paper}
The paper is organized as follows. In Section~2, we introduce notations and recall the fundamental definitions and properties of Musielak--Orlicz and Musielak--Orlicz--Sobolev spaces. 
Section~3 presents a new De Giorgi-type class and establishes key regularity results for this class, which provide the theoretical foundation for our main results, Theorems~\ref{Theorem 2.1} and~\ref{Theorem 2.2}. 
Building on this framework, Section~4 proves the boundedness of weak solutions to problem~\eqref{P} in both subcritical and critical cases (Theorems~\ref{Theorem 4.1w} and~\ref{Theorem 6.1w}). 
Furthermore, using the regularity results for the De Giorgi-type class, we establish the $C^{0,\alpha}$ regularity of weak solutions (Theorems~\ref{Theorem 4.2}--\ref{Theorem 4.4}).


%----------------
\section{Preliminaries}
Here, we introduce the fundamental notations  and definitions used throughout the paper. A central concept is the Musielak-Orlicz-Sobolev space, whose properties are essential for the results that follow.
\subsection{Notations}
The following notations will be used throughout this paper:
\begin{itemize}
    \item $\mathbb{R}^n$ denotes the $n$-dimensional Euclidean space.
    \item $\Omega \subset \mathbb{R}^n$ is a bounded domain.
    \item For a measurable set $E \subset \mathbb{R}^n$, $\mathrm{mes}(E)$ or $|E|$ denotes its $n$-dimensional Lebesgue measure.
    \item If $u$ is a measurable function defined on $\Omega$, and $E \subset \Omega$ is measurable, we define:
   \begin{equation}\label{os1}
    \max_E u := \mathop{\mathrm{ess\,sup}}_{x \in E} u(x), \quad
    \min_E u := \mathop{\mathrm{ess\,inf}}_{x \in E} u(x), \quad
    \mathrm{osc}_E u := \max_E u - \min_E u.
    \end{equation}
    \item For any $x_0 \in \mathbb{R}^n$ and $\rho > 0$, we define the open ball:
    \[
    B_\rho(x_0) := \{x \in \mathbb{R}^n : |x - x_0| < \rho\},
    \]
    which we may denote simply by $B_\rho$ when the center is clear from context. Balls $B_{\rho_1}$ and $B_{\rho_2}$ are assumed to be concentric. We denote by $\omega_n := |B_1|$ the measure of the unit ball in $\mathbb{R}^n$.
    \item For $0 < \alpha \leq 1$, we denote:
    \[
    C^{0,\alpha}(\overline{\Omega}) := \left\{ u : u \text{ is Hölder continuous on } \overline{\Omega} \text{ with exponent } \alpha \right\},
    \]
    and
    \[
    C^{0,\alpha}(\Omega) := \left\{ u : u \in C^{0,\alpha}(\overline{\Omega'}) \text{ for all } \Omega' \Subset \Omega \right\}.
    \]
    \item For any function \( f : \Omega \rightarrow \mathbb{R} \), we define the positive part and the negative part of \( f \) as:
\[
f^+ = \max\{f, 0\}, \quad f^- = \max\{-f, 0\}.
\]

\end{itemize}
\subsection{Musielak-Orlicz-Sobolev space \texorpdfstring{$W^{1,G(x,t)}$}{W1G(x,t)}}\label{sec2}
In this subsection, we recall some definitions and fundamental properties of Musielak-Orlicz and Musielak-Orlicz-Sobolev spaces. 
For a comprehensive bibliography on Musielak-Orlicz-Sobolev spaces, we refer the reader to \cite{Ala2, TD, Chlebicka2021, Diening2011, Harjulehto2019, Fan2012a, Fan2012b, Musielak1983, TS1, TS11, TS12, TS2, TS3, TS4, TS5, TS6, TS7, TS8, TS9, TS10}.

%-----------------------------
\begin{definition}
Let $\Omega$ be an open subset of $\mathbb{R}^n$. A function $G:\Omega \times \mathbb{R} \to \mathbb{R}$ is called a generalized \textnormal{N}-function if it satisfies the following conditions:
\begin{enumerate}
    \item[(1)] For a.e. $x\in \Omega$, the function $G(x,t)$ is even, continuous, strictly increasing, and convex in $t$, and for each $t\in \mathbb{R}$, $G(x,t)$ is measurable in $x$;
    \item[(2)] $\displaystyle\lim_{t \to 0} \frac{G(x,t)}{t} = 0$, for a.e. $x\in \Omega$;
    \item[(3)] $\displaystyle\lim_{t \to \infty} \frac{G(x,t)}{t} = \infty$, for a.e. $x\in \Omega$;
    \item[(4)] $G(x,t) > 0$ for all $t > 0$ and all $x\in \Omega$, and $G(x,0) = 0$ for all $x\in \Omega$.
\end{enumerate}
\end{definition}

\begin{remark}[Remark~2.1, \cite{Ala1}]
We give an equivalent definition of a generalized \textnormal{N}-function that admits an integral representation. For $x \in \Omega$ and $t \geq 0$, let $g(x,t)$ be the right-hand derivative of $G(x,\cdot)$ at $t$, and define $g(x,t) = -g(x,-t)$ for $t < 0$. Then for each $x \in \Omega$, the function $g(x,\cdot)$ is odd, real-valued, satisfies $g(x,0) = 0$, and $g(x,t) > 0$ for $t > 0$, is right-continuous and nondecreasing on $[0, +\infty)$, and
\[
G(x,t) = \int_0^{|t|} g(x,s)\, \mathrm{d}s, \quad \text{for all } x\in \Omega, \ t\in \mathbb{R}.
\]
\end{remark}
\begin{definition}\label{definv}
Let $G:\Omega\times[0,\infty)\to[0,\infty)$ be a generalized $\mathrm{N}$-function.
Since for a.e.\ $x\in\Omega$ the mapping $t\mapsto G(x,t)$ is continuous and
strictly increasing on $[0,\infty)$, we can define $G^{-1}(x,\cdot)$ as the
continuous inverse of $G(x,\cdot)$. In particular, it holds that
\[
G^{-1}(x,G(x,t)) = G(x,G^{-1}(x,t)) = t,
\quad \text{for a.e.\ } x\in\Omega \text{ and all } t\ge 0.
\]
\end{definition}

\begin{definition}
A generalized \textnormal{N}-function $G$ satisfies the \emph{$\Delta_2$-condition} if there exist $C_0 > 0$ and a nonnegative function $\varphi \in L^1(\Omega)$ such that
\[
G(x,2t) \leq C_0 G(x,t) + \varphi(x), \quad \text{for a.e. } x\in \Omega \text{ and all } t \geq 0.
\]
\end{definition}
\begin{definition}
A generalized \textnormal{N}-function $G(x,t)$ is said to satisfy:
\begin{enumerate}
    \item[\textnormal{(A0)}] if there exists $\mu \in (0,1]$ such that
    \[
    \mu \leq G^{-1}(x,1) \leq \frac{1}{\mu}, \quad \text{for a.e. } x \in \mathbb{R}^n;
    \]
    
    \item[\textnormal{(A1)}] if there exists $\mu \in (0,1]$ such that
    \[
    \mu G^{-1}(x,t) \leq G^{-1}(y,t),
    \]
    for every $t \in [1, 1/|B|]$, for a.e. $x,y \in B$, and every ball $B \subset \mathbb{R}^n$ with $|B| \leq 1$;
    
    \item[\textnormal{(A2)}] if for every $s > 0$ there exist $\mu \in (0,1]$ and $\phi \in L^1(\mathbb{R}^n) \cap L^\infty(\mathbb{R}^n)$ such that
    \[
    \mu G^{-1}(x,t) \leq G^{-1}(y,t),
    \]
    for a.e. $x,y \in \mathbb{R}^n$ and for all $t \in [\phi(x)+\phi(y), s]$.
\end{enumerate}
\end{definition}


\begin{definition}\label{prec}
Let $G_1(x,t)$ and $G_2(x,t)$ be two generalized \textnormal{N}-functions.
\begin{itemize}
    \item[(1)] We say that $G_1(x,t)$ increases essentially slower than $G_2(x,t)$ near infinity, and we write $G_1 \prec\prec G_2$, if for any $k > 0$
    \[
    \lim_{t \to \infty} \frac{G_1(x, kt)}{G_2(x,t)} = 0, \quad \text{uniformly in } x \in \Omega.
    \]
    \item[(2)] We say that $G_1(x,t)$ is \emph{weaker} than $G_2(x,t)$, denoted by $G_1 \prec G_2$, if there exist constants $C_1, C_2 > 0$ and a nonnegative function $\varphi \in L^1(\Omega)$ such that
    \[
    G_1(x,t) \leq C_1 G_2(x, C_2 t) + \varphi(x), \quad \text{for a.e. } x \in \Omega \text{ and all } t \geq 0.
    \]
\end{itemize}
\end{definition}

\begin{definition}[Complementary function]\label{CF}
Let $G:\Omega\times[0,\infty)\to[0,\infty)$ be a generalized $\mathrm{N}$-function.
The function $\widetilde{G}:\Omega\times[0,\infty)\to[0,\infty)$ defined by
\begin{equation}\label{Cf}
\widetilde{G}(x,t)
:= \sup_{s \ge 0} \bigl\{ ts - G(x,s) \bigr\},
\quad \text{for a.e.\ } x\in\Omega \text{ and all } t\ge 0,
\end{equation}
is called the \emph{complementary function} (or \emph{conjugate function}) of $G$.
Moreover, $\widetilde{G}$ is also a generalized $\mathrm{N}$-function.
\end{definition}


\begin{remark}
From the definition of the complementary function $\widetilde{G}(x,t)$, we derive the following Young-type inequality:
\begin{equation}\label{Yi}
st \leq G(x,s) + \widetilde{G}(x,t), \quad \text{for all } x \in \Omega, \ s,t \geq 0.
\end{equation}
\end{remark}

The next lemma is taken from Bahrouni--Bahrouni--Missaoui \cite[Lemma~2.3]{Ala1}.

\begin{lemma}\label{lm1}
Let $G(x,t)$ be a generalized \textnormal{N}-function. Suppose that $t \mapsto g(x,t)$ is continuous and increasing on $\mathbb{R}$ for a.e. $x \in \Omega$. Moreover, assume that there exist constants $g^-, g^+ \in \mathbb{R}$ such that
\begin{align}
\widetilde{G}(x, g(x,s)) &\leq (g^+ - 1) G(x,s), \quad \text{for all } s \geq 0, \ x \in \Omega, \label{L1} \end{align}
and
\begin{align}
\frac{g^+}{g^+ - 1} =: \widetilde{g^-} &\leq \frac{\widetilde{g}(x,s)s}{\widetilde{G}(x,s)} \leq \widetilde{g^+} := \frac{g^-}{g^- - 1}, \quad \text{for all } x \in \Omega, \ s > 0, \label{D3}
\end{align}
where $\displaystyle \widetilde{G}(x,s) = \int_0^s \widetilde{g}(x,t)\, dt$.
\end{lemma}

\begin{remark}\label{compl}
The condition \eqref{D22} implies that $G(x,t)$ and its complementary function $\widetilde{G}(x,t)$ satisfy the $\Delta_2$-condition.
\end{remark}
 Now, we  define the Musielak-Orlicz space as follows:
$$L^{G(x,t)}(\Omega):=\left\lbrace u:\Omega\longrightarrow \mathbb{R}\ \text{measurable :}\ \rho_{G}(\lambda u)<+\infty,\ \ \text{for some}\ \lambda>0\right\rbrace,$$
where
\begin{equation}\label{Mo}
  \rho_{G}(u):= \int_{\Omega}G(x, u){\rm d}x.
\end{equation}
The space $L^{G(x,t)}(\Omega)$ is endowed with the Luxemburg norm
\begin{equation}\label{No}
  \Vert u\Vert_{L^{G(x,t)}(\Omega)}:=\inf\left\lbrace \lambda:\ \rho_{G}\left(x,\frac{u}{\lambda}\right)\leq 1\right\rbrace.
\end{equation}
\begin{proposition}\label{HM}
    Let $G$ be a generalized \textnormal{N}-function satisfies the $\Delta_2$-condition, then
    $$L^{G(x,t)}(\Omega) =\left\lbrace  u: \Omega \longrightarrow \mathbb{R}\ \text{measurable :}\ \rho_{G}( u)<+\infty\right\rbrace.$$
\end{proposition}
\begin{proposition} \label{zoo}
     Let $G$ be a generalized \textnormal{N}-function satisfies \eqref{D22}, then the following assertions hold:
     \begin{itemize}
    \item [(1)] $\min \{s^{g^-}, s^{g^+}\}G(x,t)\leq  G(x,s t)\leq \max \{ s^{g^-}, s^{g^+}\}G(x,t),\text{ for a.e. } x \in \Omega$ $\text{ and all }  s, \ t \geq 0. $
     \item [(2)] $\min \{s^{\widetilde{g^-}}, s^{\widetilde{g^+}}\}\widetilde{G}(x,t)\leq  \widetilde{G}(x,s t)\leq \max \{ s^{\widetilde{g^-}}, s^{\widetilde{g^+}}\}\widetilde{G}(x,t),\text{ for a.e. } x \in \Omega$ $\text{ and all }  s, \ t \geq 0. $
 \item [(3)]  $\min \left\{\|u\|_{L^{G(x,t)}(\Omega)}^{g^-},\|u\|_{L^{G(x,t)}(\Omega)}^{g^+}\right\} \leq \rho_{G}(u) \leq\max \left\{\|u\|_{L^{G(x,t)}(\Omega)}^{g^-},\|u\|_{L^{G(x,t)}(\Omega)}^{g^+}\right\},$ \ for all $  u\in L^{G(x,t)}(\Omega)$.
 %\item [(4)] Let $\left\{ u_n\right\}_{n \in \mathbb{N}}\subseteq L^{G(x,t)}(\Omega) $ and $  u \in L^{G(x,t)}(\Omega) $, then
% $$
% \|u_n-u\|_{L^{G(x,t)}(\Omega) }\longrightarrow0 \ \Longleftrightarrow  \rho_G (u_n-u) \longrightarrow 0,\ \text{ as } n \rightarrow +\infty. $$
\end{itemize}
\end{proposition}
  %\begin{lm}
   %  Let $\h$ be a generalized \textnormal{N}-function. Then, we have
    % $$ \Vert u\Vert_{\h}\leq \rho_{\h}(u)+1,\ \ \text{for all}\ u\in L^{\h}(\Omega).$$
%\end{lm}
As a consequence of \eqref{Yi}, we have the following result:
\begin{lemma}[H\"older's type
inequality]\label{H1}
  Let $\Omega$ be an open subset of $\mathbb{R}^n$ and $G$ be a generalized \textnormal{N}-function satisfies \eqref{D22}, then
  \begin{equation}\label{Ho}
     \left\vert \int_{\Omega} uv {\rm d}x \right\vert \leq 2 \Vert u\Vert_{L^{G(x,t)}(\Omega)}\Vert v\Vert_{L^{\widetilde{G}(x,t)}(\Omega)},\ \text{for all}\ u\in L^{G(x,t)}(\Omega)\ \text{and all}\ v\in L^{\widetilde{G}(x,t)}(\Omega).
  \end{equation}
\end{lemma}
The subsequent proposition deals with some topological properties of
the Musielak-Orlicz space, see \cite[Theorem 7.7 and Theorem
8.5]{Musielak1983}.
\begin{proposition}\label{AB}\ \\
\begin{enumerate}
\item [(1)] Let $G(x,t)$ be a generalized \textnormal{N}-function and $\Omega$ an open subset of
 $\mathbb{R}^n$. Then,
 \begin{enumerate}
     \item[(a)] the space
    % $L^{\h}(\Omega)$ is complete with respect to the norm $\|\cdot\|_{\lh}$, that is,
    $\left(L^{G(x,t)}(\Omega),\|\cdot\|_{G}\right)$ is a Banach space;
     \item[(b)] if $G(x,t)$ satisfies \eqref{D22}, then
$L^{G(x,t)}(\Omega)$ is a separable and reflexive space.
 \end{enumerate}
\item [(2)] Let $G_1(x,t)$ and $G_2(x,t)$  be two generalized \textnormal{N}-functions such that $G_1(x,t) \prec G_2(x,t)$ and $\Omega$ be an open bounded subset of $\mathbb{R}^n$. Then,
$$
L^{G_2(x,t)}(\Omega) \hookrightarrow L^{G_1(x,t)}(\Omega).
$$
\end{enumerate}
\end{proposition}

Now, we are ready to define the  Musielak-Orlicz Sobolev space. Let
$G(x,t)$ be a generalized \textnormal{N}-function and $\Omega$ be an open subset of
$\mathbb{R}^n$. The  Musielak-Sobolev space is defined as follows
$$W^{1,G(x,t)}(\Omega):=\left\lbrace u\in L^{G(x,t)}(\Omega):\ |D u| \in L^{G(x,t)}(\Omega)\right\rbrace.$$
The space $W^{1,G(x,t)}(\Omega)$ is endowed with the norm
\begin{equation}\label{NM}
  \Vert u\Vert:=\Vert u\Vert_{L^{G(x,t)}(\Omega)}+\Vert D u\Vert_{L^{G(x,t)}(\Omega)},\ \ \text{for all}\ u\in W^{1,G(x,t)}(\Omega),
\end{equation}
where $\|D u\|_{L^{G(x,t)}(\Omega)} := \| |D u| \|_{L^{G(x,t)}(\Omega)}$.\\ We denote by
$W^{1,G(x,t)}_0(\Omega)$
 the completion of $C^\infty _0(\Omega)$ in $W^{1,G(x,t)}(\Omega)$.
%\begin{theorem}\label{thc}\cite[Theorem 10.2] {Musielak1983}\cite[Proposition 1.7]{Fan2012b}~~~~~~~~~~~~~~~~~~~~~~~~~~~~~~~~~~~~~~~~~~~~~~~~~~~~~~~~~~~~~~

% Let $G(x,t)$ be a generalized \textnormal{N}-function that is locally integrable such that
%\begin{equation}\label{c1c}
%\inf _{x \in \Omega} G(x, 1)>0.
%\end{equation}
%Then, the spaces $W^{1, G(x,t)}(\Omega)$ and $W_0^{1, G(x,t)}(\Omega)$ are
%%reflexive.
%\end{theorem}
%\begin{theorem}\cite[Theorem 10.3]{Musielak1983}~~~~~~~~~~~~~~~~~~~~~~~~~~~~

 %Let $\h_1, \h_2 $ be locally integrable functions such that
 %$$
 %\inf _{x \in \Omega} \h_1(x, 1)>0 \text{ and } \inf _{x \in \Omega} \h_2(x, 1)>0.
 %$$
 %Moreover, if $\h_1 \preceq \h_2$, then the following continuous embedding hold
 %$$
 %W^{1, \h_2}(\Omega) \hookrightarrow W^{1, \h_1}(\Omega).
 %$$
% \end{theorem}
\begin{remark}\label{Ref}
 If $G$  satisfies \eqref{D22}, then the space
$W^{1,G(x,t)}(\Omega)$ is a reflexive and separable Banach space with
respect to the norm $\Vert \cdot\Vert$.
\end{remark}

Now we introduce the Sobolev conjugate from \cite{Ala3}, which refines the one proposed in \cite{Cianchi2024} and yields several important results presented below.

\begin{definition}
The Sobolev conjugate of $G(x,t)$ is defined as the generalized \textnormal{N}-function $G^\ast(x,t)$ given by
			\begin{equation*}
  G ^ \ast (x,t)=G (x, N^{-1} (x,t)) , \text{ for a.e. } x \in \mathbb{R}^n \text{ and all } t \geq 0,
\end{equation*}
where
\begin{equation*}
  N(x,t)= \left( \int_{0}^{t} \left( \frac{\tau}{\Phi (x, \tau)} \right)^{\frac{1}{n-1}} d \tau \right)^{\frac{n-1}{n}}, \text{ for } x \in \mathbb{R}^n \text{ and } t\geq 0.
\end{equation*}
 \end{definition}
 \begin{proposition}[\cite{Ala3}] \label{zoo*}
     Let $G(x,t)$ be a generalized \textnormal{N}-function that satisfies \eqref{D22}--\eqref{GG2} and $G^*(x,t)$ its Sobolev conjugate function, then the following assertions hold:
     \begin{itemize}
    \item [(1)] $\min \left\{s^{g^-_*}, s^{g^+_*}\right\}G^*(x,t)\leq  G^*(x,s t)\leq \max \left\{ s^{g^-_*}, s^{g^+_*}\right\}G^*(x,t),\text{ for a.e. } x \in \Omega$ $\text{ and all }  s, \ t \geq 0$;
 \item [(2)] $\displaystyle{
        1<g^-_*\leq \frac{g^*(x,t)t}{G^*(x,t)}\leq g^+_*,\ \ \text{for all}\ x\in \Omega\ \text{and all}\ t> 0}$;
    \item[(3)] $\min \left\{\|u\|_{L^{G^*(x,t)}(\Omega)}^{g^-_*},\|u\|_{L^{G^*(x,t)}(\Omega)}^{g^+_*}\right\} \leq \rho_{G^*}(u) \leq\max \left\{\|u\|_{L^{G^*(x,t)}(\Omega)}^{g^-_*},\|u\|_{L^{G^*(x,t)}(\Omega)}^{g^+_*}\right\},$ \ for all $  u\in L^{G^*(x,t)}(\Omega)$;
\end{itemize}
where $\displaystyle{g_*^- := \frac{ng^-}{n - g^-}}$, $\displaystyle{g_*^+ := \frac{ng^+}{n - g^+}}$ and $\displaystyle{G^*(x,t)=\int_0^t g^*(x,s){\rm d}s}$.
\end{proposition}
\begin{proposition}[\cite{Cianchi2024}]\label{embb}
Let $\Omega$ be a bounded domain in $\mathbb{R}^n$ and $G(x,t)$ be a generalized \textnormal{N}-function that satisfies \eqref{GG1}--\eqref{GG2}. Let $\vartheta(x,t)$ be a generalized \textnormal{N}-function such that
\begin{equation}
\vartheta(x,t)\prec \prec G^*(x,t),\text{  and   } \int_{\Omega} \vartheta(x,t) \, \mathrm{d}x < \infty,\ \text{for } t>0.
\end{equation}
\begin{enumerate}
\item[(1)] The embedding
\begin{equation}
W^{1,G(x,t)}_0(\Omega) \hookrightarrow L^{\vartheta(x,t)}(\Omega) 
\end{equation}
is compact.

\item[(2)] Assume, in addition, that $\Omega$ is a bounded domain with a Lipschitz boundary $\partial \Omega$. Then, the embedding
\begin{equation}
W^{1,G(x,t)}(\Omega) \hookrightarrow L^{\vartheta(x,t)}(\Omega) 
\end{equation}
is compact.
\end{enumerate}
\end{proposition}
\begin{remark}
If the function $G(x,t)$ satisfies assumptions \eqref{D22}--\eqref{GG2}, then it also fulfills conditions (A0)--(A1) as stated in L. Diening and A. Cianchi paper~\cite{Cianchi2024}, and it satisfies condition (A2) due to the boundedness of $\Omega \subset \mathbb{R}^n$. Consequently, the embedding result given in Proposition~\ref{embb} remains valid under assumptions \eqref{D22}--\eqref{GG2}.
\end{remark}
%----------------------------
\section{A New Class of De Giorgi-Type Functions \texorpdfstring{$\mathcal{B}_{G(x,t)}$}{$B_{G(x,t)}$}}

In this section, we introduce and analyze a new De Giorgi-type class of functions denoted by \( \mathcal{B}_{G(x,t)} \). Our main results are encapsulated in Theorems~\ref{Theorem 2.1} and~\ref{Theorem 2.2}, which establish key regularity properties for functions belonging to this class.

The proof strategies we adopt are inspired by the foundational works of X.~Fan~\cite{Fan1999} and G.~Lieberman~\cite{Li1991}. However, our approach involves significant technical refinements that are essential to overcome the challenges arising from the nonstandard growth behavior and spatial heterogeneity inherent in the Musielak–Orlicz framework. In particular, the interaction between Orlicz-type growth and variable exponent dependence introduces substantial analytical difficulties, which require new ideas and careful adaptations of classical methods.

These extensions allow us to unify and generalize existing results under a set of natural and flexible assumptions, thereby providing a broader and more robust regularity theory applicable to a wide range of problems.

We begin by collecting several auxiliary lemmas that will play a crucial role in the subsequent analysis.

\begin{lemma}[see Lemma 4.7 of Chap. II of \cite{OL}]\label{Lemma 2.1}
Let $\{Y_h\}_{h\geq 0}$ be a sequence of non-negative numbers satisfying the recursion relation
\[
Y_{h+1} \leq c b^h Y_h^{1+\varepsilon}, \quad h = 0, 1, 2, \dots,
\]
where $c > 0$, $b > 1$, and $\varepsilon > 0$. If
\[
Y_0 \leq \theta := c^{-\frac{1}{\varepsilon}} b^{-\frac{1}{\varepsilon^2}},
\]
then
\[
Y_h \leq \theta b^{-\frac{h}{\varepsilon}}, \quad \forall h =0,1,2\ldots,
\]
and hence $Y_h \to 0$ as $h \to \infty$.
\end{lemma}

\begin{lemma}[see Lemma 3.5 of Chap. II of \cite{OL}]\label{Lemma 2.2}
For any \(u \in W^{1,1}(B_\rho)\) and arbitrary numbers \(k\) and \(l\) with \(l > k\), the following inequality hol{\rm d}s:
\begin{equation}\label{2.2}
(l - k)|A_{l,\rho}|^{1-\frac{1}{n}} \leq \beta \frac{\rho^n}{|B_\rho \setminus A_{k,\rho}|} \int_{A_{k,\rho}\setminus A_{l,\rho}} |Du| \, {\rm d}x,
\end{equation}
where \(A_{k,\rho} = \{x \in B_\rho : u(x) > k\}\), \(\beta = \beta(n) > 1\) is a constant depending only on \(n\).
\end{lemma}

\begin{lemma}[see Lemma 2.3 of \cite{Fan1999} or Lemma 4.8 of Chap. II of \cite{OL}]\label{Lemma 2.3}
Suppose a function \(u(x)\) is measurable and bounded in some ball \(B_{R_0}\). Consider balls \(B_R\) and \(B_{bR}\) which have a common center with \(B_{R_0}\), where \(b > 1\) is a fixed constant, and suppose that for any \(0 < R \leq b^{-1} R_0\), at least one of the following two inequalities is valid:
\[
\operatorname{osc}_{B_R}u \leq c_1 R^\varepsilon, \quad \text{or}\quad \operatorname{osc}_{B_R}u \leq \theta \operatorname{osc}_{ B_{bR}}u,
\]
where \(c_1, \varepsilon \leq 1\) and \(\theta < 1\) are positive constants. Then \(u \in C^{0,\alpha}(B_{R_0})\), where \(\alpha = \min\{\varepsilon, -\log_b \theta\}\).
\end{lemma}
The following lemmas are fundamental to our subsequent analysis.

\begin{lemma} \label{lem22}
Under assumptions \eqref{D22}--\eqref{GG1} and \eqref{GG3}, for any $x_0 \in \Omega$, $t > 1$, and $R \in (0,1)$, define
\begin{align*}
f^-(R) &:= \inf_{x \in B_R(x_0)} r(x, t), \\
f^+(R) &:= \sup_{x \in B_R(x_0)} r(x, t).
\end{align*}
Then the following properties hold.
\begin{enumerate}
\item[$(1)$] For $t>1$, $$\operatorname{osc}_{B_R}r(x,t)=f^+(R)-f^-(R)=\eta(R).$$
    \item[$(2)$] For all $x \in B_R(x_0)$ and $t > 1$,
          \[
          g^- \leq f^-(R) \leq r(x, t) \leq f^+(R) \leq g^+.
          \]
    \item[$(3)$] For all $x \in B_R(x_0)$ and $t > 1$,
          \[
          F^{-1} t^{f^-(R)} \leq G(x, t) \leq F t^{f^+(R)}.
          \]
    \item[$(4)$] For all sufficiently small $R > 0$,
          \[
          R^{f^-(R) - f^+(R)} \leq L, \quad \text{where } L := \exp(L_0).
          \]
    \item[$(5)$] The oscillation decays to zero:
          \[
          \lim_{R \to 0^+} \left( f^+(R) - f^-(R) \right) = 0.
          \]
\end{enumerate}
\end{lemma}

\begin{proof}
\textbf{Proof of (1):} Let $t>1$. We first show that $$\eta(R) \leq f^+(R) - f^-(R).$$
For any $x, y \in B_R(x_0)$, we have
\[
r(x,t) \leq f^+(R) \quad \text{and} \quad r(y,t) \geq f^-(R),
\]
so that
\[
r(x,t) - r(y,t) \leq f^+(R) - f^-(R).
\]
Similarly,
\[
r(y,t) - r(x,t) \leq f^+(R) - f^-(R),
\]
and hence
\[
|r(x,t) - r(y,t)| \leq f^+(R) - f^-(R).
\]
Taking the supremum over all $x, y \in B_R(x_0)$, it yields that
\[
\eta(R) = \sup_{x,y \in B_R(x_0)} |r(x,t) - r(y,t)| \leq f^+(R) - f^-(R).
\]

Next, we show that $$f^+(R) - f^-(R) \leq \eta(R).$$ 
Let $\varepsilon > 0$ be arbitrary. By the definitions of supremum and infimum, there exist points $x_\varepsilon, y_\varepsilon \in B_R(x_0)$ such that
\[
r(x_\varepsilon, t) > f^+(R) - \frac{\varepsilon}{2} \quad \text{and} \quad r(y_\varepsilon, t) < f^-(R) + \frac{\varepsilon}{2}.
\]
Then,
\[
r(x_\varepsilon, t) - r(y_\varepsilon, t) > \left(f^+(R) - \frac{\varepsilon}{2}\right) - \left(f^-(R) + \frac{\varepsilon}{2}\right) = f^+(R) - f^-(R) - \varepsilon.
\]
Since $|r(x_\varepsilon, t) - r(y_\varepsilon, t)| \geq r(x_\varepsilon, t) - r(y_\varepsilon, t)$, we have
\[
|r(x_\varepsilon, t) - r(y_\varepsilon, t)| > f^+(R) - f^-(R) - \varepsilon.
\]
But by the definition of $\eta(R)$,
\[
\eta(R) \geq |r(x_\varepsilon, t) - r(y_\varepsilon, t)| > f^+(R) - f^-(R) - \varepsilon.
\]
Since $\varepsilon > 0$ is arbitrary, it follows that
\[
\eta(R) \geq f^+(R) - f^-(R).
\]

Combining the two inequalities, we obtain
\[
\eta(R) = f^+(R) - f^-(R).
\]
By definition, $\operatorname{osc}_{B_R} r(x,t) = f^+(R) - f^-(R)$, so
\[
\operatorname{osc}_{B_R} r(x,t) = f^+(R) - f^-(R) = \eta(R),
\]
as desired.


\textbf{Proof of (2):}
 From the definitions of the infimum and supremum, we immediately obtain  
\[
f^-(R) \leq r(x,t) \leq f^+(R), \qquad \forall\, x \in B_R(x_0).
\]  
Since $B_R(x_0)\subset \Omega$, this establishes the desired chain of inequalities.

\textbf{Proof of (3):}
Fix $x \in B_R(x_0)$ and $t > 1$. Consider the function $h(s) = \ln G(x, s)$. Its derivative is
\[
h'(s) = \frac{g(x, s)}{G(x, s)} = \frac{r(x, s)}{s}.
\]
From part (2), we have $f^-(R) \leq r(x, s) \leq f^+(R)$ for all $s > 1$. This implies the differential inequalities
\[
\frac{f^-(R)}{s} \leq h'(s) \leq \frac{f^+(R)}{s}.
\]
Integrating these inequalities from $s=1$ to $s=t$, it yields ttha
\begin{align*}
f^-(R) \ln t &\leq \ln G(x, t) - \ln G(x, 1) \leq f^+(R) \ln t.
\end{align*}
Exponentiating and using the bounds $F^{-1} \leq G(x, 1) \leq F$ from \eqref{GG1}, gives the desired result
\[
F^{-1} t^{f^-(R)} \leq G(x, t) \leq F t^{f^+(R)}.
\]

\textbf{Proof of (4):}
From \eqref{GG3} and $(1)$, we have
\[
\eta(R) \leq \frac{L_0}{|\ln(2R)|},
\]
hence, for $R < \frac{1}{2}$, since $|\ln(2R)| > |\ln R|$, one has 
\begin{equation}\label{eq1eq1eq1}
\eta(R) < \frac{L_0 }{ |\ln R|}. 
\end{equation}
Consequently,
\begin{align*}
R^{f^-(R) - f^+(R)} &= \exp\left( -\eta(R) \ln R \right) = \exp\left( \eta(R) |\ln R| \right) \\
&\leq \exp\left( \frac{L_0}{|\ln R|} \cdot |\ln R| \right) = \exp(L_0) = L.
\end{align*}

\textbf{Proof of (5):}
From \eqref{GG3}, it follows directly that
\[
0 \leq \lim_{R \to 0^+} \eta(R) \leq \lim_{R \to 0^+} \frac{L_0}{|\ln(2R)|} = 0,
\]
which completes the proof.
\end{proof}
\begin{lemma} \label{remff}
Assume \eqref{D22}--\eqref{GG1} and \eqref{GG3} hold. For any $M \geq 1$, there exists $R_M > 0$ such that for all $R < R_M$,
\[
M^{f^+(R) - f^-(R)} \leq 2.
\]
Specifically, one may take $R_M = \min\left\{ R_0,\, \exp\left( -\frac{L_0 \ln M}{\ln 2} \right) \right\}$, where $R_0$ and $L_0$ are the constants from \eqref{GG3}.
\end{lemma}

\begin{proof}
Recall, from Lemma \ref{lem22}, that $\eta(R) = f^+(R) - f^-(R)$. By \eqref{eq1eq1eq1}, we have $\eta(R) \leq \frac{L_0}{|\ln R|}$ for $R < R_0$. Consequently,
\[
M^{\eta(R)} \leq \exp\left( \frac{L_0 \ln M}{|\ln R|} \right).
\]
Choose $R_M = \min\left\{ R_0,\, \exp\left( -\frac{L_0 \ln M}{\ln 2} \right) \right\}$. Then, for all $R < R_M$,
\[
\frac{L_0 \ln M}{|\ln R|} \leq \ln 2,
\]
which implies
\[
M^{\eta(R)} \leq \exp(\ln 2) = 2.
\]
\end{proof}
\begin{remark}
In view of assumption~\eqref{GG3}, together with Lemma~\ref{Lemma 2.3} and Lemma~\ref{lem22}-(1), the regularity of weak solutions to Problem~\eqref{P} is strongly influenced by the behavior of the function $r(x,t)$, in particular by its oscillation and regularity properties. This highlights the central role of assumption~\eqref{GG3}.
\end{remark}
\begin{proof}[\textbf{Proof of Theorem \ref{Theorem 2.1}}]
Let $u \in \mathcal{B}_{G(x,t)}(\Omega, M, \gamma, \gamma_1, \delta)$, where $G(x,t)$ satisfies conditions \eqref{D22}--\eqref{GG1} and \eqref{GG3}. Without loss of generality, assume $M \geq 1$, $\gamma \geq 1$ and $L\geq 1$.

To prove Theorem~\ref{Theorem 2.1}, it suffices to show that for every point $x_0 \in \Omega$, there exists a ball $B_{R_0}(x_0) \subset \Omega$ such that $u \in C^{0,\alpha}(B_{R_0}(x_0))$, where $\alpha$ is the constant stated in the theorem.

Fix $x_0 \in \Omega$ and choose $R_0 \in (0, R_M)$ where $R_M$ is from Lemma \ref{remff}, such that $\overline{B}_{R_0}(x_0) \subset \Omega$ and $R_M < 1$. Then, for an arbitrary $R \in (0, R_0]$, at least one of the functions $w := \pm u$ satisfies the following condition.
\begin{equation}\label{2.8}
\left\vert\left\lbrace x \in B_{\frac{R}{2}} : w(x) > \max_{B_R} w(x) - \tfrac{1}{2} \mathrm{osc}_{B_R} u \right\rbrace\right\vert \leq \tfrac{1}{2} \, \vert B_{\frac{R}{2}}\vert.
\end{equation}
From this point on, we denote by $w$ the function (either $u$ or $-u$) that satisfies inequality \eqref{2.8}. Define
\begin{equation}\label{2.9}
\tau := \max\left\{2, \frac{2}{\delta} \right\}, \quad \psi := \tau^{-1} \, \mathrm{osc}_{B_R} u,
\end{equation}
and
\begin{equation}\label{2.10}
k' := \max_{B_R} w(x) - \psi.
\end{equation}
It follows that
\begin{equation}\label{2.11}
k' \geq \max_{B_R} w(x) - \tfrac{1}{2} \, \mathrm{osc}_{B_R} u,
\end{equation}
and
\begin{equation}\label{2.12}
k' \geq \max_{B_R} w(x) - \delta M.
\end{equation}
\begin{remark}
Note that inequality~\eqref{2.12} guarantees that conditions~\eqref{2.4} and~\eqref{2.5} hold for all \( k \geq k' \) and for all \( \rho \leq R \).
\end{remark}
Under these assumptions and notations, the proof of Theorem~\ref{Theorem 2.1} is completed through the lemmas presented below, specifically Lemmas~\ref{Lemma 2.5} to~\ref{Lemma 2.8}.
\begin{lemma}\label{Lemma 2.5}
There is a positive constant \(\theta = \theta(n,g^+, L, F,\gamma)\) such that if
\begin{equation}\label{2.13}
|A_{k^0, \frac{R}{2}}| \leq \theta R^n,
\end{equation}
then at least one of the following two inequalities hol{\rm d}s:
\begin{equation}\label{2.14}
H \leq \left( \frac{\gamma  + \gamma_1 + 1}{\gamma} \right)^{\frac{1}{f^+(R)}} R,
\end{equation}
or
\begin{equation}\label{2.15}
\max_{B_{\frac{R}{4}}} w(x) \leq k^0 + \frac{H}{2},
\end{equation}
where
\[
0 < H < \psi, \quad k^0 = \max_{B_R} w(x) - H.
\]
\end{lemma}

\begin{proof}
If inequality \eqref{2.14} does not hold, i.e., if we assume
\begin{equation}\label{2.144}
\left( \frac{\gamma + \gamma_1 + 1}{\gamma} \right) <  \left(\frac{H}{R}\right)^{f^+(R)}.
\end{equation}
For \(h = 0, 1, 2, \ldots\), set
\[
\rho_h = \frac{R}{4} + \frac{R}{2^{h+2}}, \quad k_h = k^0 + \frac{H}{2} - \frac{H}{2^{h+1}},
\]
\[
Y_h = R^{-n}|A_{k_h, \rho_h}|, \quad D_{h+1} = A_{k_h, \rho_{h+1}} \backslash A_{k_{h+1}, \rho_{h+1}}.
\]
{\bf Claim :} \(Y_h \to 0\) as \(h \to \infty\).\\
Applying inequality \eqref{2.4} to \(k = k_h\), \(\rho = \rho_h\) and \(\rho - \sigma \rho = \rho_{h+1}\) for \(h = 0, 1, 2, \ldots\), and using Proposition~\ref{zoo} and Lemma~\ref{lem22} (with $t=\frac{H}{R}>1$), we obtain
\begin{align}\label{2.17}
\begin{split}
 \int_{A_{k_h, \rho_{h+1}}} G(x,|Dw|){\rm d}x {\rm d}x
&\leq \gamma \int_{A_{k_h, \rho_h}}G\left(x,\left| \frac{w(x) - k_h}{\rho_h - \rho_{h+1}} \right|\right) {\rm d}x + \gamma_1 |A_{k_h, \rho_h}|\\
& \leq \gamma \int_{A_{k_h, \rho_h}}G\left(x,\left| \frac{\max_{B_R} w(x) - \left( k^0 + \frac{H}{2} - \frac{H}{2^{h+1}}\right)}{\frac{R}{2^{h+2}}-\frac{R}{2^{h+3}}} \right|\right) {\rm d}x + \gamma_1 |A_{k_h, \rho_h}|\\
& \leq \gamma \int_{A_{k_h, \rho_h}}G\left(x, \frac{ \frac{H}{2} + \frac{H}{2^{h+1}}}{\frac{R}{2^{h+3}}} \right) {\rm d}x + \gamma_1 |A_{k_h, \rho_h}|\\
& \leq \gamma \int_{A_{k_h, \rho_h}}G\left(x, \frac{ \frac{(2^{h}+1)H}{2^{h+1}}}{\frac{R}{2^{h+3}}} \right) {\rm d}x + \gamma_1 |A_{k_h, \rho_h}|\\
& \leq \gamma \int_{A_{k_h, \rho_h}}G\left(x,  \frac{2^{h+3}H}{R} \right) {\rm d}x + \gamma_1 |A_{k_h, \rho_h}|\\
& \leq \gamma \left(2^{h+3}\right)^{g^+} \int_{A_{k_h, \rho_h}}G\left(x,  \frac{H}{R} \right) {\rm d}x + \gamma_1 |A_{k_h, \rho_h}|\\
& \leq \gamma \left[\left(2^{h+3}\right)^{g^+}\int_{A_{k_h, \rho_h}}F\left(\frac{H}{R} \right)^{f^+(R)}   {\rm d}x  + \frac{\gamma_1}{\gamma} |A_{k_h, \rho_h}|\right]\\
& \leq \gamma F\left[\left(2^{h+3}\right)^{g^+} H^{f^+(R)}R^{-f^+(R)}|A_{k_h, \rho_h}| + \frac{\gamma_1}{\gamma} |A_{k_h, \rho_h}|\right].
\end{split}
\end{align}
In light of Lemma~\ref{lem22}, we see that
\begin{align}\label{Ge2h}
\begin{split}
    \int_{A_{k_h, \rho_{h+1}}} G(x,|Dw|){\rm d}x &\geq F^{-1}\left[\int_{A_{k_h, \rho_{h+1}}} |Dw|^{f^-(R)}{\rm d}x - |A_{k_h, \rho_{h+1}}|\right] \\
    & \geq F^{-1}\left[\int_{D_{h+1}} |Dw|^{f^-(R)}{\rm d}x - |A_{k_h, \rho_{h}}|\right].
       \end{split}
\end{align}
Therefore, by inequality \eqref{2.17}, it follows that
\begin{align}\label{2.18}
\begin{split}
\int_{D_{h+1}} |Dw|^{f^-(R)}{\rm d}x &\leq \gamma F^2\left[\left(2^{h+3}\right)^{g^+}H^{f^+(R)}R^{-f^+(R)}|A_{k_h, \rho_h}|   + \frac{\gamma_1}{\gamma} |A_{k_h, \rho_h}|\right]+|A_{k_h, \rho_h}|\\
&\leq \gamma F^2\left[\left(2^{h+3}\right)^{g^+}H^{f^+(R)}R^{-f^+(R)}|A_{k_h, \rho_h}|   + \frac{\gamma+\gamma_1+1}{\gamma} |A_{k_h, \rho_h}|\right]\\
&\leq \gamma F^2\left[\left(2^{h+3}\right)^{g^+}H^{f^+(R)}R^{-f^+(R)}   + \frac{\gamma+\gamma_1+1}{\gamma} \right]|A_{k_h, \rho_h}|\\
&\leq \gamma F^2\left[\left(2^{h+3}\right)^{g^+}H^{f^+(R)}R^{-f^+(R)}   + \frac{\gamma+\gamma_1+1}{\gamma} \right]R^nY_h.
\end{split}
\end{align}
Hence, by inequalities \eqref{2.144} and \eqref{2.18}, we have 
\begin{align}\label{Ge6}
\begin{split}
\int_{D_{h+1}} |Dw|^{f^-(R)}{\rm d}x 
&\leq  \gamma F^2\left[\left(2^{h+3}\right)^{g^+}H^{f^+(R)}R^{-f^+(R)}   + \frac{\gamma+\gamma_1+1}{\gamma} \right]R^nY_h\\
& \leq \gamma F^2\left[\left(2^{h+3}\right)^{g^+}H^{f^+(R)}R^{-f^+(R)}   + \left(HR^{-1}\right)^{f^+(R)} \right]R^nY_h\\
&\leq \gamma F^2\left(2^{h+4}\right)^{g^+} H^{f^+(R)} R^{n-f^+(R)}Y_h.
\end{split}
\end{align}
Consequently, using Hölder's inequality and inequality \eqref{Ge6}, we obtain
\begin{align}\label{2.20}
\begin{split}
\int_{D_{h+1}} |Dw| {\rm d}x &\leq \left( \int_{D_{h+1}} |Dw|^{f^-(R)} {\rm d}x \right)^{\frac{1}{f^-(R)}} |D_{h+1}|^{\frac{f^-(R)-1}{f^-(R)}}\\
&\leq \left( \gamma F^2\left(2^{h+4}\right)^{g^+} H^{f^+(R)} R^{n-f^+(R)}Y_h\right)^{\frac{1}{f^-(R)}} (R^n Y_h)^{\frac{f^-(R)-1}{f^-(R)}}\\
&= \left( \gamma F^2\left(2^{h+4}\right)^{g^+} H^{f^+(R)} \right)^{\frac{1}{f^-(R)}}R^{\frac{nf^-(R)-f^+(R)}{f^-(R)}}Y_h.
\end{split}
\end{align}
Applying inequality \eqref{2.2} to \(k = k_h\), \(l = k_{h+1}\) and \(\rho = \rho_{h+1}\), we get
\begin{align}\label{2.21}
\begin{split}
\int_{D_{h+1}} |Dw| {\rm d}x &\geq \beta^{-1} (k_{h+1} - k_h)(R^n Y_{h+1})^{1-\frac{1}{n}} \rho_{h+1}^{-n} |B_{\rho_{h+1}} \backslash A_{k_h, \rho_{h+1}}|\\
&\geq \beta^{-1} 2^{-(h+2)} HR^{n-1} Y_{h+1}^{1-\frac{1}{n}} 2^n R^{-n} \left(|B_{\frac{R}{4}}| - |A_{k^0, \frac{R}{2}}|\right).
\end{split}
\end{align}
If inequality \eqref{2.13} hol{\rm d}s and
\[
\theta \leq \frac{1}{2} \cdot 4^{-n} \omega_n,
\]
then
\[
|B_{R/4}| - |A_{k^0, \frac{R}{2}}| \geq 4^{-n}R^n\omega_n - \frac{1}{2} \cdot 4^{-n}\omega_nR^n = 2^{-2n-1}\omega_nR^n
\]
and consequently from inequality \eqref{2.21}, it follows that
\begin{equation}\label{eqeq1}
\int_{D_{h+1}} |Dw| {\rm d}x \geq \beta^{-1}2^{-(h+3+n)}HR^{n-1}\omega_n Y_{h+1}^{1-\frac{1}{n}}.
\end{equation}
From inequalities \eqref{2.20}, \eqref{eqeq1} and the fact that
\[0<H<\psi\leq M,\ M\geq 1,\ \gamma\geq 1,\ L\geq 1,\ R<1,\] 
we deduce, by Lemma \ref{remff}, that
\begin{align}
\begin{split}
Y_{h+1}^{1-\frac{1}{n}} &\leq \beta 2^{h+3+n}H^{-1}R^{1-n}\omega_n^{-1} \int_{D_{h+1}} |Dw| {\rm d}x\\
&\leq \beta 2^{h+3+n}H^{-1}R^{1-n}\omega_n^{-1}\left( \gamma F^2\left(2^{h+4}\right)^{g^+} H^{f^+(R)} \right)^{\frac{1}{f^-(R)}}R^{\frac{nf^-(R)-f^+(R)}{f^-(R)}}Y_h\\
&= \beta 2^{h+3+n}\omega_n^{-1}H^{\frac{f^+(R)-f^-(R)}{f^{-}(R)}}\left( \gamma F^2\left(2^{h+4}\right)^{g^+}  \right)^{\frac{1}{f^-(R)}}R^{\frac{nf^-(R)-f^+(R)}{f^-(R)}-n+1}Y_h\\
& \leq \beta 2^{h+3+n}\omega_n^{-1}M^{\frac{f^+(R)-f^-(R)}{f^{-}(R)}}\left( \gamma F^2\left(2^{h+4}\right)^{g^+}  \right)^{\frac{1}{f^-(R)}} R^{\frac{f^-(R)-f^+(R)}{f^-(R)}}Y_h\\
& \leq \gamma F^2\left(2^{h+4}\right)^{g^+}\beta 2^{h+3+n}\omega_n^{-1}2^{\frac{1}{f^{-}(R)}} R^{f^-(R)-f^+(R)}Y_h\\
& \leq \gamma F^2\left(2^{h+4}\right)^{g^+}\beta 2^{h+4+n}\omega_n^{-1} R^{f^-(R)-f^+(R)}Y_h\\
& \leq 2^{(g^++1)h}\gamma F^2\beta 2^{4g^++4+n}\omega_n^{-1} R^{f^-(R)-f^+(R)}Y_h\\
& \leq 2^{(g^++1)h}\gamma F^2\beta 2^{4g^++4+n}\omega_n^{-1} L Y_h\\
&\leq c_1 b_1^h Y_h,
\end{split}
\end{align}
where \(b_1 = 2^{(g^++1)} > 1\), \(c_1 = \gamma F^2\beta 2^{4g^++4+n}\omega_n^{-1} L+1 > 1\) and \(\beta = \beta(n)\) as in Lemma \ref{Lemma 2.2}. Hence
\[
Y_{h+1} \leq c b^h Y_h^{1+\frac{1}{n-1}} = c b^h Y_h^{1+\varepsilon},
\]
where \(b = b_1^{\frac{n}{n-1}} = b(n,g^+) > 1\), \(c = c_1^{\frac{n}{n-1}} = c(n,g^+, L, F,\gamma)\) and \(\varepsilon = \frac{1}{n-1}\).
Now we choose
\[
\theta = \min\left\{ \frac{1}{2} \cdot 4^{-n}\omega_n, \quad c^{-\frac{1}{\varepsilon}}b^{-\frac{1}{\varepsilon^2}} \right\}.
\]
Then \(\theta = \theta(n,g^+, L, F,\gamma)\) and when
\[
|A_{k^0, \frac{R}{2}}| \leq \theta R^n \quad \text{and} \quad H > \left( \frac{\gamma  + \gamma_1 + 1}{\gamma} \right)^{\frac{1}{f^+(R)}}R,
\]
we have
\[
Y_0 = R^{-n}|A_{k^0, \frac{R}{2}}| \leq \theta \leq c^{-\frac{1}{\varepsilon}}b^{-\frac{1}{\varepsilon^2}}
\]
and consequently, by Lemma \ref{Lemma 2.1}, \(Y_h \to 0\) as \(h \to \infty\). This shows the Claim.\\
From the claim, it yields that
\[
\max_{B_{\frac{R}{4}}} w(x) \leq \lim_{h \to \infty} k_h = k^0 + \frac{H}{2}.
\]
Lemma \ref{Lemma 2.5} is proved.
\end{proof}

\begin{lemma}\label{Lemma 2.6}
For any given \(\theta > 0\), there is a natural number \(s = s(n,g^-,g^+,F,L,\beta,\gamma) > 2\) such that either
\begin{equation}\label{2.29}
\psi \leq 2^s \left( \frac{\gamma + \gamma_1 + 1}{\gamma} \right)^{\frac{1}{f^+(R)}} R
\end{equation}
or
\begin{equation}\label{2.30}
|A_{k^0, \frac{R}{2}}| \leq \theta R^n
\end{equation}
hol{\rm d}s, where
\[
k^0 = \max_{B_R} w(x) - 2^{-s+1}\psi.
\]
\end{lemma}

\begin{proof}
If inequality \eqref{2.29} does not hold, i.e., if we assume
\begin{equation}\label{2.32}
\psi > 2^s \left( \frac{\gamma + \gamma_1 + 1}{\gamma} \right)^{\frac{1}{f^+(R)}} R,
\end{equation}
where the natural number \( s \) will be chosen later, we proceed as follows.\\
For \( t = 0, \ldots, s - 1 \), define
\[
k_t = \max_{B_R} w(x) - 2^{-t} \psi, \quad D_t = A_{k_t, \frac{R}{2}} \setminus A_{k_{t+1}, \frac{R}{2}}.
\]
By \eqref{2.32}, one has  
\begin{equation}\label{2.32m}
\frac{ 2^{-t+1}\psi}{R}> 1,\quad \text{for } t = 0, \ldots, s - 1.
\end{equation}
Applying inequality \eqref{2.4} to \(\rho = R\), \(\rho - \sigma \rho = \frac{R}{2}\) and \(k = k_t\) for \(t = 0, \ldots, s - 2\), and using Proposition~\ref{zoo}, Lemma ~\ref{lem22} (with $t= \frac{2^{-t+1}\psi}{R} $), we  obtain
\begin{align}\label{Ge1}
\begin{split}
\int_{A_{k_t, \frac{R}{2}}} G(x,|Dw|){\rm d}x &\leq \gamma\left[\int_{A_{k_t, R}}G\left(x,\frac{\omega(x)-k_t}{\frac{R}{2}}\right){\rm d}x +\frac{\gamma_1}{\gamma} |A_{k_t, R}|\right]\\
& \leq\gamma\left[\int_{A_{k_t, R}}G\left(x,\frac{2^{-t}\psi}{\frac{R}{2}}\right){\rm d}x +\frac{\gamma_1}{\gamma} |A_{k_t, R}|\right]\\
&\leq \gamma\left[\int_{A_{k_t, R}}G\left(x,\frac{2^{-t+1}\psi}{R}\right){\rm d}x+\frac{\gamma_1}{\gamma} |A_{k_t, R}|\right]\\
& \leq \gamma \left[\int_{A_{k_t, R}}F \left(\frac{2^{-t+1}\psi}{R}\right)^{f^+(R)}{\rm d}x+\frac{\gamma_1}{\gamma} |A_{k_t, R}|\right]\\
& \leq \gamma F\left[\left(2^{-t+1}\psi\right)^{f^+(R)}R^{-f^+(R)}|A_{k_t, R}|+\frac{\gamma_1}{\gamma} |A_{k_t, R}|\right].
\end{split}
\end{align}
On the other hand, from Lemma~\ref{lem22}, we deduce that
\begin{align}\label{Ge2}
\begin{split}
    \int_{A_{k_t, \frac{R}{2}}} G(x,|Dw|){\rm d}x &\geq F^{-1}\left[\int_{A_{k_t, \frac{R}{2}}} |Dw|^{f^-(R)}{\rm d}x - |A_{k_t, \frac{R}{2}}|\right] \\
    &\geq F^{-1}\left[\int_{D_{t}} |Dw|^{f^-(R)}{\rm d}x - |A_{k_t, R}|\right].
   \end{split}
\end{align}
Thus, combining  \eqref{2.32}, \eqref{Ge1}, and \eqref{Ge2}, we get
\begin{align}\label{Ge3}
\begin{split}
  \int_{D_t} |Dw|^{f^-(R)}{\rm d}x & \leq \gamma F^2\left[\left(2^{-t+1}\psi\right)^{f^+(R)}R^{-f^+(R)}|A_{k_t, R}|+\frac{\gamma_1}{\gamma} |A_{k_t, R}|\right]+|A_{k_t, R}|\\
  & \leq \gamma F^2\left[\left(2^{-t+1}\psi\right)^{f^+(R)}R^{-f^+(R)}|A_{k_t, R}|+\frac{\gamma+\gamma_1+1}{\gamma} |A_{k_t, R}|\right]\\
  &  \leq\gamma F^2\left[\left(2^{-t+1}\psi\right)^{f^+(R)}R^{-f^+(R)}+\left( 2^{-s}\psi R^{-1}\right)^{f^+(R)} \right]|A_{k_t, R}|\\
  &\leq  \gamma F^2\left[\left(2^{-t+1}\psi\right)^{f^+(R)}R^{-f^+(R)}+\left( 2^{-s}\psi R^{-1}\right)^{f^+(R)} \right]\omega_n R^n\\
  &\leq  \gamma F^2\left[\left(2^{-t+1}\psi\right)^{f^+(R)}+\left(2^{-t}\psi\right)^{f^+(R)} \right]\omega_n R^{n-f^+(R)}\\
  &\leq  \gamma F^22^{f^+(R)+1}\omega_n\left(2^{-t}\psi\right)^{f^+(R)} R^{n-f^+(R)}\\
  &\leq  \gamma F^22^{g^++1}\omega_n\left(2^{-t}\psi\right)^{f^+(R)} R^{n-f^+(R)},\quad (t = 0, \ldots, s - 2).
  \end{split}
\end{align}
Now, applying inequality \eqref{2.2} to \(w\) with \(k = k_t\), \(l = k_{t+1}\), \(\rho = \frac{R}{2}\), for \(t = 0, \ldots, s - 2\), we obtain
\begin{align}\label{2.34}
\begin{split}
|A_{k_{s-1}, \frac{R}{2}}|^{1-\frac{1}{n}} &\leq |A_{k_{t+1}, \frac{R}{2}}|^{1-\frac{1}{n}}\\
&\leq \beta 2^{-n} R^n (k_{t+1} - k_t)^{-1} |B_{\frac{R}{2}} \backslash A_{k_t, \frac{R}{2}}|^{-1} \int_{D_t} |Dw| {\rm d}x\\
&\leq \beta 2^{-n} R^n (2^{-(t+1)}\psi)^{-1} (2^{-1} \omega_n (2^{-1} R)^n)^{-1} \left( \int_{D_t} |Dw|^{f^-(R)} {\rm d}x \right)^{\frac{1}{f^-(R)}} |D_t|^{\frac{f^-(R)-1}{f^-(R)}}\\
&\leq 4 \beta \omega_n^{-1} 2^t \psi^{-1} \left( \int_{D_t} |Dw|^{f^-(R)} {\rm d}x \right)^{\frac{1}{f^-(R)}} |D_t|^{\frac{f^-(R)-1}{f^-(R)}}.
\end{split}
\end{align}
It follows, from inequalities \eqref{Ge3}, \eqref{2.34} and Lemma \ref{remff}, that
\begin{align}\label{2.35}
\begin{split}
|A_{k_{s-1}, \frac{R}{2}}|^{\frac{(n-1)f^-(R)}{n}} &\leq (4\beta\omega_n^{-1}2^t\psi^{-1})^{f^-(R)} \int_{D_t} |Dw|^{f^-(R)} {\rm d}x \cdot |D_t|^{f^-(R) - 1}\\
&\leq (4\beta\omega_n^{-1}2^t\psi^{-1})^{f^-(R)}\gamma F^22^{g^++1}\omega_n\left(2^{-t}\psi\right)^{f^+(R)} R^{n-f^+(R)}|D_t|^{f^-(R) - 1}\\
&\leq (2^n\beta\omega_n^{-1})^{g^+}\gamma F^22^{g^++1}\omega_n\left(2^{-t}\psi\right)^{f^+(R)-f^-(R)} R^{n-f^+(R)}|D_t|^{f^-(R) - 1}\\
&\leq (2^n\beta\omega_n^{-1})^{g^+}\gamma F^22^{g^++1}\omega_n M^{f^+(R)-f^-(R)} R^{n-f^+(R)}|D_t|^{f^-(R) - 1}\ \ (\text{since}\ 2^{-t}\psi \leq \psi \leq M )\\
&\leq (2^n\beta\omega_n^{-1})^{g^+}\gamma F^22^{g^++2}\omega_n  R^{n-f^+(R)}|D_t|^{f^-(R) - 1},\quad (t = 0, \ldots, s - 2).
\end{split}
\end{align}
Consequently, by inequality \eqref{2.35}, we deduce that
\begin{equation}\label{2.36}
|A_{k_{s-1}, \frac{R}{2}}|^{\frac{(n-1)f^-(R)}{n}} \leq c_1 R^{n - f^+(R)} |D_t|^{f^-(R) - 1}, \quad (t = 0, \ldots, s - 2),
\end{equation}
where \(c_1 = (2^n\beta\omega_n^{-1})^{g^+}\gamma F^22^{g^++2}\omega_n + 1 = c(n,g^+,F,\beta,\gamma) > 1\).\\
Thus, due to inequality \eqref{2.36}, we infer that
\begin{equation}\label{2.37}
|A_{k_{s-1}, \frac{R}{2}}|^{\frac{(n-1)f^-(R)}{n(f^-(R) - 1)}} \leq c_2R^{\frac{n-f^+(R)}{(f^-(R) - 1)}} |D_t|, \quad (t = 0, \ldots, s - 2),
\end{equation}
where \(c_2:=  c_1^{\frac{1}{g^--1}}=c(n,g^-,g^+,F,\beta,\gamma)\).\\
Summing inequality \eqref{2.37} with respect to \(t\) from $0$ to \(s - 2\) and noting that
\[
\sum_{t=0}^{s-2} |D_t|\leq |B_{\frac{R}{2}}| = \omega_n \left(\frac{R}{2}\right)^n,
\]
we obtain
\[
|A_{k_{s-1}, \frac{R}{2}}|^{\frac{(n-1)f^-(R)}{n(f^-(R) - 1)}} \leq \frac{c_22^{-n}\omega_n}{s  - 1} R^{\frac{n-f^+(R)}{(f^-(R) - 1)}+n}=\frac{c_22^{-n}\omega_n}{s - 1} R^{\frac{nf^-(R)-f^+(R)}{(f^-(R) - 1)}},
\]
and therefore, using Lemma \ref{lem22}, we get
\begin{align}\label{2.38}
\begin{split}
|A_{k_{s-1}, \frac{R}{2}}| &\leq \left(\frac{c_22^{-n}\omega_n}{s - 1}\right)^{\frac{n(f^-(R) - 1)}{(n-1)f^-(R)}} R^{\frac{nf^-(R)-f^+(R)}{(f^-(R) - 1)}\times\frac{n(f^-(R) - 1)}{(n-1)f^-(R)}}\\
&\leq \left(\frac{c_22^{-n}\omega_n}{s - 1}\right)^{\frac{n(f^-(R) - 1)}{(n-1)f^-(R)}} R^{\frac{nf^-(R)-f^+(R)}{f^-(R)}\times\frac{n}{(n-1)}}\\
& \leq  \left(\frac{c_22^{-n}\omega_n}{s - 1}\right)^{\frac{n(f^-(R) - 1)}{(n-1)f^-(R)}} R^{\frac{n\left(f^-(R)-f^+(R)\right)}{f^-(R)}\times\frac{n}{(n-1)}} R^n\\
& \leq  \left(\frac{c_22^{-n}\omega_n}{s - 1}\right)^{\frac{n(f^-(R) - 1)}{(n-1)f^-(R)}} L^{\frac{n}{f^-(R)}\times\frac{n}{(n-1)}} R^n\\
& \leq  \left(\frac{c_22^{-n}\omega_n}{s - 1}\right)^{\frac{n(f^-(R) - 1)}{(n-1)f^-(R)}} L^{\frac{n}{g^-}\times\frac{n}{(n-1)}} R^n.
\end{split}
\end{align}
Now we choose a natural number \(s\) such that \(\frac{c_22^{-n}\omega_n}{s - 1}<1\) and
\begin{equation}\label{2.39}
\left(\frac{c_22^{-n}\omega_n}{s - 1}\right)^{\frac{n(g^- - 1)}{(n-1)g^-}} L^{\frac{n}{g^-}\times\frac{n}{(n-1)}}\leq \theta.
\end{equation}
Note that \(s = s(\theta, c_2, n, g^-,L) = s(\theta,n,g^-,g^+,F,L,\beta,\gamma)\). Then for such \(s\), when inequality \eqref{2.29} does not hold, from inequalities \eqref{2.38} and \eqref{2.39}, we have
\[
|A_{k_{s-1}, \frac{R}{2}}| \leq \theta R^n
\]
which means \(|A_{k^0, \frac{R}{2}}| \leq \theta R^n\), i.e., inequality \eqref{2.30} hol{\rm d}s, since
\[
k^0 = \max_{B_R} w(x) - 2^{-s+1}\psi = k_{s-1}.
\]
Lemma \ref{Lemma 2.6} is proved.
\end{proof}
\begin{lemma}\label{Lemma 2.7}
There is a number \(s = s(n, g^-, g^+,  L,F, \gamma) > 2\) such that
\begin{equation}\label{2.40}
\psi \leq 2^s \max \left\{ \max_{B_R} w(x) - \max_{B_{\frac{R}{4}}} w(x), \left( \frac{\gamma  + \gamma_1 + 1}{\gamma} \right)^{\frac{1}{f^+(R)}} R \right\},\text{    for any } R \in (0, R_0]. 
\end{equation}
\end{lemma}
\begin{proof}
Let \( R \in (0, R_0] \), and let 
\[
\theta = \theta(n, g^-, g^+, F,L, \gamma)
\]
be the constant appearing in Lemma~\ref{Lemma 2.5}. Applying Lemma~\ref{Lemma 2.6} with this choice of \( \theta \), we obtain a constant 
\[
s = s(\theta, n, g^-, g^+, F,L, \gamma) = s(n, g^-, g^+, F,L, \gamma)
\]
such that at least one of the inequalities \eqref{2.29} or \eqref{2.30} holds.
If \eqref{2.29} holds, then \eqref{2.40} follows immediately. Now suppose that \eqref{2.29} does not hold. Then, by Lemma~\ref{Lemma 2.6}, inequality \eqref{2.30} must be satisfied; that is,
\[
|A_{k^0, \frac{R}{2}}| \leq \theta R^n,
\]
where \( k^0 = \max_{B_R} w(x) - 2^{-s+1} \psi \). Set \( H > 2^{-s+1} \psi \).
 Then, since inequality \eqref{2.29} does not hold, we obtain
\[
H > 2^{-s+1} \cdot 2^s \left( \frac{\gamma  + \gamma_1 + 1}{\gamma} \right)^{\frac{1}{f^+(R)}} R > \left( \frac{\gamma + \gamma_1 + 1}{\gamma} \right)^{\frac{1}{f^+(R)}} R.
\]
Thus, by Lemma \ref{Lemma 2.5}, one has
\[
\max_{B_{\frac{R}{4}}} w(x) \leq k^0 + \frac{H}{2} = \max_{B_R} w(x) - 2^{-s}\psi
\]
i.e.,
\[
\psi \leq 2^s \left( \max_{B_R} w(x) - \max_{B_{\frac{R}{4}}} w(x) \right).
\]
This shows that inequality \eqref{2.40} hol{\rm d}s. Lemma \ref{Lemma 2.7} is proved.
\end{proof}
\begin{lemma}\label{Lemma 2.8}
For any \(R \in (0, R_0]\), at least one of the following two inequalities hol{\rm d}s:
\begin{equation}\label{2.41}
\operatorname{osc}_{B_R}u \leq \tau 2^s \frac{\gamma + \gamma_1 + 1}{\gamma} R,
\end{equation}
or
\begin{equation}\label{2.42}
\operatorname{osc}_{B_{\frac{R}{4}}}u \leq (1 - \tau^{-1} 2^{-s}) \operatorname{osc}_{B_R}u,
\end{equation}
where \(\tau = \max \left\{2, \frac{2}{\delta}\right\}\), \(s = s(n, g^-, g^+, F,L, \gamma)\) is the constant as in Lemma \ref{Lemma 2.7}.
\end{lemma}
\begin{proof}
By Lemma \ref{Lemma 2.7}, inequality \eqref{2.40} hol{\rm d}s, this implies that at least one of the following two inequalities hol{\rm d}s:
\begin{equation}\label{2.43}
\psi \leq 2^s \left( \frac{\gamma  + \gamma_1 + 1}{\gamma} \right)^{\frac{1}{f^+(R)}} R,
\end{equation}
or
\begin{equation}\label{2.44}
\psi \leq 2^s \left( \max_{B_R} w(x) - \max_{B_{\frac{R}{4}}} w(x) \right).
\end{equation}
When inequality \eqref{2.43} hol{\rm d}s, we have
\[
\operatorname{osc}_{B_R}u = \tau \psi \leq \tau 2^s \left( \frac{\gamma  + \gamma_1 + 1}{\gamma} \right)^{\frac{1}{f^+(R)}} R \leq \tau 2^s \frac{\gamma + \gamma_1 + 1}{\gamma} R,
\]
i.e., inequality \eqref{2.41} hol{\rm d}s. When inequality \eqref{2.44} hol{\rm d}s, we have
\begin{align*}
\operatorname{osc}_{B_R}u, &= \tau \psi \leq \tau 2^s \left( \max_{B_R} w(x) - \max_{B_{\frac{R}{4}}} w(x) \right)\\
&\leq \tau 2^s [\operatorname{osc}_{B_R}w - \operatorname{osc}_{B_{\frac{R}{4}}}w]\\
&= \tau 2^s \left[\operatorname{osc}_{B_R}u - \operatorname{osc}_{B_{\frac{R}{4}}}u\right]
\end{align*}
this implies inequality \eqref{2.42}. Lemma \ref{Lemma 2.8} is proved.
\end{proof}
We now return to the proof of Theorem \ref{Theorem 2.1}. Combining the result of Lemma \ref{Lemma 2.3} with the conclusion of Lemma \ref{Lemma 2.8}, we deduce that 
\[
u \in C^{0,\alpha}(B_{R_0}(x_0)),
\] where
\[
\alpha = \min\left\{1, -\log_4(1 - \tau^{-1} 2^{-s})\right\} = \alpha(n, g^-, g^+, F, L,\gamma, \delta).
\]
By the arbitrarily of \(x_0 \in \Omega\), we have \(u \in C^{0,\alpha}(\Omega)\). The proof of Theorem \ref{Theorem 2.1} is completed.
\end{proof}
\begin{proof}[\textbf{Proof of Theorem \ref{Theorem 2.2}}]
    Following the same lines as the proof of Theorem \ref{Theorem 2.1}, we can establish the proof of Theorem \ref{Theorem 2.2}.
\end{proof}
\begin{remark}\label{Remark 2.3}
In analogy with Section~7 of Chapter~2 in \cite{OL}, one can define the class 
\[
\mathcal{B}_{G(x,t)}(\Omega \cup S_1, M, \gamma, \gamma_1, \delta),
\]
where \( S_1 \subset \partial \Omega \), and establish the corresponding results.
\end{remark}
\section{Application to Problems in Musielak--Orlicz--Sobolev Spaces}

This section is devoted to the regularity of weak solutions to problem \eqref{P}, with the ultimate goal of proving Theorems~\ref{Theorem 4.3} and~\ref{Theorem 4.4}.

A fundamental step in establishing Hölder continuity ($C^{0,\alpha}$) is to first prove that weak solutions are bounded. Leveraging the embedding results of Diening and Cianchi~\cite{Cianchi2024}, we begin by demonstrating the boundedness of weak solutions to \eqref{P} under both Dirichlet and Neumann boundary conditions, encompassing both critical and subcritical growth conditions on the nonlinearity. This is the subject of Theorems~\ref{Theorem 4.1w} and~\ref{Theorem 6.1w}.

Our strategy for proving boundedness is inspired by the work of Ho, Kim, Winkert, and Zhang~\cite{KKPZ}. We extend their approach to provide a unified treatment for both critical and subcritical cases, thereby yielding a more general result.

\subsection{Boundedness of Weak Solutions to Problems with Dirichlet Boundary Condition}

In this subsection, we focus on proving the boundedness of weak solutions to the Dirichlet problem~\eqref{PD}, as stated in Theorem~\ref{Theorem 4.1w}. To this end, we need the following lemma concerning the geometric convergence of sequences of numbers will be needed for the De Giorgi iteration arguments below. It can be found in Ho-Sim \cite[Lemma 4.3]{Ho-Sim-2015}. The case $\mu_1=\mu_2$ is contained in Lady{\v{z}}enskaja-Solonnikov-Ural{\cprime}ceva \cite[Chapter II, Lemma 5.6]{OL}, see also DiBenedetto \cite[Chapter I, Lemma 4.1]{EDI}.
\begin{lemma}\label{leRecur}
	Let $\{Z_h\}_{h\geq 0}$ be a sequence of positive numbers satisfying the recursion inequality
\[
Z_{h+1} \leq K b^h \left(Z_h^{1+\mu_1} + Z_h^{1+\mu_2}\right), \quad h = 0,1,2,\dots,
\]
for some constants $K > 0$, $b > 1$, and exponents $\mu_2 \geq \mu_1 > 0$. Suppose that either
\[
Z_0 \leq \min \left(1, (2K)^{-\frac{1}{\mu_1}} b^{-\frac{1}{\mu_1^2}} \right),
\]
or
\[
Z_0 \leq \min \left( (2K)^{-\frac{1}{\mu_1}} b^{-\frac{1}{\mu_1^2}}, (2K)^{-\frac{1}{\mu_2}} b^{-\frac{1}{\mu_1 \mu_2} - \frac{\mu_2 - \mu_1}{\mu_2^2}} \right).
\]
Then there exists an index $h_0=0,1,2,\ldots$ such that $Z_h \leq 1$ for all $h \geq h_0$, and moreover:
\[
Z_h \leq \min \left(1, (2K)^{-\frac{1}{\mu_1}} b^{-\frac{1}{\mu_1^2}} b^{-\frac{h}{\mu_1}} \right), \quad \forall h \geq h_0.
\]
In particular, $Z_h \to 0$ as $h \to \infty$.
\end{lemma}
\begin{proof}[\textbf{Proof of Theorem \ref{Theorem 4.1w}}]
The compactness of \( \overline{\Omega} \) implies that for any \( R > 0 \), there exists a finite open cover \( \{B_i(R)\}_{i=1}^{m} \) of balls with radius \( R \), such that
\[
\overline{\Omega} \subset \bigcup_{i=1}^{m} B_i,
\]
where each \( \Omega_i := B_i \cap \Omega \) (for \( i = 1, \ldots, m \)) is also a Lipschitz domain (see, for example, Papageorgiou–Winkert~\cite[ Theorem 1.4.86]{Papageorgiou-Winkert-2018}).

Let $u$ be a weak solution of problem \eqref{PD}. Let $k_*\geq 1$ be sufficiently large such that
	\begin{equation}\label{k*}
		\int_{A_{k_*}}G(x,|D u|){\rm d}x+\int_{A_{k_*}}H_1(x,\vert u\vert){\rm d}x<1,
	\end{equation}
	where $A_k:=\{x\in\Omega\,:\, u(x)>k\}$ for $k\in\mathbb{R}$.\\	
	We define
	\begin{equation}\label{Zn.def}
		Z_h:=\int_{A_{k_h}}G(x,|D u|){\rm d}x+\int_{A_{k_h} }H_1(x,(u-k_{h})){\rm d}x,
	\end{equation}
	where
	\begin{equation*}
		k_h:=k_*\left(2-\frac{1}{2^h}\right)
		\quad \text{for }h =0,1,2,\ldots.
	\end{equation*}
	Obviously, $\text{for }h =0,1,2,\ldots$, it hol{\rm d}s
	\begin{equation}\label{k_n}
		k_h \nearrow 2k_* \quad \text{and} \quad k_* \leq k_h <2k_*.
	\end{equation}
	Since $k_{h}<k_{h+1}$ and $A_{k_{h+1}}\subset A_{k_h}$ for all $h=0,1,2,\ldots$, we have
	\begin{equation}\label{Zn.decreasing}
		Z_{h+1}\leq Z_h\quad\text{for all }h =0,1,2,\ldots.
	\end{equation}
	Moreover, for $x \in A_{k_{h+1}}$,  we see that
	\begin{equation*}
		u(x)- k_h \geq u(x)\left(1-\frac{k_h}{k_{h+1}}\right) = \frac{u(x)}{2^{h+2}-1}.
	\end{equation*}
	Hence, we obtain
	\begin{equation} \label{est.u}
		u(x)\leq (2^{h+2}-1)(u(x)-k_h)\quad \text{for a.\,a.\,}x\in A_{k_{h+1}}\text{ and for all }h=0,1,2,\ldots.
	\end{equation}
	Furthermore, by Propositions \ref{zoo} and \ref{zoo*}, we have
	\begin{align}\label{|A_{k_{h+1}}|}
		\begin{split}
			|A_{k_{h+1}}|
			& \leq \int_{A_{k_{h+1}}}H_1\left (x,\frac{u-k_h}{k_{h+1}-k_h}\right )d x\\
            & \leq \max\left\lbrace \frac{1}{(k_{h+1}-k_h)^{h^-_1}},\frac{1}{(k_{h+1}-k_h)^{h^+_1}} \right\rbrace\int_{A_{k_{h}}} H_1\left (x,u-k_h\right )d x\\
             & \leq \max\left\lbrace \frac{2^{h^-_1(h+1)}}{k_*^{h^-}},\frac{2^{h^+(h+1)}}{k_*^{h^+_1}} \right\rbrace\int_{A_{k_{h}}} H_1\left (x,u-k_h\right )d x\\
             & \leq  \frac{2^{h^+_1(h+1)}}{k_*^{h^-_1}}\int_{A_{k_{h}}} H_1\left (x,u-k_h\right )d x\\
			& \leq \frac{2^{h^+_1(h+1)}}{k_*^{h^-_1}} Z_h\\
			& \leq 2^{h^+_1(h+1)} Z_h\quad \text{for all }h=0,1,2,\ldots.
		\end{split}
	\end{align}
In the remainder of the proof, we denote by \( C_i \) (for \( i = 1, 2, \ldots \)) positive constants that are independent of \( h \) and \( k_* \).\\	
    {\bf Claim 1:}
	There exist positive constants $\mu_1$ and $\mu_2$ such that
	\begin{equation*}%\label{u^p^*}
		\int_{A_{k_{h+1}}}H_1(x,(u-k_{h+1})){\rm d}x
		\leq  C_1 2^{\frac{h\left(h^+_1\right)^2}{g^-}}\left(Z_{h}^{1+\mu_1}+Z_{h}^{1+\mu_2}\right)\quad \text{for all }h=0,1,2,\ldots.
	\end{equation*}	
First, note that
	\begin{align}\label{decompose1}
		\begin{split}
			\int_{A_{k_{h+1}}}H_1(x,(u-k_{h+1})){\rm d}x
			&=\int_{\Omega}H_1(x,(u-k_{h+1})^+){\rm d}x\\
			&\leq \sum_{i=1}^m\int_{\Omega_{i}}H_1(x,(u-k_{h+1})^+){\rm d}x.
		\end{split}
	\end{align}
Let $i\in\{1,\ldots,m\}$. From \eqref{k*}, Propositions \ref{zoo}, \ref{zoo*} and \ref{embb} for $\Omega=\Omega_{i}$ we have
	\begin{align*}
		\int_{\Omega_{i}}H_1(x,(u-k_{h+1})^+){\rm d}x
		&\leq \|(u-k_{h+1})^+\|_{L^{H_1(x,t)}(\Omega_{i})}^{h^-_1}\\
		& \leq C_2\left[\|D (u-k_{h+1})^+\|_{L^{G(x,t)}(\Omega_{i})}+\|(u-k_{h+1})^+\|_{L^{G(x,t)}(\Omega_{i})}\right]^{h^-_1}.
	\end{align*}
Then, again by Propositions \ref{zoo}, \ref{zoo*} and \ref{embb} along with \eqref{k*}, this leads to
	\begin{align*}
		\int_{\Omega_{i}}H_1(x,(u-k_{h+1})^+){\rm d}x&\leq C_3\left(\int_{\Omega_{i}}G(x,|D (u-k_{h+1})^+|){\rm d}x+\int_{\Omega_{i}}G(x,(u-k_{h+1})^+)d x\right)^{\frac{h^-}{g^+}}\\
		&\leq C_3\left(\int_{A_{k_{h+1}}}G(x,|D u|)d x+\int_{A_{k_{h+1}}}H_1(x,(u-k_{h+1})){\rm d}x+|A_{k_{h+1}}|\right)^{\frac{h^-_1}{g^+}}.
	\end{align*}
From this, \eqref{Zn.def}, \eqref{Zn.decreasing} and \eqref{|A_{k_{h+1}}|}, we obtain
	\begin{align*}%\label{grad3}
		\int_{\Omega_{i}}H_1(x,(u-k_{h+1})^+){\rm d}x
		\leq C_42^{\frac{hh^+_1h^-_1}{g^+}}Z_{h}^{\frac{h^-_1}{g^+}}.
	\end{align*}
	Combining this with \eqref{decompose1}, we deduce that 
	\begin{equation*}
		\int_{\Omega}H_1(x,(u-k_{h+1})^+{\rm d}x
		\leq C_52^{\frac{h\left(h^+_1\right)^2}{g^-}}\left(Z_{h}^{1+\mu_1}+Z_{h}^{1+\mu_2}\right),
	\end{equation*}
	where
	\begin{align*}
		0<\mu_1:=\frac{h^-_1}{g^+}-1=:\mu_2.
	\end{align*}
 This proves Claim 1.\\
 {\bf Claim 2:} It hol{\rm d}s that
	\begin{equation*}
		\int_{A_{k_{h+1}}}G(x,|D u|){\rm d}x\leq  C_6 2^{h\left[\frac{\left(h^+\right)^2}{g^-}+h^+\right]}\left(Z_{h-1}^{1+\mu_1}+Z_{h-1}^{1+\mu_2}\right)\quad \text{for all }h=1,2,\ldots.
	\end{equation*}
Testing \eqref{4.9} with $\varphi =(u-k_{h+1})^{+} \in W_0^{1,G(x,t)}(\Omega)$, it yields 
	\begin{align*}
		\int_{\Omega} A(x,u,D u) \cdot D(u-k_{h+1})^{+} d x
		=\int_{\Omega }B(x,u,D u)(u-k_{h+1})^{+} {\rm d}x,
	\end{align*}
which can be written as
	\begin{equation}\label{D.var.Eq}
		\int_{A_{k_{h+1}}} A(x,u,D u) \cdot D u d x
		=\int_{A_{k_{h+1}} }B(x,u,D u)(u-k_{h+1}) d x.
	\end{equation}
We proceed to estimate the terms in \eqref{D.var.Eq}, starting with the left-hand side.

Since \( u \geq u - k_{h+1} > 0 \), and \( u > k_{h+1} \geq 1 \) on the set \( A_{k_{h+1}} \). Furthermore, the function \( t \mapsto H_1(x, t) \) is increasing for all \( x \in \Omega \). Therefore, by assumption \eqref{6.31} and Propositions \ref{zoo} and \ref{zoo*}, we obtain
\begin{align*}
\int_{A_{k_{h+1}}} A(x,u,D u)\cdot D u \, \mathrm{d}x
&\geq a_4 \int_{A_{k_{h+1}}} G(x,|D u|) \, \mathrm{d}x - a_5 \int_{A_{k_{h+1}}} H_1(x,u) \, \mathrm{d}x - \int_{A_{k_{h+1}}} a_6 \, \mathrm{d}x \\
&\geq a_4 \int_{A_{k_{h+1}}} G(x,|D u|) \, \mathrm{d}x 
- \max\{a_5, a_6\} \int_{A_{k_{h+1}}} H_1(x,u) \, \mathrm{d}x.
\end{align*}
We now estimate the right-hand side of \eqref{D.var.Eq}. Let \( \varepsilon \in (0,1] \) be chosen such that  
\[
b_1 \varepsilon^{\frac{h^+_1}{h^+_1 - 1}} \leq \frac{a_4}{4}.
\]  
 Then, employing assumption \eqref{6.51}, Young’s inequality \eqref{Yi}, the estimate \eqref{L1}, and Propositions \ref{zoo} and \ref{zoo*}, we find
\begin{align*}
\int_{A_{k_{h+1}}} B(x,u,D u)(u - k_{h+1}) \, \mathrm{d}x 
&\leq b_1 \int_{A_{k_{h+1}}} M(x,|D u|) u \, \mathrm{d}x 
+ b_2 \int_{A_{k_{h+1}}} h_1(x,u) u \, \mathrm{d}x 
+ \int_{A_{k_{h+1}}} b_3(u - k_{h+1}) \, \mathrm{d}x \\
&\leq b_1 \int_{A_{k_{h+1}}} \frac{\varepsilon}{\varepsilon} M(x,|D u|) u \, \mathrm{d}x 
+ \max\{h^+_1 b_2, b_3\} \int_{A_{k_{h+1}}} H_1(x,u) \, \mathrm{d}x \\
&\leq b_1 \int_{A_{k_{h+1}}} \widetilde{H_1}\left(x, \varepsilon M(x,|D u|) \right) \, \mathrm{d}x 
+ b_1 \int_{A_{k_{h+1}}} H_1\left(x, \frac{1}{\varepsilon} u \right) \, \mathrm{d}x \\
&\quad + \max\{h^+_1 b_2, b_3\} \int_{A_{k_{h+1}}} H_1(x,u) \, \mathrm{d}x \\
&\leq b_1 \varepsilon^{\frac{h^+_1}{h^+_1 - 1}} \int_{A_{k_{h+1}}} \widetilde{H_1}\left(x, M(x,|D u|) \right) \, \mathrm{d}x 
+ b_1 \varepsilon^{-h^+_1} \int_{A_{k_{h+1}}} H_1(x,u) \, \mathrm{d}x \\
&\quad + \max\{h^+_1 b_2, b_3\} \int_{A_{k_{h+1}}} H_1(x,u) \, \mathrm{d}x \\
&\leq b_1 \varepsilon^{\frac{h^+_1}{h^+_1 - 1}} \int_{A_{k_{h+1}}} G(x,|D u|) \, \mathrm{d}x 
+ C_7 \int_{A_{k_{h+1}}} H_1(x,u) \, \mathrm{d}x,
\end{align*}
where 
\[
C_7 := \max\left\{h^+_1 b_2, b_3, b_1 \varepsilon^{-h^+_1} \right\}.
\]
Using the choice of \( \varepsilon \), we conclude that
\[
\int_{A_{k_{h+1}}} B(x,u,D u)(u - k_{h+1}) \, \mathrm{d}x 
\leq \frac{a_4}{4} \int_{A_{k_{h+1}}} G(x,|D u|) \, \mathrm{d}x 
+ C_7 \int_{A_{k_{h+1}}} H_1(x,u) \, \mathrm{d}x.
\]
Combining the estimates above with \eqref{D.var.Eq} and then using \eqref{est.u}, we obtain
	\begin{align*}
		\int_{A_{k_{h+1}}}G(x,|D u|){\rm d}x &
		\leq  C_8\int_{A_{k_{h+1}}}H_1(x,u) {\rm d}x\\
		& \leq C_8\int_{A_{k_{h+1}}}H_1\left(x,(2^{h+2}-1)(u-k_h)\right) {\rm d}x.
	\end{align*}
	Hence,
	\begin{equation*}
		\int_{A_{k_{h+1}}}G(x,|D u|) {\rm d}x
		\leq C_92^{hh^+_1}\int_{\Omega}H_1(x,(u-k_h)^+) {\rm d}x.
	\end{equation*}
Then, Claim~2 follows from the preceding inequality together with Claim~1.\\
From Claims 1 and 2, along with \eqref{Zn.decreasing}, we conclude that
	\begin{equation}\label{Recur}
		Z_{h+1}\leq C_{10} b^h\left(Z_{h-1}^{1+\mu_1}+Z_{h-1}^{1+\mu_2}\right)\quad \text{for all }h=1,2,\ldots,
	\end{equation}
	where
	\begin{align*}
		b:=2^{\big[\frac{\left(h^+_1\right)^2}{g^-}+h^+_1\big]}>1.
	\end{align*}
	This yields to
	\begin{equation*}
		Z_{2(h+1)}\leq C_{10} b^{2h+1}\left(Z_{2h}^{1+\mu_1}+Z_{2h}^{1+\mu_2}\right)\quad \text{for all }h=0,1,2,\ldots,
	\end{equation*}
	that is,
	\begin{equation}\label{Recur1}
		\tilde{Z}_{h+1}\leq bC_{10} \tilde{b}^h\left(\tilde{Z}_h^{1+\mu_1}+\tilde{Z}_h^{1+\mu_2}\right) \quad \text{for all }h=0,1,2,\ldots,
	\end{equation}
	where $\tilde{Z}_h:=Z_{2h}$ and $\tilde{b}:=b^2$. Applying Lemma \ref{leRecur} to \eqref{Recur1} yields
	\begin{equation}\label{Recur+1}
		Z_{2h}=\tilde{Z}_h \to 0 \quad \text {as }  h\to \infty
	\end{equation}
	provided that
	\begin{equation}\label{Z_0}
		\tilde{Z}_{0}\leq \min\left\{(2bC_{10})^{-\frac{1}{\mu_1}}\ \tilde{b}^{-\frac{1}{\mu_1^{2}}},\left(2bC_{10}\right)^{-\frac{1}{\mu_2}}\ \tilde{b}^{-\frac{1}{\mu_1\mu_2}-\frac{\mu_2-\mu_1}{\mu_2^{2}}}\right\}.
	\end{equation}
	Again, from \eqref{Recur} we obtain
	\begin{equation*}
		Z_{2(h+1)+1}\leq C_{10} b^{2(h+1)}\left(Z_{2h+1}^{1+\mu_1}+Z_{2h+1}^{1+\mu_2}\right)\quad \text{for all } h=0,1,2,\ldots,
	\end{equation*}
	which can be written as
	\begin{equation}\label{Recur2}
		\hat{Z}_{h+1}\leq \tilde{b}C_{10} \tilde{b}^h\left(\hat{Z}_h^{1+\mu_1}+\hat{Z}_h^{1+\mu_2}\right) \quad \text{for all }h=1,2,\ldots,
	\end{equation}
	where $\hat{Z}_h:=Z_{2h+1}$. From Lemma \ref{leRecur} applied to \eqref{Recur2} it follows that
	\begin{equation}\label{Recur+2}
		Z_{2h+1}=\hat{Z}_h \to 0 \quad \text{as } n\to \infty
	\end{equation}
	provided that
	\begin{equation}\label{Z_0'}
		\hat{Z}_{0}\leq \min\left\{(2\tilde{b}C_{10})^{-\frac{1}{\mu_1}}\ \tilde{b}^{-\frac{1}{\mu_1^{2}}},\left(2\tilde{b}C_{10}\right)^{-\frac{1}{\mu_2}}\ \tilde{b}^{-\frac{1}{\mu_1\mu_2}-\frac{\mu_2-\mu_1}{\mu_2^{2}}}\right\}.
	\end{equation}
	Note that
	\begin{align*}
		\hat{Z}_{0}=Z_1\leq Z_0=\tilde{Z}_0\leq \int_{A_{k_*}}G(x,|D u|) d x+\int_{A_{k_*} }H(x,u){\rm d}x.
	\end{align*}
	Therefore, by choosing $k_*>1$ sufficiently large we have
	\begin{align*}
		\int_{A_{k_*}}G(x,|D u|) d x+\int_{A_{k_*} }H_1(x,u){\rm d}x
		&\leq \min\left\{1,(2\tilde{b}C_{10})^{-\frac{1}{\mu_1}}\ \tilde{b}^{-\frac{1}{\mu_1^{2}}},\left(2\tilde{b}C_{10}\right)^{-\frac{1}{\mu_2}}\ \tilde{b}^{-\frac{1}{\mu_1\mu_2}-\frac{\mu_2-\mu_1}{\mu_2^{2}}}\right\}.
	\end{align*}
	Hence, \eqref{k*}, \eqref{Z_0} and \eqref{Z_0'} are fulfilled and we obtain \eqref{Recur+1} and \eqref{Recur+2}. This means that
	\begin{align*}
		Z_h=\int_{A_{k_h}}G(x,|D u|) d x+\int_{A_{k_h} }H_1(x,(u-k_h)) d x\to 0 \quad \text{as } h\to\infty.
	\end{align*}
	In particular, we have
	\begin{align*}
		\int_{\Omega}H_1(x,(u-2k_*)^+)d x=0.
	\end{align*}
	Consequently, by Proposition \ref{zoo}, $(u-2k_*)^{+}=0$ a.\,e.\,in $\Omega$ and so
	\begin{align*}
		\underset{\Omega}{\mathop{\rm ess\,sup}}\ u \leq 2k_*.
	\end{align*}
Replacing $u$ by $-u$ in the arguments above, we also obtain
	\begin{align*}
		\underset{\Omega}{\mathop{\rm ess\,sup}}\ (-u) \leq 2k_*.
	\end{align*}
Hence, $\|u\|_\infty\leq 2k_*$. This finishes the proof.
\end{proof}
\subsection{Boundedness of Weak Solutions to Problems with Neumann Boundary Condition}\label{subsection3}

In this subsection, we address the boundedness of weak solutions to the Neumann problem~\eqref{PN}, as formulated in Theorem~\ref{Theorem 6.1w}.

\begin{proof}[\textbf{Proof of Theorem \ref{Theorem 6.1w}}]
	As before, since $\overline{\Omega}$ is compact, for any \( R > 0 \), there exists a finite open cover \( \{B_i(R)\}_{i=1}^m \) of balls \( B_i := B_i(R) \), each of radius \( R \), such that 
\[
\overline{\Omega} \subset \bigcup_{i=1}^{m} B_i,
\]
and each intersection \( \Omega_i := B_i \cap \Omega \) (for \( i = 1, \ldots, m \)) is also a Lipschitz domain.\\ 
We denote by \( I \subset \{1, \ldots, m\} \) the set of indices for which \( B_i \cap \partial\Omega \neq \emptyset \), and we choose \( R > 0 \) sufficiently small.\\
Let \( u \) be a weak solution to problem~\eqref{PN}, and let \( k_* \geq 1 \) be sufficiently large such that 
	\begin{equation}\label{N.k*11}
		\int_{A_{k_*}}G(x,|D u|){\rm d}x+\int_{A_{k_*}}H_1(x,u){\rm d}x+\int_{\Gamma_{k_*}}H_2(x,u){\rm d}\sigma<1,
	\end{equation}
	where
	\begin{align*}
		A_k:=\{x\in\Omega \, : \, u(x)>k\} \quad\text{and}\quad \Gamma_k:=\{x\in\partial\Omega \,:\, u(x)>k\} \quad\text{for }k\in\mathbb{R}.
	\end{align*}
	We define
	\begin{equation}\label{N.Zn.def11}
		Z_h:=\int_{A_{k_h}}G(x,|D u|){\rm d}x+\int_{A_{k_h} }H_1\left(x,(u-k_h)\right){\rm d}x+\int_{\Gamma_{k_n}}H_2\left(x,(u-k_h)\right){\rm d}\sigma,
	\end{equation}
	where $\{k_{h}\}_{h\geq 0}$ is a sequence defined as in the proof of Theorem \ref{Theorem 4.1w}. Since $k_{h}<k_{h+1}$, we have $A_{k_{h+1}}\subset A_{k_{h}}$ and $\Gamma_{k_{h+1}}\subset \Gamma_{k_{h}}$ for all $h=0,1,2,\ldots$. Hence, it hol{\rm d}s
	\begin{equation}\label{N.Zn.decreasing11}
	Z_{h+1}\leq Z_h\quad\text{for all }h=0,1,2,\ldots.
	\end{equation}
	Moreover, we have the following estimates (see the proof of Theorem \ref{Theorem 4.1w})
	\begin{align}
		u(x)&\leq \left(2^{h+2}-1\right)(u(x)-k_{h})\quad \text{for a.\,a.\,}x\in A_{k_{h+1}} \text{ and for all } h=0,1,2,\ldots,\label{N.est.u111}\\
		u(x)&\leq \left(2^{h+2}-1\right)(u(x)-k_{h})\quad \text{for a.\,a.\,}x\in \Gamma_{k_{h+1}}\text{ and for all } h=0,1,2,\ldots. \label{N.est.u211}
	\end{align}
	As shown in \eqref{|A_{k_{h+1}}|} the following estimate hol{\rm d}s true
	\begin{equation} \label{N.|A_k|1}
	|A_{k_{h+1}}|  \leq \frac{2^{h_1^+(h+1)}}{k_*^{h^-_1}} Z_h\leq 2^{h^+_1(h+1)} Z_h\quad \text{for all }h=0,1,2,\ldots.
	\end{equation}
	In the following, we will denote by $C_i$ ($i=1,2,3,\ldots$) positive constants independent of $h$ and $k_*$.\\
	{\bf Claim 1:} It hol{\rm d}s that
	\begin{equation*}%\label{N.u^p^*}
		\int_{A_{k_{h+1}}}H_1\left(x,(u-k_{h+1})\right){\rm d}x\leq  C_1 2^{\frac{h\left(h^+_1\right)^2}{g^-}}\left(Z_{h}^{1+\nu_1}+Z_{h}^{1+\nu_2}\right)\quad \text{for all } h=0,1,2,\ldots,
	\end{equation*}
	where
	\begin{align*}
		0<\nu_1:=\frac{h_1^-}{g^+}-1=:\nu_2.
	\end{align*}
	The proof of Claim 1 is the same as Claim 1 in Theorem \ref{Theorem 4.1w}.\\	
	{\bf Claim 2:} There exist positive constants $\nu_3$ and $\nu_4$ such that
	\begin{equation*}%\label{N.u^p_*.Ga}
		\int_{\Gamma_{k_{h+1}}}H_2\left(x,(u-k_{h+1})\right){\rm d} \sigma\leq  C_2 2^{\frac{h\left(h^+_2\right)^2}{g^-}}\left(Z_{h}^{1+\nu_3}+Z_{h}^{1+\nu_4}\right)\quad \text{for all } h=0,1,2,\ldots.
	\end{equation*}
	First, we see that
	\begin{align}\label{N.deco.Bdr11}
		\begin{split}
			\int_{\Gamma_{k_{h+1}}}H_2(x,(u-k_{h+1})){\rm d} \sigma
			&=\int_{\partial\Omega}H_2(x,(u-k_{h+1})^+){\rm d} \sigma\\
			& \leq \sum_{i=1}^m\int_{\partial\Omega\cap \partial \Omega_{i}}H_2(x,(u-k_{h+1})^+){\rm d} \sigma.
		\end{split}
	\end{align}
	Let \( i \in \{1, \ldots, m\} \). From Propositions~\ref{zoo} and \ref{zoo*}, inequality \eqref{N.k*11}, and Proposition~\ref{embb}, for \( \Omega = \Omega_i \), we have
	\begin{align*}
		& \int_{\partial\Omega\cap \partial \Omega_{i}}H_2(x,(u-k_{h+1})^+){\rm d} \sigma\\
		& \leq \|(u-k_{h+1})^+\|_{L^{H_2(x,t)}(\partial\Omega\cap \partial \Omega_{i})}^{h_2^-}\\
        & \leq \|(u-k_{h+1})^+\|_{L^{H_2(x,t)}( \partial \Omega_{i})}^{h_2^-}\\
		&\leq C_3\left[\|D (u-k_{n+1})^+\|_{L^{G(x,t)}(\Omega_{i})}+\|(u-k_{h+1})^+\|_{L^{G(x,t)}(\Omega_{i})}\right]^{h_2^-}.
	\end{align*}
	Then, by Propositions \ref{zoo}, \ref{zoo*} and inequality\eqref{N.k*11} it follows
	\begin{align*}
		& \int_{\partial\Omega\cap \partial \Omega_{i}}H_2(x,(u-k_{h+1})^+){\rm d} \sigma\\
		&\leq C_4\left(\int_{\Omega_{i}}G(x,|\nabla (u-k_{h+1})^+|){\rm d}x+\int_{\Omega_{i}}G(x,(u-k_{h+1})^+){\rm d}x\right)^{\frac{h_2^-}{g^+}}\\
	\notag&\leq C_4\left(\int_{A_{k_{h+1}}}G(x,|D u|){\rm d}x+\int_{A_{k_{h+1}}}H_2(x,(u-k_{h+1})){\rm d}x+|A_{k_{h+1}}|\right)^{\frac{h_2^-}{g^+}}.
	\end{align*}
	From this combined with \eqref{N.Zn.def11}, \eqref{N.Zn.decreasing11} as well as \eqref{N.|A_k|1} we obtain
	\begin{align}\label{N.grad311}
		\int_{\partial\Omega\cap \partial \Omega_{i}}H_2(x,(u-k_{h+1})^+){\rm d} \sigma&\leq C_52^{\frac{hh_2^+h_2^-}{g^+}}Z_{h}^{\frac{h_2^-}{g^+}}.
	\end{align}
Consequently, by \eqref{N.deco.Bdr11} and \eqref{N.grad311}, we conclude that
	\begin{equation*}
		\int_{\partial\Omega}H_2((u-k_{h+1})^+){\rm d}\sigma\leq C_62^{\frac{h\left(h_2^+\right)^2}{g^-}}\left(Z_{h}^{1+\nu_3}+Z_{h}^{1+\nu_4}\right),
	\end{equation*}
	where
	\begin{align*}
		0<\nu_3:=\frac{h_2^-}{g^+}-1=:\nu_4.
	\end{align*}
 This proves Claim 2.\\	
	{\bf Claim 3:} It hol{\rm d}s that
	\begin{equation*}
		\int_{A_{k_{h+1}}}G(x,|D u|){\rm d}x\leq  C_7 2^{h\left[\frac{\left(h_1^++h_2^+\right)^2}{g^-}+g_*^+\right]}\left(Z_{h-1}^{1+\mu_1}+Z_{h-1}^{1+\mu_2}\right)\quad \text{for all }h=1,2,3,\ldots,
	\end{equation*}
	where
	\begin{align*}
		0<\mu_1:=\nu_1=\nu_2=\nu_3=\nu_4=:\mu_2.
	\end{align*}
	Taking $\varphi =(u-k_{h+1})^{+} \in W^{1,G(x,t)}(\Omega)$ as test function in \eqref{FNV} gives
	\begin{align*}
		\int_{\Omega}
		A(x,u,D u) \cdot \nabla
		\varphi {\rm d}x =\int_{\Omega }B(x,u,D u)\varphi {\rm d}x+\int_{\partial\Omega}C(x,u)\varphi {\rm d}\sigma,
	\end{align*}
	which can be written as
	\begin{align}\label{estimate_311}
		\begin{split}
			&\int_{A_{k_{h+1}}} A(x,u,D u) \cdot D u {\rm d}x\\ 
			&=\int_{A_{k_{h+1}} }B(x,u,D u)(u-k_{h+1}) {\rm d}x +\int_{\Gamma_{k_{h+1}} }C(x,u)(u-k_{h+1}) {\rm d} \sigma.
		\end{split}
	\end{align}
	As done in the proof of Theorem \ref{Theorem 4.1w}, we apply assumptions \eqref{6.31}, \eqref{6.51} and \eqref{6.61}, as well as Young's inequality \eqref{Yi}, \eqref{L1} and Propositions \ref{zoo} and \ref{zoo*}, and we take $\varepsilon \in (0,1]$ such that
    $$b_1\varepsilon^{\frac{h_1^+}{h_1^+-1}}\leq\frac{a_4}{4},$$
  we have the estimates
	\begin{align}\label{estimate_411}
		\begin{split}
			&\int_{A_{k_{h+1}}} A(x,u,D u)\cdot D u {\rm d}x\\
			&\geq a_4\int_{A_{k_{h+1}} }G(x,|D u|) {\rm d}x-\max\{a_5,a_6\}\int_{A_{k_{h+1}} }H_1(x,u){\rm d}x,
		\end{split}
	\end{align}
	and
	\begin{align}\label{estimate_511}
		\begin{split}
			&\int_{A_{k_{h+1}}} B(x,u,D u)(u-k_{h+1}){\rm d}x\\
            & \leq b_1\int_{A_{k_{h+1}}}  M(x,|D u|)(u-k_{h+1}) {\rm d}x + b_2 \int_{A_{k_{h+1}}}  h_1(x, u) (u-k_{h+1}){\rm d}x+ b_3\int_{A_{k_{h+1}}} (u-k_{h+1}) d x\\
            & \leq b_1\int_{A_{k_{h+1}}}  M(x,|D u|)u {\rm d}x + b_2 \int_{A_{k_{h+1}}}  h_1(x, u) u{\rm d}x+b_3\int_{A_{k_{h+1}}}  u(x) d x\\
             & \leq b_1\int_{A_{k_{h+1}}}  M(x,|D u|)u {\rm d}x + b_2 h_1^+\int_{A_{k_{h+1}}}  H_1(x, u) {\rm d}x+b_3\int_{A_{k_{h+1}}}  H_1(x,u) d x\\
			& \leq b_1\int_{A_{k_{h+1}} }\frac{\varepsilon}{\varepsilon}M(x,|D u|)u{\rm d}x+\max\{h_1^+b_2,b_3\}\int_{A_{k_{h+1}}}H_1(x,u){\rm d}x \\
        & \leq b_1\int_{A_{k_{h+1}} }\widetilde{H_1}\left(x,\varepsilon M(x,|D u|)\right){\rm d}x + b_1\int_{A_{k_{h+1}} }H_1\left(x,\frac{1}{\varepsilon}u\right){\rm d}x\\
        &+\max\{h_1^+b_2,b_3\}\int_{A_{k_{h+1}}}H_1(x,u){\rm d}x \\
        & \leq b_1\varepsilon^{\frac{h_1^+}{h_1^+-1}}\int_{A_{k_{h+1}} }\widetilde{H_1}\left(x, M(x,|D u|)\right){\rm d}x + b_1\varepsilon^{-h_1^+}\int_{A_{k_{h+1}} }H_1\left(x,u\right){\rm d}x\\
        &+\max\{h_1^+b_2,b_3\}\int_{A_{k_{h+1}}}H_1(x,u){\rm d}x \\
        & \leq b_1\varepsilon^{\frac{h_1^+}{h_1^+-1}}\int_{A_{k_{h+1}} }G(x,|D u|){\rm d}x +\max\{h_1^+b_2,b_3,b_1\varepsilon^{-h_1^+}\}\int_{A_{k_{h+1}}}H_1(x,u){\rm d}x\\ 
	& = \frac{a_4}{4}\int_{A_{k_{h+1}} }G(x,|D u|){\rm d}x +C_7\int_{A_{k_{h+1}}}H_1(x,u){\rm d}x.
		\end{split}
	\end{align}
 Note that $u\geq u-k_{h+1} >0$ and $u> k_{h+1} \geq 1$ on $\Gamma_{k_{h+1}}$. So, applying assumption \eqref{6.61}, it yields that
	\begin{align}\label{estimate_611}
		\begin{split}
			\int_{\Gamma_{k_{h+1}} }C(x,u)(u-k_{h+1}) {\rm d}\sigma
			& \leq c_1 \int_{\Gamma_{k_{h+1}} }h_2(x,u) (u-k_{h+1}){\rm d}\sigma+c_2 \int_{\Gamma_{k_{h+1}} }(u-k_{h+1}) {\rm d}\sigma\\
            & \leq c_1 \int_{\Gamma_{k_{h+1}} }h_2(x,u) u{\rm d}\sigma+c_2 \int_{\Gamma_{k_{h+1}} }u {\rm d}\sigma\\
            & \leq c_1 h_2^+ \int_{\Gamma_{k_{h+1}} }H_2(x,u) {\rm d}\sigma+c_2 \int_{\Gamma_{k_{h+1}} }H_2(x,u) {\rm d}\sigma\\
			& \leq (c_1h_2^++c_2)\int_{\Gamma_{k_{h+1}} }H_2(x,u) {\rm d}\sigma.
		\end{split}
	\end{align}
	Combining \eqref{estimate_311}, \eqref{estimate_411}, \eqref{estimate_511} and \eqref{estimate_611}, and using \eqref{N.est.u111} as well as \eqref{N.est.u211}, we obtain
	\begin{align*}
		&\int_{A_{k_{h+1}}}G(x,|D u|){\rm d}x\\
		& \leq  C_8\int_{A_{k_{h+1}}}H_1(x,u){\rm d}x+C_8\int_{\Gamma_{k_{h+1}}}H_2(x,u){\rm d} \sigma\\
		&\leq C_8\int_{A_{k_{h+1}}}H_1\left(x,\left(2^{h+2}-1\right)(u-k_{h})\right){\rm d}x\\
		&\qquad +C_8\int_{\Gamma_{k_{h+1}}}H_2\left(x,\left(2^{h+2}-1\right)(u-k_{h})\right){\rm d} \sigma.
	\end{align*}
	Hence,
	\begin{equation*}
		\int_{A_{k_{h+1}}}G(x,|D u|){\rm d}x
		\leq C_92^{h(h_1^++h_2^+)}\left[\int_{A_{k_{h}}}H_1(x,(u-k_{h})){\rm d}x+\int_{\Gamma_{k_h}}H_2(x,(u-k_{h}) ){\rm d}\sigma\right].
	\end{equation*}
	Then, Claim 3 follows from the last inequality and Claims 1 and 2.\\
	From Claims 1, 2, and 3, along with \eqref{N.Zn.decreasing11} one has
	\begin{equation}\label{N.Recur1}
		Z_{h+1}\leq C_{10} b^h\left(Z_{h-1}^{1+\mu_1}+Z_{h-1}^{1+\mu_2}\right)\quad \text{for all } h=1,2,3,\ldots,
	\end{equation}
	where
	\begin{align*}
		b:=2^{\left[\frac{\left(h_1^++h_2^+\right)^2}{g^-}+g_*^+\right]}>1.
	\end{align*}
	Repeating the same arguments used in the proof of Theorem~\ref{Theorem 4.1w}, by choosing \( k_* > 1 \) sufficiently large such that
	\begin{align*}
		&\int_{A_{k_*}}G(x,|D u|){\rm d}x+\int_{A_{k_*} }H_1(x,u){\rm d}x+\int_{\Gamma_{k_*} }H_2(x,u){\rm d} \sigma\\
		&\leq \min\left\{1,(2\tilde{b}C_{10})^{-\frac{1}{\mu_1}}\ \tilde{b}^{-\frac{1}{\mu_1^{2}}},\left(2\tilde{b}C_{10}\right)^{-\frac{1}{\mu_2}}\ \tilde{b}^{-\frac{1}{\mu_1\mu_2}-\frac{\mu_2-\mu_1}{\mu_2^{2}}}\right\},
	\end{align*}
	where $\tilde{b}:=b^2$, we deduce from \eqref{N.Recur1} that
	\begin{align*}
		Z_h=\int_{A_{k_{h}}}G(x,|D u|){\rm d}x+\int_{A_{k_{h}} }H_1(x,(u-k_{h})){\rm d}x+\int_{\Gamma_{k_{h}} }H_2(x,(u-k_{h})){\rm d} \sigma\to 0
	\end{align*}
	as $h\to \infty$. This implies that
	\begin{align*}
		\int_{\Omega}H_1(x,(u-2k_*)^+) {\rm d}x+\int_{\partial\Omega}H_2(x,(u-2k_*)^+) {\rm d} \sigma=0.
	\end{align*}
	Therefore, $(u-2k_*)^{+}=0$ a.\,e.\,in $\Omega$ and $(u-2k_*)^{+}=0$ a.\,e.\,on $\partial\Omega$. This means that
	\begin{align*}
		\underset{\Omega}{\mathop{\rm ess\,sup}}\ u +\underset{\partial\Omega}{\mathop{\rm ess\,sup}}\ u\leq 4k_*.
	\end{align*}
By replacing \( u \) with \( -u \) in the arguments above, we can show in the same way that
\[
\underset{\Omega}{\mathop{\rm ess\,sup}} (-u) + \underset{\partial\Omega}{\mathop{\rm ess\,sup}} (-u) \leq 4k_*.
\]
Hence,
\[
\|u\|_{L^\infty(\Omega)} + \|u\|_{L^\infty(\partial\Omega)} \leq 4k_*.
\]
The proof is complete.
\end{proof}

%----------------------------------------------------

%%***********************************************************************************************************************************

%***********************************************************************************************************************************
%------------------
\subsection{Hölder Continuity of Weak Solutions to Problems with Dirichlet or Neumann Boundary Conditions}
We now turn to the Hölder continuity of weak solutions to equation~\eqref{P}. Since the boundedness condition~\eqref{M} required in Theorem~\ref{Theorem 4.2} is ensured by Theorems~\ref{Theorem 4.1w} and~\ref{Theorem 6.1w}, we are in a position to establish the Hölder regularity of weak solutions to Problem~\eqref{P} under both Dirichlet~\eqref{PD} and Neumann~\eqref{PN} boundary conditions. This result is summarized in Theorems~\ref{Theorem 4.2}–\ref{Theorem 4.4}.
\begin{proof}[\textbf{Proof of Theorem \ref{Theorem 4.2}}]
Let \(\bar{B}_t \subset B_s \subset \Omega\) and select a \(C^\infty\)-function \(\xi\) satisfying
\begin{equation}\label{55eq}
0 \leq \xi \leq 1, \quad \text{supp } \xi \subset B_s, \quad \xi \equiv 1 \text{ on } B_t, \quad |D\xi| \leq \frac{2}{s - t}.
\end{equation}
Define
\begin{equation}\label{4.29}
\delta = \min \left\{ \frac{a_4}{4b_3M}, 2 \right\}
\end{equation}
and set \(v := \xi^{g^+} (u - k)^+ = \xi^{g^+} \max \{u - k, 0\}\), where
\begin{equation}\label{4.30}
k \geq \max_{B_s} u(x) - \delta M.
\end{equation}
Clearly, \(v \in W_0^{1,G(x,t)}(\Omega)\). Substituting into equation \eqref{4.9} yields
\begin{align}\label{eq55}
      \int_{A_{k,s}} \xi^{g^+} A(x, u, Du) \cdot Du  \mathrm{d}x &+ \int_{A_{k,s}} g^+\xi^{g^+-1} (u - k) A(x, u, Du) \cdot D\xi  \mathrm{d}x \\
      &= \int_{A_{k,s}} B(x, u, Du) \xi^{g^+} (u - k)  \mathrm{d}x.    
\end{align}
In light of conditions \eqref{6.31}--\eqref{6.41} and \eqref{6.51}, we derive the inequality
\begin{align}\label{4.31}
\begin{split}
a_4 \int_{A_{k,s}} G(x,|Du|) \xi^{g^+}  \mathrm{d}x &\leq a_5 \int_{A_{k,s}} H_1(x,|u|) \xi^{g^+}  \mathrm{d}x + a_6 \int_{A_{k,s}} \xi^{g^+}  \mathrm{d}x \\
&\quad + a_1 g^+ \int_{A_{k,s}} g(x,|Du|) \xi^{g^+-1} |D\xi|(u - k)  \mathrm{d}x \\
&\quad + a_2 g^+\int_{A_{k,s}} P(x,|u|) \xi^{g^+-1} |D\xi|(u - k)  \mathrm{d}x \\
&\quad + a_3g^+ \int_{A_{k,s}} |D\xi| \xi^{g^+-1}(u - k)  \mathrm{d}x \\
&\quad + b_1 \int_{A_{k,s}}M(x,|Du|) \xi^{g^+} (u - k)  \mathrm{d}x \\
&\quad + b_2 \int_{A_{k,s}}h_1(x,|u|) \xi^{g^+} (u - k)  \mathrm{d}x \\
&\quad + b_3 \int_{A_{k,s}} \xi^{g^+} (u - k)  \mathrm{d}x.
\end{split}
\end{align}
We proceed to estimate each term on the right-hand side of \eqref{4.31}.

 Employing \eqref{L1} and Young's inequality \eqref{Yi}, select \(\varepsilon_1 \in (0,1)\) such that
\[
a_1 (g^+)^2 \varepsilon_1^{\frac{g^+}{g^+-1}} \leq \frac{a_4}{4},
\]
 gives that
\begin{align}\label{4.32}
\begin{split}
&a_1 g^+ \int_{A_{k,s}} g(x,|Du|) \xi^{g^+-1} |D\xi|(u - k)  \mathrm{d}x \\
&\quad \leq a_1 g^+ \int_{A_{k,s}} \widetilde{G}\left(x, g(x,|Du|) \xi^{g^+-1}\varepsilon_1\right)  \mathrm{d}x + a_1 g^+ \int_{A_{k,s}} G\left(x, \frac{|D\xi|(u - k)}{\varepsilon_1}\right)  \mathrm{d}x \\
&\quad \leq a_1 g^+ \varepsilon_1^{\frac{g^+}{g^+-1}} \int_{A_{k,s}} \widetilde{G}\left(x, g(x,|Du|) \right) \xi^{g^+}  \mathrm{d}x + a_1 g^+ \varepsilon_1^{-g^+} \int_{A_{k,s}} G\left(x, \frac{2(u - k)}{s-t}\right)  \mathrm{d}x \\
&\quad \leq a_1 (g^+)^2 \varepsilon_1^{\frac{g^+}{g^+-1}} \int_{A_{k,s}} G(x,|Du|)  \xi^{g^+}  \mathrm{d}x + a_1 g^+ \varepsilon_1^{-g^+} 2^{g^+} \int_{A_{k,s}} G\left(x, \frac{u - k}{s-t}\right)  \mathrm{d}x \\
&\quad \leq \frac{a_4}{4} \int_{A_{k,s}} G(x,|Du|) \xi^{g^+}  \mathrm{d}x + d_1 \int_{A_{k,s}} G\left(x, \frac{u - k}{s - t} \right)  \mathrm{d}x,
\end{split}
\end{align}
with \(d_1 = d_1(a_1, a_2, a_3, g^-, g^+) > 0\).

 An application of Young's inequality \eqref{Yi} and Proposition \ref{zoo}, with \(\varepsilon_2 \in (0,1)\) chosen so that
\[
a_3g^+ 2^{g^+} \varepsilon_2^{g^-} \leq 1,
\]
yields that
\begin{align}\label{4.33}
\begin{split}
&a_3g^+ \int_{A_{k,s}} |D\xi|(u - k) \xi^{g^+-1}  \mathrm{d}x \leq a_3g^+ \int_{A_{k,s}} |D\xi|(u - k)  \mathrm{d}x \\
&\quad \leq a_3g^+ \int_{A_{k,s}} G\left(x, \varepsilon_2 |D\xi| (u - k)\right)  \mathrm{d}x + a_3g^+ \int_{A_{k,s}} \widetilde{G}\left(x, \frac{1}{\varepsilon_2} \right)  \mathrm{d}x \\
&\quad \leq a_3g^+\varepsilon_2^{g^-} 2^{g^+} \int_{A_{k,s}} G\left(x, \frac{u - k}{s-t}\right)  \mathrm{d}x + a_3g^+ \int_{A_{k,s}} \widetilde{G}\left(x, \frac{1}{\varepsilon_2} \right)  \mathrm{d}x \\
&\quad \leq \int_{A_{k,s}} G\left(x, \frac{u - k}{s - t} \right)  \mathrm{d}x + c_1 |A_{k,s}|,
\end{split}
\end{align}
where \(c_1 = c_1(a_3, g^+) > 0\).

 From \eqref{H111}, \eqref{M}, and Proposition \ref{zoo}, we find
\begin{align}\label{11eq}
    \begin{split}
        a_5 \int_{A_{k,s}} H_1(x,|u|) \xi^{g^+}  \mathrm{d}x &\leq a_5 \int_{A_{k,s}} H_1(x, M) \xi^{g^+}  \mathrm{d}x \\
        &\leq a_5 \int_{A_{k,s}} F_1 M^{h_1^+} \xi^{g^+}  \mathrm{d}x \leq c_2 |A_{k,s}|,
    \end{split}
\end{align}
with \(c_2 = c_2(a_5, F_1, h_1^+, M) > 0\).

 Invoking Young's inequality \eqref{Yi}, Proposition \ref{zoo}, inequalities \eqref{6.81}, \eqref{11eq}, and selecting \(\varepsilon_3 \in (0,1)\) such that
\[
a_2g^+ 2^{g^+} \varepsilon_3^{-g^+} \leq 1,
\]
we obtain
\begin{align}\label{22eq}
    \begin{split}
        &a_2 g^+\int_{A_{k,s}} P(x,|u|) \xi^{g^+-1} |D\xi|(u - k)  \mathrm{d}x \\
        &\quad \leq a_2 g^+ \int_{A_{k,s}} \widetilde{G}\left(x, P(x,|u|) \xi^{g^+-1}\varepsilon_3\right)  \mathrm{d}x + a_2 g^+ \int_{A_{k,s}} G\left(x, \frac{|D\xi|(u - k)}{\varepsilon_3}\right)  \mathrm{d}x \\
        &\quad \leq a_2 g^+ \varepsilon_3^{\frac{g^+}{g^+-1}} \int_{A_{k,s}} \widetilde{G}\left(x, P(x,|u|) \right) \xi^{g^+}  \mathrm{d}x + a_2 g^+ \varepsilon_3^{-g^+} \int_{A_{k,s}} G\left(x, \frac{2(u - k)}{s-t}\right)  \mathrm{d}x \\
        &\quad \leq a_2 g^+ \varepsilon_3^{\frac{g^+}{g^+-1}} \int_{A_{k,s}} H_1(x,|u|)  \xi^{g^+}  \mathrm{d}x + a_2 g^+ \varepsilon_3^{-g^+} 2^{g^+} \int_{A_{k,s}} G\left(x, \frac{u - k}{s-t}\right)  \mathrm{d}x \\
        &\quad \leq c_3 |A_{k,s}| + d_2 \int_{A_{k,s}} G\left(x, \frac{u - k}{s - t} \right)  \mathrm{d}x,
    \end{split}
\end{align}
where \(c_3 = c_3(a_2, F_1, h_1^+, M, g^+) > 0\) and \(d_2 = d_2(a_2, g^+) > 0\).

 In light of \eqref{6.71}, \eqref{H111}, \eqref{M}, Proposition \ref{zoo}, and Young's inequality \eqref{Yi}, choose \(\varepsilon_4 \in (0,1)\) such that
\[
b_1 \varepsilon_4^{\frac{h^+_1}{h^+_1-1}} \leq \frac{a_4}{4},
\]
we see that
\begin{align}\label{33eq}
    \begin{split}
        &b_1 \int_{A_{k,s}} M(x,|Du|) \xi^{g^+} (u - k)  \mathrm{d}x \\
        &\quad \leq b_1 \int_{A_{k,s}} \widetilde{H_1}\left(x, \varepsilon_4 M(x,|Du|) \xi^{g^+}\right)  \mathrm{d}x + b_1 \int_{A_{k,s}} H_1\left(x, \frac{u - k}{\varepsilon_4}\right)  \mathrm{d}x \\
        &\quad \leq b_1 \varepsilon_4^{\frac{h^+_1}{h^+_1-1}} \int_{A_{k,s}} \widetilde{H_1}\left(x, M(x,|Du|) \right) \xi^{g^+\frac{h^+_1}{h^+_1-1}}  \mathrm{d}x + b_1 \int_{A_{k,s}} H_1\left(x, \frac{\delta M}{\varepsilon_4}\right)  \mathrm{d}x \\
        &\quad \leq b_1 \varepsilon_4^{\frac{h^+_1}{h^+_1-1}} \int_{A_{k,s}} G(x,|Du|) \xi^{g^+}  \mathrm{d}x + b_1 (\delta M)^{h_1^+} \varepsilon_4^{-h^+_1} F_1 |A_{k,s}| \\
        &\quad \leq \frac{a_4}{4} \int_{A_{k,s}} G(x,|Du|) \xi^{g^+}  \mathrm{d}x + c_4 |A_{k,s}|,
    \end{split}
\end{align}
with \(c_4 = c_4(b_1, M, \delta, F_1, h_1^+) > 0\).

 From \eqref{L1}, \eqref{H111}, \eqref{M}, Proposition \ref{zoo}, and Young's inequality \eqref{Yi}, we get
\begin{align}\label{44eq}
    \begin{split}
        &b_2 \int_{A_{k,s}} h_1(x,|u|) \xi^{g^+} (u - k)  \mathrm{d}x \\
        &\quad \leq b_2 \int_{A_{k,s}} \widetilde{H_1}\left(x, h_1(x,|u|)\right)  \mathrm{d}x + b_2 \int_{A_{k,s}} H_1(x, (u - k))  \mathrm{d}x \\
        &\quad \leq b_2 h^+_1 \int_{A_{k,s}} H_1\left(x, |u|\right)  \mathrm{d}x + b_2 \int_{A_{k,s}} H_1(x, \delta M)  \mathrm{d}x \\
        &\quad \leq b_2 h^+_1 F_1 M^{h_1^+} |A_{k,s}| + b_2 F_1 (\delta M)^{h^+_1} |A_{k,s}| \\
        &\quad \leq \left(b_2 F_1 M^{h_1^+}\right) \left(\delta ^{h^+_1} + h_1^+ \right) |A_{k,s}| \leq c_5 |A_{k,s}|,
    \end{split}
\end{align}
where \(c_5 = c_5(b_2, M, \delta, F_1, h_1^+) > 0\).

 The remaining terms are bounded directly. From \eqref{4.29} and \eqref{4.30}, we have
\begin{align}
    a_6 \int_{A_{k,s}} \xi^{g^+}  \mathrm{d}x &\leq a_6 |A_{k,s}|, \\
    \int_{A_{k,s}} \xi^{g^+} (u - k)  \mathrm{d}x &\leq \int_{A_{k,s}} (u - k)  \mathrm{d}x \leq (\max_{B_s} u(x) - k) |A_{k,s}| \leq \delta M |A_{k,s}| \leq \frac{a_4}{4b_3} |A_{k,s}|. \label{4.35}
\end{align}

Combining estimates \eqref{4.31}--\eqref{4.35}, we conclude that
\[
\int_{A_{k,s}} G(x,|Du|) \xi^{g^+}  \mathrm{d}x \leq \gamma \int_{A_{k,s}} G\left(x, \frac{u - k}{s - t} \right)  \mathrm{d}x + \gamma_1 |A_{k,s}|,
\]
and consequently,
\begin{equation}\label{4.36}
\int_{A_{k,t}} G(x,|Du|)  \mathrm{d}x \leq \gamma \int_{A_{k,s}} G\left(x, \frac{u - k}{s - t} \right)  \mathrm{d}x + \gamma_1 |A_{k,s}|,
\end{equation}
for \(\bar{B}_t \subset B_s \subset \Omega\) and \(k\) satisfying \eqref{4.30}, where \[\gamma := \gamma(a_2, a_3, a_4, g^-, g^+)\quad \text{and} \quad \gamma_1 := \gamma_1(a_2, a_3, a_5, a_6, b_1, b_2, b_3, g^+, h^+_1, F_1, M)\] are positive constants. Inequality \eqref{4.36} implies \(u \in \mathcal{B}_{G(x,t)}(\Omega, M, \gamma, \gamma_1, \delta)\). A similar argument shows \[u \in \mathcal{B}_{G(x,t)}(\overline{\Omega}, M, \gamma, \gamma_1, \delta),\] which completes the proof of Theorem \ref{Theorem 4.2}.
\end{proof}


\begin{proof}[\textbf{Proof of Theorem \ref{Theorem 4.21}}]
The proof proceeds along the same lines as that of Theorem~\ref{Theorem 4.2}, and we therefore omit the details. The only modification arises in the treatment of inequality~\eqref{4.31}. Specifically, by using equation~\eqref{FNV} together with the structural conditions \eqref{6.31}--\eqref{6.41}, \eqref{6.51} and \eqref{6.61}, we obtain the following estimate
\begin{align}\label{4.43}
\begin{split}
a_4 \int_{A_{k,s}} G(x,|Du|) \xi^{g^+} {\rm d}x &\leq a_5 \int_{A_{k,s}} H_1(x,|u|) \xi^{g^+} {\rm d}x +a_6 \int_{A_{k,s}} \xi^{g^+} {\rm d}x+ a_1 g^+ \int_{A_{k,s}} g(x,|Du|) \xi^{g^+-1} |D\xi|(u - k) {\rm d}x\\
&+a_2 g^+\int_{A_{k,s}} P(x,|u|) \xi^{g^+-1} |D\xi|(u - k) {\rm d}x+ a_3g^+ \int_{A_{k,s}} |D\xi| \xi^{g^+-1}(u - k) {\rm d}x\\
&+ b_1 \int_{A_{k,s}}M(x,|Du|) \xi^{g^+} (u - k) {\rm d}x+ b_2 \int_{A_{k,s}}h_1(x,|u|) \xi^{g^+} (u - k) {\rm d}x\\
&+ b_3 \int_{A_{k,s}} \xi^{g^+} (u - k) {\rm d}x+ c_1 \int_{\Gamma_{k,s}}h_2(x,|u|) \xi^{g^+} (u - k) {\rm d}x\\
&+ c_2 \int_{\Gamma_{k,s}} \xi^{g^+} (u - k) {\rm d}x\\
&\leq a_5 \int_{A_{k,s}} H_1(x,|u|) \xi^{g^+} {\rm d}x +a_6 \int_{A_{k,s}} \xi^{g^+} {\rm d}x+ a_1 g^+ \int_{A_{k,s}} g(x,|Du|) \xi^{g^+-1} |D\xi|(u - k) {\rm d}x\\
&+a_2 g^+\int_{A_{k,s}} P(x,|u|) \xi^{g^+-1} |D\xi|(u - k) {\rm d}x+ a_3g^+ \int_{A_{k,s}} |D\xi| \xi^{g^+-1}(u - k) {\rm d}x\\
&+ b_1 \int_{A_{k,s}}M(x,|Du|) \xi^{g^+} (u - k) {\rm d}x+ b_2 \int_{A_{k,s}}h_1(x,|u|) \xi^{g^+} (u - k) {\rm d}x\\
&+ (b_3 +c_2)\int_{A_{k,s}} \xi^{g^+} (u - k) {\rm d}x+ c_1 \int_{A_{k,s}}h_2(x,|u|) \xi^{g^+} (u - k) {\rm d}x,
\end{split}
\end{align}
where,
\(
\Gamma_{k,s} := \left\{ x \in B_s \cap \partial\Omega : u(x) > k \right\}.
\)
\end{proof}
\begin{proof}[\textbf{Proof of Theorems~\ref{Theorem 4.3} and~\ref{Theorem 4.4}}]
Theorems~\ref{Theorem 4.3} and~\ref{Theorem 4.4} follow as direct consequences of Theorems~\ref{Theorem 2.1}--\ref{Theorem 4.21}.
\end{proof}
   %---------------------------------
%\section{New and Recent Results Covered by Our Work}
%In this section, we briefly present some recent and novel results addressed in our work, with particular emphasis on the Hölder continuity of weak solutions.

%As mentioned in the introduction, we investigate the Hölder continuity of weak solutions to the following elliptic problem in divergence form:
%\begin{equation*}
%\operatorname{div} A(x, u, Du) + B(x, u, Du) = 0, \quad x \in \Omega, \tag{$\mathcal{P}$}
%\end{equation*}
%where the functions
%\[
%A : \Omega \times \mathbb{R} \times \mathbb{R}^n \rightarrow \mathbb{R}^n \quad \text{and} \quad
%B : \Omega \times \mathbb{R} \times \mathbb{R}^n \rightarrow \mathbb{R}
%\]
%are Carathéodory functions satisfying the \( G(x,t) \)-growth conditions given in \eqref{6.31}--\eqref{6.51}. Since the structural assumptions on \( A(x,u,Du) \) and \( B(x,u,Du) \) are formulated in terms of a generalized \textnormal{N}-function \( G(x,t) \), we provide several examples of such functions to illustrate the scope of the recent and new results encompassed by our work.
%-------------------------
\ \\
 \textbf{Author Contributions:} The authors contributed equally to this work.\\
\ \\
\textbf{ Data Availability:} Data sharing is not applicable to this article as no new data were created or analyzed in this
 study.\\
 \ \\
 \textbf{Declarations}\\
 \ \\
\textbf{ Ethical Approval:} Not applicable.\\
\ \\
\textbf{Conflict of interest:} The authors declare that they have no Conflict of interest.
%-------------------------

















\begin{thebibliography}{99}
{\small

\bibitem{Ac1} 
E. Acerbi, N. Fusco, 
Partial regularity under anisotropic \((p,q)\) growth conditions, 
{\it J. Differ. Equ.} \textbf{107} (1994), 46–67.

\bibitem{Ac2} 
E. Acerbi, N. Fusco, 
A transmission problem in the calculus of variations, 
{\it Calc. Var. Partial Differ. Equ.} \textbf{2} (1994), 1–16.


\bibitem{Acerbi2007} 
E. Acerbi, G. Mingione,  
Regularity Results for a Class of Functionals with Non-Standard Growth, \textit{Arch. Ration. Mech. Anal.} \textbf{156} (2001), 121–140.


\bibitem{Ala1}A.E. Bahrouni, A. Bahrouni, H. Missaoui, A new class of anisotropic double phase problems: exponents depending on solutions and their gradients, \textit{preprint} (2024), https://arxiv.org/abs/2409.16662.

\bibitem{Ala2}A.E. Bahrouni, A. Bahrouni, P. Winkert, Double phase problems with variable exponents depending on the solution and the gradient in the whole space $\mathbb{R}^N$, \textit{Nonlinear Analysis: Real World Applications} \textbf{85} (2025), 104334.

\bibitem{Ala3}A.E. Bahrouni, A. Bahrouni, The Concentration-Compactness Principle for Musielak-Orlicz spaces and Applications, \textit{preprint} (2025).


\bibitem{Baroni2016}
P. Baroni, M. Colombo, G. Mingione, 
Regularity for general functionals with double phase, 
{\it Calc. Var. Partial Differ. Equ.} \textbf{55} (2016), 62.

%\bibitem{BarlinGmeinederIrvingKristensen2022}M. Bärlin, F. Gmeineder, C. Irving, J. Kristensen,A-harmonic approximation and partial regularity, revisited,{\it arXiv preprint} arXiv:2212.12821 [math.AP] (2022).

%\bibitem{CarozzaKristensenPassarelli2011}M. Carozza, J. Kristensen, A. Passarelli di Napoli,Higher differentiability of minimizers of convex variational integrals,{\it Ann. Inst. H. Poincaré C Anal. Non Linéaire} \textbf{28} (2011), no. 3, 395-411.

%\bibitem{CarozzaKristensenPassarelli2014}M. Carozza, J. Kristensen, A. Passarelli di Napoli,Regularity of minimizers of autonomous convex variational integrals,{\it Annali della Scuola Normale Superiore di Pisa, Classe di Scienze} \textbf{13} (2014), no. 4, 1065-1089.

\bibitem{Chlebicka2022} 
I. Chlebicka,  
Gradient estimates for problems with Orlicz growth,  
{\it Nonlinear Anal.} \textbf{194} (2020), 111364.

\bibitem{Chlebicka2021} 
I. Chlebicka, P. Gwiazda, A. Świerczewska-Gwiazda, A. Wróblewska-Kamińska, 
{\it Partial differential equations in anisotropic Musielak-Orlicz spaces}, 
Springer Monographs in Mathematics, Springer Nature Switzerland (2021).

\bibitem{Cianchi2024}
A. Cianchi, L. Diening,
{\it Sobolev embeddings in Musielak-Orlicz spaces},
Adv. Math. {\bf 447} (2024), Paper No. 109679, 51 pp.


\bibitem{Simon1} S. Ciani, E. Henriques, I. I. Skrypnik,  The weak Harnack inequality for unbounded minimizers of elliptic functionals with generalized Orlicz growth, \textit{Advances in Calculus of Variations} \textbf{18}(3) (2025), 639--652.


\bibitem{Colombo2015} 
M. Colombo, G. Mingione,  
Regularity for double phase variational problems,  
{\it Arch. Ration. Mech. Anal.} \textbf{215} (2015), 443--496.

\bibitem{Colombo2015b}
M. Colombo, G. Mingione,
Bounded minimisers of double phase variational integrals,
{\it Arch. Ration. Mech. Anal.} \textbf{218} (2015), 219-273.

\bibitem{Coscia1999}
A. Coscia, G. Mingione,
Hölder continuity of the gradient of $p(x)$-harmonic mappings,
{\it C. R. Acad. Sci. Paris Sér. I Math.} \textbf{328} (1999), 363-368.

\bibitem{Cupini2017}
G. Cupini, P. Marcellini, E. Mascolo,
Local boundedness of minimizers with limit growth conditions,
{\it J. Optim. Theory Appl.} \textbf{166} (2017), 1-22.


\bibitem{DeFilippisKochKristensen2024}
C. De Filippis, L. Koch, J. Kristensen,
Quantified Legendreness and the regularity of minima,
{\it Arch. Ration. Mech. Anal.} \textbf{248} (2024), no. 4, Paper No. 69, 51 pp.

\bibitem{defi}
 C. De Filippis, G. Mingione, Regularity for Double Phase Problems at Nearly Linear Growth, {\it Arch. Ration. Mech. Anal}. \textbf{247}(5) (2023).
 
\bibitem{deGiorgi1957} 
E. De Giorgi,  
Sulla differenziabilità e l'analiticità delle estremali degli integrali multipli regolari,  
{\it Mem. Accad. Sci. Torino. Cl. Sci. Fis. Mat. Nat.} \textbf{3} (1957), 25--43.

\bibitem{EDI}E. DiBenedetto, Degenerate parabolic equations, \textit{Springer, New York} (1993).

\bibitem{Diening2011} 
L. Diening, P. Harjulehto, P. Hästö, M. R\u{z}ička, 
{\it Lebesgue and Sobolev Spaces with Variable Exponents}, 
Lecture Notes in Math. \textbf{2017}, Springer, Heidelberg (2011).

\bibitem{TD} 
T. Donaldson, 
Nonlinear elliptic boundary value problems in Orlicz–Sobolev space, 
{\it J. Differ. Equ.} \textbf{10} (1971), 507–528.

\bibitem{DonaldsonTrudinger1971} 
T. K. Donaldson, N. S. Trudinger,
Orlicz-Sobolev spaces and imbedding theorems,
{\it J.  Funct. Anal.} \textbf{8} (1971), 52--75.

\bibitem{Eleuteri2016}
M. Eleuteri, P. Marcellini, E. Mascolo,
Lipschitz estimates for systems with ellipticity conditions at infinity,
{\it Ann. Mat. Pura Appl.} \textbf{195} (2016), 1575-1603.

\bibitem{LE} 
L.C. Evans, 
A new proof of local $C^{1,\alpha}$ regularity for solutions of certain degenerate elliptic p.d.e., 
{\it J. Differ. Equ.} \textbf{45} (1982), 356–373.

\bibitem{Fan2012a} 
X. Fan, 
An imbedding theorem for Musielak-Sobolev spaces, 
{\it Nonlinear Anal.} \textbf{75} (2012), 1959–1971.

\bibitem{Fan2012b} 
X. Fan, 
Differential equations of divergence form in Musielak-Sobolev spaces and a sub-supersolution method, 
{\it J. Math. Anal. Appl.} \textbf{386} (2012), 593–604.

\bibitem{Fan1999} 
X. L. Fan, D. Zhao, 
A class of De Giorgi type and Hölder continuity, 
{\it Nonlinear Anal.} \textbf{36} (1999), 295–318.

\bibitem{FanZhao1999} 
X. Fan, D. Zhao,  
On the spaces $L^{p(x)}(\Omega)$ and $W^{m,p(x)}(\Omega)$,   
{\it J. Math. Anal. Appl.} \textbf{263} (1999), 424--446.

\bibitem{mg3} 
I. Fonseca, J. Maly, G. Mingione, 
Scalar minimizers with fractal singular sets, 
{\it Arch. Ration. Mech. Anal.} \textbf{172} (2004), 295-307.

\bibitem{NF} 
N. Fusco, C. Sbordone, 
Some remarks on the regularity of minima of anisotropic integrals, 
{\it Comm. Partial Differential Equations} \textbf{18} (1993), 153–167.

\bibitem{MG1} 
M. Giaquinta, 
{\it Multiple Integrals in the Calculus of Variations and Nonlinear Elliptic Systems}, 
Princeton Univ. Press, Princeton (1983).

\bibitem{MG2} 
M. Giaquinta, 
Growth conditions and regularity, a counterexample, 
{\it Manuscr. Math.} \textbf{59} (1987), 245–248.

\bibitem{MG3} 
M. Giaquinta, E. Giusti, 
On the regularity of the minima of variational integrals, 
{\it Acta Math.} \textbf{148} (1982), 31–46.

\bibitem{GmeinederKristensen2024}
F. Gmeineder, J. Kristensen,
Quasiconvex functionals of (p,q)-growth and the partial regularity of relaxed minimizers,
{\it Arch. Ration. Mech. Anal.} \textbf{248} (2024), no. 5, Paper No. 80, 125 pp.

%\bibitem{GmeinederKristensen2019}F. Gmeineder, J. Kristensen, Partial regularity for BV minimizers,{\it Arch. Ration. Mech. Anal.} \textbf{232} (2019), no. 3, 1429-1473.

%\bibitem{GmeinederKristensen2019preprint}F. Gmeineder, J. Kristensen,Regularity for higher order quasiconvex problems with linear growth from below,{\it arXiv preprint} arXiv:1903.08124 [math.AP] (2019).

\bibitem{Harjulehto2017}
P. Harjulehto, T. Kuusi, T. Lukkari, N. Marola, M. Parviainen, . Harnack's inequality for quasiminimizers with nonstandard growth conditions, \textit{J. Math. Anal. Appl.}, \textbf{344}(1) (2008), 504-520.

\bibitem{Harjulehto2019} 
P. Harjulehto, P. Hästö, 
{\it Orlicz Spaces and Generalized Orlicz Spaces}, 
Springer, Cham (2019).

\bibitem{HsK1} P. Hästö, J. Ok, Maximal regularity for local minimizers of non-autonomous func
tionals. \textit{J.Eur.Math.Soc.}, (2022).

\bibitem{HsK2}P. Harjulehto, P. Hästö, Boundary regularity under generalized growth conditions.
\textit{Z.Anal.Anwend.}, (2019).

\bibitem{HHL21}P. Harjulehto, P. Hästo,  M. Lee, Hölder continuity of $\omega $-minimizers of functionals with generalized Orlicz growth,  \textit{Ann.Sc.Norm.Super.Pisa Cl.Sci.} (2021), 549--582.

\bibitem{MH} 
M.C. Hong, 
Some remarks on the minimizers of variational integrals with nonstandard growth conditions, 
{\it Boll. Unione Mat. Ital. Ser. A} \textbf{6} (7) (1992), 91–102.

\bibitem{KKPZ}K. Ho, Y. H. Kim, P. Winkert, C. Zhang, The boundedness and Hölder continuity of weak solutions to elliptic equations involving variable exponents and critical growth, \textit{ J. Differ. Equ.} \textbf{313} (2022), 503-532

\bibitem{Ho-Sim-2015} 
K. Ho, I. Sim, 
Corrigendum to ``Existence and some properties of solutions for degenerate elliptic equations with exponent variable'' [Nonlinear Anal. 98 (2014) 146–164], 
{\it Nonlinear Anal.} \textbf{128} (2015), 423–426.

%\bibitem{KristensenMelcher2008} J. Kristensen, C. Melcher,Regularity in oscillatory nonlinear elliptic systems,{\it Mathematische Zeitschrift} \textbf{260} (2008), no. 4, 813-847.

\bibitem{KristensenMingione2008}
J. Kristensen, G. Mingione,
Boundary regularity of minima,
{\it Atti della Accademia Nazionale dei Lincei, Classe di Scienze Fisiche, Matematiche e Naturali, Rendiconti Lincei, Matematica e Applicazioni} \textbf{19} (2008), no. 4, 265-277.

\bibitem{KristensenMingione2010}
J. Kristensen, G. Mingione,
Boundary regularity in variational problems,
{\it Arch. Ration. Mech. Anal.} \textbf{198} (2010), no. 2, 369-455.


\bibitem{KristensenTaheri2003}
J. Kristensen, A. Taheri,
Partial regularity of strong local minimizers in the multi-dimensional calculus of variations,
{\it Arch. Ration. Mech. Anal.} \textbf{170} (2003), no. 1, 63-89.

\bibitem{OL} 
O.A. Ladyzhenskaya, N.N. Ural'tseva, 
{\it Linear and Quasilinear Elliptic Equations}, 
2nd Russian ed., Nauka, Moscow (1973).

\bibitem{Li1988} 
G.M. Lieberman, 
Boundary regularity for solutions of degenerate elliptic equations, 
{\it Nonlinear Anal. TMA} \textbf{11} (1988), 1203–1219.

\bibitem{Li1991} 
G.M. Lieberman, 
The natural generalization of the natural conditions of Ladyzhenskaya and Ural'tseva for elliptic equations, 
{\it Comm. Partial Differential Equations} \textbf{16} (1991), 311–361.

\bibitem{Manfredi1986}
J.J. Manfredi,
Regularity for minima of functionals with p-growth,
{\it  J. Differ. Equ.} \textbf{76} (1986), 203-212.

\bibitem{Manfredi1988}
J.J. Manfredi,
Regularity of the gradient for a class of nonlinear elliptic systems,
{\it Arch. Ration. Mech. Anal.} \textbf{102} (1988), 93-102.

\bibitem{marc2} 
P. Marcellini, 
Regularity of minimizers of integrals of the calculus of variations with nonstandard growth conditions, 
{\it Arch. Ration. Mech. Anal.} \textbf{105} (1989), no. 3, 267–284.

\bibitem{Mar1}  P. Marcellini, Un example de solution discontinue d’un problème variationnel dans le cas scalaire, \textit{Preprint 11, 
Istituto Matematico “U. Dini”, Università di Firenze}, (1987).

\bibitem{marc1} 
P. Marcellini, 
Regularity and existence of solutions of elliptic equations with $p,q$-growth conditions, 
{\it J. Differ. Equ.} \textbf{90} (1991), no. 1, 1–30.


\bibitem{mr3} 
P. Marcellini, 
Growth conditions and regularity for weak solutions to nonlinear elliptic PDEs, 
{\it J. Math. Anal. Appl.} \textbf{501} (2021), 124408, 32 pp.

\bibitem{Mingione2006}
G. Mingione,
Regularity of minima: an invitation to the dark side of the calculus of variations,
{\it Appl. Math.} \textbf{51} (2006), 355-426.

\bibitem{mgg} 
G. Mingione, V.D. R\u{a}dulescu, 
Recent developments in problems with nonstandard growth and nonuniform ellipticity, 
{\it J. Math. Anal. Appl.} \textbf{501} (2021), 125197, 41 pp.


\bibitem{TS2} 
Y. Mizuta, E. Nakai, T. Ohno, T. Shimomura, 
Riesz potentials and Sobolev embeddings on Morrey spaces of variable exponent, 
{\it Complex Var. Elliptic Equ.} \textbf{56}(7–9) (2011), 671–695.

\bibitem{TS3} 
Y. Mizuta, E. Nakai, T. Ohno, T. Shimomura, 
Maximal functions, Riesz potentials and Sobolev embeddings on Musielak-Orlicz-Morrey spaces of variable exponent in $\mathbb{R}^n$, 
{\it Rev. Mat. Complut.} \textbf{25}(2) (2012), 413–434.

\bibitem{TS4} 
Y. Mizuta, E. Nakai, T. Ohno, T. Shimomura, 
Campanato-Morrey spaces for the double phase functionals with variable exponents, 
{\it Nonlinear Anal.} \textbf{201} (2020), 111827.

\bibitem{TS5} 
Y. Mizuta, T. Ohno, T. Shimomura, 
Sobolev's inequalities and vanishing integrability for Riesz potentials of functions in the generalized Lebesgue space $L^{p(\cdot)}(\log L)^{q(\cdot)}$, 
{\it J. Math. Anal. Appl.} \textbf{345} (2008), 70–85.

\bibitem{TS6} 
Y. Mizuta, T. Ohno, T. Shimomura, 
Sobolev inequalities for Musielak-Orlicz spaces, 
{\it Manuscripta Math.} \textbf{155} (2018), 209–227.

\bibitem{TS7} 
Y. Mizuta, T. Ohno, T. Shimomura, 
Sobolev's theorem for double phase functionals, 
{\it Math. Inequal. Appl.} \textbf{23} (2020), 17–33.

\bibitem{TS8} 
Y. Mizuta, T. Shimomura, 
Sobolev embeddings for Riesz potentials of functions in Morrey spaces of variable exponent, 
{\it J. Math. Soc. Japan} \textbf{60} (2008), 583–602.

\bibitem{TS9} 
Y. Mizuta, T. Shimomura, 
Sobolev's inequality for Riesz potentials of functions in Morrey spaces of integral form, 
{\it Math. Nachr.} \textbf{283}(9) (2010), 1336–1352.

\bibitem{TS10} 
Y. Mizuta, T. Shimomura, 
Hardy-Sobolev inequalities in the unit ball for double phase functionals, 
{\it J. Math. Anal. Appl.} \textbf{499} (2021), 124133.

\bibitem{CM} 
C.B. Morrey, 
{\it Multiple Integrals in the Calculus of Variations}, 
Springer, Berlin (1968).

\bibitem{Moser1960} 
J. Moser,  
A new proof of De Giorgi's theorem concerning the regularity problem for elliptic differential equations,  
{\it Comm. Pure Appl. Math.} \textbf{13} (1960), 457--468.

\bibitem{Moser1961} 
J. Moser,  
On Harnack's theorem for elliptic differential equations,  
{\it Comm. Pure Appl. Math.} \textbf{14} (1961), 577--591.

\bibitem{Musielak1983} 
J. Musielak, 
{\it Orlicz Spaces and Modular Spaces}, 
Lecture Notes in Math. \textbf{1034}, Springer, Berlin (1983).

\bibitem{Nash1958} 
J. Nash,  
Continuity of solutions of parabolic and elliptic equations,  
{\it Amer. J. Math.} \textbf{80} (1958), 931--954.

\bibitem{TS1} 
T. Ohno, T. Shimomura, 
Sobolev-type inequalities on Musielak–Orlicz–Morrey spaces of an integral form, 
{\it Bull. Malays. Math. Sci. Soc.} \textbf{46} (1) (2023), 31.

\bibitem{TS11} 
T. Ohno, T. Shimomura, 
Sobolev's inequality for Musielak-Orlicz-Morrey spaces over metric measure spaces, 
{\it J. Aust. Math. Soc.} \textbf{110} (2021), 371–385.

\bibitem{TS12} 
T. Ohno, T. Shimomura, 
Sobolev-type inequalities on variable exponent Morrey spaces of an integral form, 
{\it Ricerche Mat.} \textbf{71} (2022), 189–204.

\bibitem{Papageorgiou-Winkert-2018} 
N.S. Papageorgiou, P. Winkert, 
{\it Applied Nonlinear Functional Analysis: An Introduction}, 
De Gruyter, Berlin (2018).

\bibitem{Pucci} 
 P. Pucci, R. Servadei, Regularity of weak solutions of homogeneous or inhomogeneous quasilinear elliptic equations, \textit{Indiana University mathematics journal} \textbf{57} (2008), 3329-3363.

\bibitem{Uhlenbeck1977} 
K. Uhlenbeck,  
Regularity for a class of non-linear elliptic systems,  
{\it Acta Math.} \textbf{138} (1977), 219--240.

\bibitem{Uraltseva1968} 
N. N. Uraltseva,  
Degenerate quasilinear elliptic systems,  
{\it Zap. Nauchn. Sem. Leningrad. Otdel. Mat. Inst. Steklov. (LOMI)} \textbf{7} (1968), 184--222.

\bibitem{WaLiuZhao} 
B. Wang, D. Liu, P. Zhao,
H\"older continuity for nonlinear elliptic problem in Musielak-Orlicz-Sobolev space,
{\it  J. Differ. Equ.} \textbf{266} (2019), 4835-4863.

\bibitem{Zhikov1987}
V.V. Zhikov, Averaging of functionals of the calculus of variations and elasticity theory, \textit{Izv. Akad. Nauk
SSSR Ser. Mat.} \textbf{50} (1986), 675--710.

\bibitem{Zhikov1995}
V. V. Zhikov,
On Lavrentiev's phenomenon,
\textit{Russian J. Math. Phys} \textbf{3} (1995), 249--269.

}
\end{thebibliography}
%%%%%%%%%%%%%%%%%%%%%%%%%%%%%%%%%%%%%%%%%%%%%%%%
%%%%%%%%%%%%%%%%%%%%%%%%%%%%%%%%%%%%%%%%%%%%%%%%%%%%%%%%%%%%%
\end{document}

The qualitative theory of partial differential equations is concerned with the behavior of weak solutions 
$u \in W^{1,1}_{\mathrm{loc}}(\Omega)$ of problems of the form
\begin{equation}\label{main_eq}
-\operatorname{div} A(x,u,D u) = B(x,u,D u) \quad \text{in } \Omega.
\end{equation}
Existence results based on the calculus of variations and topological methods typically yield solutions only in Sobolev spaces. 
However, for both theoretical and applied purposes, it is essential to determine whether these solutions possess higher regularity. 
In particular, Hölder continuity $u \in C^{0,\alpha}_{\mathrm{loc}}(\Omega)$ plays a central role, lying between boundedness ($L^\infty$) and differentiability ($C^{1,\alpha}$), and being crucial for uniqueness, stability, and free boundary problems.\\

The modern theory of regularity began with the seminal works of De~Giorgi~\cite{deGiorgi1957} and Nash~\cite{Nash1958}, later completed by Moser~\cite{Moser1960,Moser1961}. 
They established that weak solutions of linear elliptic equations in divergence form,
\[
-\operatorname{div}(A(x)D u)=0,
\]
with merely bounded measurable coefficients $A(x)$ satisfying the uniform ellipticity condition 
$\lambda |\xi|^2 \le \langle A(x)\xi,\xi\rangle \le \Lambda |\xi|^2$, 
are Hölder continuous. 
This result was striking since it required no smoothness assumptions on $A(x)$. 
It also introduced two fundamental tools: De Giorgi’s measure-theoretic argument and Moser’s iteration technique.

Attention then shifted to quasilinear equations, notably the $p$-Laplacian
\[
\Delta_p u := \operatorname{div}(|D u|^{p-2}D u).
\]
Uraltseva~\cite{Uraltseva1968} and Uhlenbeck~\cite{Uhlenbeck1977} proved that solutions of $-\Delta_p u=0$ are $C^{1,\alpha}$, revealing deeper regularity phenomena. 
However, this regularity is delicate: for equations of the type $-\Delta_p u=f$, if $f \in L^q$ with $q<n/p$, solutions may fail to be even bounded~\cite{Brasco2018}. 
This illustrates that the regularity program must typically proceed in two stages: first, establish boundedness via De Giorgi’s lemma or Moser iteration; then derive Hölder continuity.

These ideas have since been extended to a variety of operators with nonstandard growth. 
In the Orlicz framework, where the growth is determined by an $N$-function $G(t)$, Lieberman~\cite{Li1991} established Hölder continuity under the $\Delta_2$-condition. 
Subsequent developments addressed models with growth depending on the spatial variable $x$. 
For double-phase functionals,
\[
\operatorname{div}\!\left(|D u|^{p-2}D u + a(x)|D u|^{q-2}D u\right),
\]
Colombo and Mingione~\cite{Colombo2015} showed that Hölder regularity requires the coefficient $a(x)$ to be Hölder continuous. 
For variable exponent problems, introduced by Fan and Zhao~\cite{FanZhao1999}, regularity is preserved provided that the exponent $p(x)$ satisfies the log-Hölder condition~\cite{Acerbi2007}. 
For recent surveys and comprehensive bibliographies on the regularity of nonlinear elliptic equations, see~\cite{mr3, Mingione2006}. 
For detailed discussions of particular models, we refer to: \cite{deGiorgi1957, Nash1958, Moser1960, Moser1961} (linear theory); 
\cite{Uraltseva1968, Uhlenbeck1977, Manfredi1986, Manfredi1988, marc1, marc2, Ac1, Ac2, mg3, MG3, NF, LE, MH, Pucci} ($p$-Laplacian and $(p,q)$-growth); 
\cite{Li1988, Li1991, Harjulehto2017, Adamowicz2015, Cupini2017, DallAglio2002, Eleuteri2016, Chlebicka2021, Chlebicka2022, mgg, Brasco2018, Coscia1999} (Orlicz and nonstandard growth);
\cite{Colombo2015, Colombo2015b, Baroni2016} (double phase problems); 
and \cite{Fan1999, FanZhao1999, Acerbi2007, Diening2021, KKPZ} (variable exponent problems).\\

In this paper, we establish Hölder continuity of weak solutions to~\eqref{main_eq} under Musielak–Orlicz growth assumptions on $A$ and $B$. 
Our framework unifies and extends the classical theories of the $p$-Laplacian, variable exponent, and Orlicz problems, covering both subcritical and critical regimes with Dirichlet or Neumann boundary conditions.

To the best of our knowledge, the only existing result in this direction is due to Wang, Liu, and Zhao~\cite{WaLiuZhao}, who proved Hölder continuity under a structural condition on the function $G(x,t)$. 
Specifically, they assumed the existence of a strictly increasing differentiable function $\xi : [0,+\infty) \to [0,+\infty)$ satisfying
\[
n\xi(s) > s \xi'(s) > \xi(s) \quad \text{for all } s \geq 0,
\]
together with the conditions
\begin{itemize}
    \item[(A11)] $G(x, s t) \geq \xi(s) G(x,t)$ for all $s \geq 0$, $t \in \mathbb{R}$, and $x \in \Omega$;
    \item[(A12)] $\xi, \xi^{-1}, \xi_{\ast}, \widehat{\widetilde{\xi}} \in \Delta_{\mathbb{R}^{+}}$.
\end{itemize}
While interesting, this framework is restrictive: in particular, (A11) excludes many fundamental models, including variable exponent spaces, $(p,q)$-growth equations, and double-phase functionals—the very cases that a general Musielak–Orlicz theory should encompass. 
Thus, although \cite{WaLiuZhao} provides valuable insight, it lacks the generality required to unify classical and modern examples.\\

The novelty of our work lies in establishing Hölder regularity for divergence-form problems under broad Musielak–Orlicz growth conditions, within a unified framework that recovers and extends the $p$-Laplacian, variable exponent, and Orlicz cases. 
Our approach, inspired by Lieberman~\cite{Li1991} and Fan--Zhao~\cite{FanZhao1999}, builds on the classical techniques of De~Giorgi~\cite{deGiorgi1957} and Ladyzhenskaya--Ural'tseva~\cite{OL}, which introduced the class $\mathcal{B}(\Omega, M, \gamma, \gamma_1, \delta, 1/q)$ and proved Hölder continuity for all its members under standard $m$-growth conditions (see, e.g.,~\cite{MG1, MG3, OL, CM}). 
For critical cases, we further rely on the method of Ho, Kim, Winkert, and Zhang~\cite{KKPZ}.\\

Three main challenges arise in the Musielak–Orlicz setting.  
First, the energy is nonhomogeneous, depending on both $x$ and $t$, which prevents the direct use of classical De Giorgi iteration.  
Second, regularity is highly sensitive to the structural parameters governing the growth of $G(x,t)$.  
Third, even when De Giorgi-type methods apply, regularity is not automatic: counterexamples by Giaquinta~\cite{MG2} and Marcellini~\cite{marc1, marc2} show that weak solutions may fail to be regular under seemingly natural assumptions.  
This makes the task of proving regularity under nonstandard growth conditions especially delicate.
